% Lesson 53 — General revision: Revise previous aspects of grammar — Anglais Tle A — Cours signé OnBuch+
% Programme officiel MINESEC. Compiler avec : tectonic EN53-general-revision-53.tex
\newcommand{\DOCMATIERE}{Anglais}
\newcommand{\DOCNIVEAU}{Tle A}
\newcommand{\DOCMODULE}{Chapter 5 : Media and Communication}
\newcommand{\DOCLECON}{EN53}
\newcommand{\DOCTITRE}{Lesson 53 — General revision: Revise previous aspects of grammar}
% OnBuch+ — préambule « cours signés » ENRICHI (compilé avec tectonic / XeLaTeX)
% Système visuel « Pop » d'OnBuch+ : fond crème (jamais blanc), bordures encre
% épaisses, ombres « dures » décalées, titres Archivo Black en capitales.
% Version MATHÉMATIQUES : graphes (pgfplots/tikz), boîtes pédagogiques
% (définition, propriété, méthode, attention, le savais-tu, exercice résolu,
% exercice, TP, bilan), page de garde et signature OnBuch+ ancrée en bas.
%
% Macros attendues dans main.tex AVANT \input{preamble} :
%   \DOCMATIERE  \DOCNIVEAU  \DOCMODULE  \DOCLECON (numéro)  \DOCTITRE
\documentclass[11pt]{article}

\usepackage{fontspec}
\usepackage[english,french]{babel} % français = langue principale ; l'anglais via \eng{..} / textebox
\usepackage[a4paper,margin=2cm,top=3.4cm,bottom=2.8cm,headheight=1.4cm,footskip=1.3cm]{geometry}
\usepackage{amsmath,amssymb,amsfonts}
\usepackage{mathtools} % \xmapsto, \coloneqq... souvent utilisés par les modèles
\usepackage{cancel} % simplifications barrées (bilans de forces)
% \abs{z} (module, valeur absolue) et \norm{u} (norme), utilisés par les modèles
\providecommand{\abs}{}\let\abs\relax\DeclarePairedDelimiter{\abs}{\lvert}{\rvert}
\providecommand{\norm}{}\let\norm\relax\DeclarePairedDelimiter{\norm}{\lVert}{\rVert}
\usepackage[table]{xcolor} % [table] : fournit aussi \rowcolors (alternance de lignes)
\usepackage{pagecolor}
\usepackage{tikz}
\usetikzlibrary{arrows.meta,positioning,calc,shapes.geometric,shapes.misc,shapes.arrows,decorations.pathmorphing,decorations.pathreplacing,decorations.markings,patterns,shadows,mindmap,trees,fit,backgrounds,matrix,intersections,angles,quotes}
\usepackage{pgfplots}
\usepackage[version=4]{mhchem} % équations nucléaires \ce{^{235}_{92}U}
\usepackage[european,siunitx]{circuitikz} % schémas électriques
\pgfplotsset{compat=1.18}
\usepgfplotslibrary{fillbetween,groupplots} % aires entre courbes (calcul intégral)
\usepackage{enumitem}
\usepackage{fancyhdr}
% [most] charge toutes les bibliothèques tcolorbox (listings, minted, raster,
% documentation...) et gonfle inutilement la mémoire TeX sur un document long
% (~300 boîtes) : on ne charge que ce qui est réellement utilisé.
\usepackage[skins,breakable]{tcolorbox}
\usepackage{graphicx}
\usepackage{colortbl}
\usepackage{booktabs}
\usepackage{tabularx}
\usepackage{diagbox} % case d'en-tête diagonale des tableaux à double entrée
% Colonnes extensibles centrée / gauche / droite (C, L, R), souvent utilisées par les modèles
\newcolumntype{C}{>{\centering\arraybackslash}X}
\newcolumntype{L}{>{\raggedright\arraybackslash}X}
\newcolumntype{R}{>{\raggedleft\arraybackslash}X}
\usepackage{array}
\usepackage{multicol}
\usepackage{multirow}
\usepackage{siunitx}
% parse-numbers=false : accepte aussi \SI{2,5 \times 10^{-3}}{...}
\sisetup{locale=FR,output-decimal-marker={,},group-minimum-digits=4,parse-numbers=false}
\DeclareSIUnit\ohm{\ensuremath{\symup{\Omega}}} % U+2126 absent de la police
\DeclareSIUnit\division{div} % réglages d'oscilloscope (ms/div, V/div)
\DeclareSIUnit\tr{tr} % tours (vitesse de rotation, ex. \SI{1500}{\tr\per\minute})
\DeclareSIUnit\molar{M} % concentration molaire (ex. \SI{3}{\molar}, \SI{1}{\micro\molar})
\DeclareSIUnit\an{an} % année (le modèle confond parfois avec \year, un compteur TeX) — ex. \SI{5}{\kilo\meter\per\an}
\DeclareSIUnit\FCFA{FCFA} % franc CFA (ex. \SI{500}{\FCFA\per\kilo\gram})
\DeclareSIUnit\degreeBrix{\textdegree Brix} % teneur en sucre d'un sirop/jus (réfractométrie)
\DeclareSIUnit\month{mois} % le modèle utilise parfois \month, un compteur TeX, comme unité de durée
\usepackage{qrcode}
% \degree : commande fréquemment utilisée par les modèles pour un angle ou une
% température en dehors de \SI{}{} (non fournie par les paquets chargés).
\providecommand{\degree}{^{\circ}}
% \qed (fin de démonstration), utilisé par les modèles sans amsthm
\providecommand{\qed}{\hfill$\square$}
% \barre : double barre des valeurs interdites dans un tableau de variations (tkz-tab)
\providecommand{\barre}{\ensuremath{\|}}
% environnement proof (amsthm non chargé) utilisé par les modèles
\ifdefined\proof\else\newenvironment{proof}[1][Démonstration]{\par\noindent\textit{#1.}\ }{\qed\par}\fi
\usepackage{lastpage}
\usepackage{hyperref}
\hypersetup{hidelinks}

% ============================================================
% Palette « Pop » OnBuch+ — mêmes hex que la classe P (plus_ui.dart)
% ============================================================
\definecolor{poporange}{HTML}{F59321}
\definecolor{popdark}{HTML}{E07A0C}
\definecolor{popcream}{HTML}{FFF6E8}
\definecolor{popink}{HTML}{1C1714}
\definecolor{poppurple}{HTML}{7A5AE0}
\definecolor{popgreen}{HTML}{2E9E6A}
\definecolor{poppink}{HTML}{E0457B}
\definecolor{popblue}{HTML}{2F7DE1}
\definecolor{popgold}{HTML}{C98A12}
\definecolor{popmuted}{HTML}{8A7F76}
\definecolor{popline}{HTML}{EADFD2}
\definecolor{popgray}{HTML}{8A7F76}
\colorlet{popgrey}{popgray}
% Alias tolérants (le modèle nomme parfois les couleurs différemment)
\colorlet{popwhite}{white}
\colorlet{popblack}{popink}
\colorlet{popred}{poppink}
\colorlet{popyellow}{popgold}
% Teintes claires (fonds de tableaux/encadrés) — le modèle nomme parfois
% les couleurs avec un suffixe L (ex. popblueL) sans les avoir définies.
\colorlet{poporangeL}{poporange!20}
\colorlet{popdarkL}{popdark!20}
\colorlet{poppurpleL}{poppurple!20}
\colorlet{popgreenL}{popgreen!20}
\colorlet{poppinkL}{poppink!20}
\colorlet{popblueL}{popblue!20}
\colorlet{popgoldL}{popgold!20}
\colorlet{popredL}{popred!20}
\colorlet{popgrayL}{popgray!20}
% Teintes claires pour fonds de tableaux / zones
\colorlet{popblueL}{popblue!10!white}
\colorlet{poporangeL}{poporange!14!white}
\colorlet{popgreenL}{popgreen!12!white}
\colorlet{poppurpleL}{poppurple!12!white}
\colorlet{poppinkL}{poppink!12!white}
\colorlet{popgoldL}{popgold!14!white}

\pagecolor{popcream}

% ============================================================
% Typographie
% ============================================================
\defaultfontfeatures{Ligatures=TeX}
\setmainfont{PlusJakartaSans-Medium.ttf}[Path=../../../fonts/,BoldFont=PlusJakartaSans-Bold.ttf]
% Symboles, API et emojis absents de la police principale : repli sur DejaVu Sans (caractères actifs)
\newfontface\symfont{DejaVuSans.ttf}[Path=../../../fonts/]
\makeatletter
\newcommand\symactive[1]{\catcode#1=13 \begingroup\lccode`\~=#1 \lowercase{\endgroup\def~}{{\symfont\char#1}}}
\newcommand\symblank[1]{\catcode#1=13 \begingroup\lccode`\~=#1 \lowercase{\endgroup\def~}{}}
\newcommand\symhyph[1]{\catcode#1=13 \begingroup\lccode`\~=#1 \lowercase{\endgroup\def~}{\mbox{-}}}
\newcount\symc
\newcommand\symrange[2]{\symc=#1 \@whilenum\symc<#2 \do{\expandafter\symactive\expandafter{\the\symc}\advance\symc by 1 }}
\newcommand\symblankrange[2]{\symc=#1 \@whilenum\symc<#2 \do{\expandafter\symblank\expandafter{\the\symc}\advance\symc by 1 }}
\makeatother
\symrange{"0250}{"02FF}\symactive{"2713}\symactive{"2714}\symactive{"2717}\symactive{"2718}\symactive{"2610}\symactive{"2611}\symactive{"2612}
\symhyph{"2011}\symhyph{"2E1A}\symblankrange{"1F300}{"1FAFF}\symblank{"FE0F}
% Maths en sans-serif (Fira Math) avec chiffres et lettres droites de Plus
% Jakarta, pour que les formules se fondent dans le texte.
\usepackage[math-style=ISO,bold-style=ISO]{unicode-math}
\setmathfont{FiraMath-Regular.otf}[Path=../../../fonts/]
\setmathfont{PlusJakartaSans-Medium.ttf}[Path=../../../fonts/,range={up/num,up/{Latin,latin},\mathup}]
\usepackage{icomma} % virgule décimale sans espace en mode math
% Caractères mathématiques Unicode absents des polices de texte (Plus Jakarta,
% Archivo Black) : rendus par leur équivalent mathématique, en texte comme en
% formule — sinon ils s'affichent en carré vide (ex. « Arithmétique dans ℤ »).
\usepackage{newunicodechar}
\newunicodechar{ℝ}{\ensuremath{\mathbb{R}}}
\newunicodechar{ℤ}{\ensuremath{\mathbb{Z}}}
\newunicodechar{ℂ}{\ensuremath{\mathbb{C}}}
\newunicodechar{ℕ}{\ensuremath{\mathbb{N}}}
\newunicodechar{ℚ}{\ensuremath{\mathbb{Q}}}
\newunicodechar{𝔻}{\ensuremath{\mathbb{D}}}
\newunicodechar{α}{\ensuremath{\alpha}}
\newunicodechar{β}{\ensuremath{\beta}}
\newunicodechar{γ}{\ensuremath{\gamma}}
\newunicodechar{δ}{\ensuremath{\delta}}
\newunicodechar{ε}{\ensuremath{\varepsilon}}
\newunicodechar{θ}{\ensuremath{\theta}}
\newunicodechar{λ}{\ensuremath{\lambda}}
\newunicodechar{μ}{\ensuremath{\mu}}
\newunicodechar{σ}{\ensuremath{\sigma}}
\newunicodechar{τ}{\ensuremath{\tau}}
\newunicodechar{φ}{\ensuremath{\varphi}}
\newunicodechar{ψ}{\ensuremath{\psi}}
\newunicodechar{ω}{\ensuremath{\omega}}
\newunicodechar{Σ}{\ensuremath{\Sigma}}
% majuscules produites par \MakeUppercase dans les titres (α-aminés -> Α)
\newunicodechar{Α}{\ensuremath{\alpha}}
\newunicodechar{Β}{\ensuremath{\beta}}
\newunicodechar{Γ}{\ensuremath{\Gamma}}
\newunicodechar{Δ}{\ensuremath{\Delta}}
\newunicodechar{Ε}{\ensuremath{\varepsilon}}
\newunicodechar{Θ}{\ensuremath{\Theta}}
\newunicodechar{Λ}{\ensuremath{\Lambda}}
\newunicodechar{Μ}{\ensuremath{\mu}}
\newunicodechar{Τ}{\ensuremath{\tau}}
\newunicodechar{Φ}{\ensuremath{\Phi}}
\newunicodechar{Ψ}{\ensuremath{\Psi}}
\newunicodechar{Ω}{\ensuremath{\Omega}}
\newunicodechar{↦}{\ensuremath{\mapsto}}
\newunicodechar{∈}{\ensuremath{\in}}
\newunicodechar{∉}{\ensuremath{\notin}}
\newunicodechar{∧}{\ensuremath{\wedge}}
\newunicodechar{⇔}{\ensuremath{\Leftrightarrow}}
\newunicodechar{⟺}{\ensuremath{\Longleftrightarrow}}
\newunicodechar{⇒}{\ensuremath{\Rightarrow}}
\newunicodechar{⟹}{\ensuremath{\Longrightarrow}}
\newunicodechar{∘}{\ensuremath{\circ}}
\newunicodechar{∪}{\ensuremath{\cup}}
\newunicodechar{∩}{\ensuremath{\cap}}
\newunicodechar{≡}{\ensuremath{\equiv}}
\newunicodechar{⊂}{\ensuremath{\subset}}
\newunicodechar{⊥}{\ensuremath{\perp}}
\newunicodechar{∃}{\ensuremath{\exists}}
\newunicodechar{∀}{\ensuremath{\forall}}
\newunicodechar{∛}{\ensuremath{\sqrt[3]{\,}}}
\newunicodechar{∜}{\ensuremath{\sqrt[4]{\,}}}
\newunicodechar{⃗}{\ensuremath{{}^{\to}}}
\newunicodechar{ⁿ}{\ifmmode^{n}\else\textsuperscript{\ensuremath{n}}\fi}
\newunicodechar{ˣ}{\ifmmode^{x}\else\textsuperscript{\ensuremath{x}}\fi}
\newunicodechar{ᵇ}{\ifmmode^{b}\else\textsuperscript{\ensuremath{b}}\fi}
\newunicodechar{ʳ}{\ifmmode^{r}\else\textsuperscript{\ensuremath{r}}\fi}
\newunicodechar{ᵘ}{\ifmmode^{u}\else\textsuperscript{\ensuremath{u}}\fi}
\newunicodechar{ʸ}{\ifmmode^{y}\else\textsuperscript{\ensuremath{y}}\fi}
\newunicodechar{ᵃ}{\ifmmode^{a}\else\textsuperscript{\ensuremath{a}}\fi}
\newunicodechar{ᵉ}{\ifmmode^{e}\else\textsuperscript{\ensuremath{e}}\fi}
\newunicodechar{ᵏ}{\ifmmode^{k}\else\textsuperscript{\ensuremath{k}}\fi}
\newunicodechar{ᵖ}{\ifmmode^{p}\else\textsuperscript{\ensuremath{p}}\fi}
\newunicodechar{ᴺ}{\ifmmode^{N}\else\textsuperscript{\ensuremath{N}}\fi}
\newunicodechar{ᶜ}{\ifmmode^{c}\else\textsuperscript{\ensuremath{c}}\fi}
\newunicodechar{ʰ}{\ifmmode^{h}\else\textsuperscript{\ensuremath{h}}\fi}
\newunicodechar{ᵠ}{\ifmmode^{q}\else\textsuperscript{\ensuremath{q}}\fi}
\newunicodechar{ₙ}{\ifmmode_{n}\else\textsubscript{\ensuremath{n}}\fi}
\newunicodechar{ₐ}{\ifmmode_{a}\else\textsubscript{\ensuremath{a}}\fi}
\newunicodechar{ᵢ}{\ifmmode_{i}\else\textsubscript{\ensuremath{i}}\fi}
\newunicodechar{ᵣ}{\ifmmode_{r}\else\textsubscript{\ensuremath{r}}\fi}
\newunicodechar{ₖ}{\ifmmode_{k}\else\textsubscript{\ensuremath{k}}\fi}
\newunicodechar{ᵦ}{\ifmmode_{\beta}\else\textsubscript{\ensuremath{\beta}}\fi}
\newunicodechar{ₓ}{\ifmmode_{x}\else\textsubscript{\ensuremath{x}}\fi}
\newfontfamily\popdisplay{ArchivoBlack-Regular.ttf}[Path=../../../fonts/]
\newfontfamily\poplogo{SpaceGrotesk-Bold.ttf}[Path=../../../fonts/]
\newfontfamily\popsemi{PlusJakartaSans-SemiBold.ttf}[Path=../../../fonts/]
\newfontfamily\popxbold{PlusJakartaSans-ExtraBold.ttf}[Path=../../../fonts/]
\newfontfamily\popxboldls{PlusJakartaSans-ExtraBold.ttf}[Path=../../../fonts/,LetterSpace=3.0]

\setlist[enumerate]{leftmargin=1.7em,itemsep=5pt,topsep=4pt}
\setlist[itemize]{leftmargin=1.7em,itemsep=4pt,topsep=4pt,label={\color{poporange}\textbullet}}
\setlength{\parindent}{0pt}
\setlength{\parskip}{5pt}
\renewcommand{\arraystretch}{1.25}

% pgfplots : style Pop par défaut
\pgfplotsset{
  every axis/.append style={axis line style={popink,line width=1pt},tick style={popink},
    grid style={popline,line width=0.5pt},label style={font=\small},tick label style={font=\footnotesize}},
  popcurve/.style={poporange,line width=1.8pt,no markers,smooth},
}
% Même style disponible dans un \draw TikZ simple (hors pgfplots)
\tikzset{popcurve/.style={poporange,line width=1.8pt}}
\tikzset{popgrid/.style={popline,very thin}}
\colorlet{popcurve}{poporange} % le nom du style est parfois utilisé comme couleur

% ============================================================
% Pill / squiggle
% ============================================================
\newcommand{\poppill}[3]{%
  \tikz[baseline=-0.55ex]\node[draw=popink,line width=1.2pt,rounded corners=2.6pt,%
    fill=#2,inner xsep=6.5pt,inner ysep=3pt]{%
    \popxboldls\fontsize{7}{7}\selectfont\color{#3}\MakeUppercase{#1}};%
}
\newcommand{\popsquiggle}[1]{%
  \tikz[baseline=-0.4ex]{
    \draw[#1,line width=2.6pt,line cap=round]
      plot [smooth] coordinates {(0,0.09) (0.34,0.01) (0.68,0.21) (1.02,0.01) (1.36,0.10)};
  }%
}
% Mot-clé surligné (marqueur orange sous le texte)
\newcommand{\cle}[1]{{\popxbold\color{popdark}#1}}

% ============================================================
% En-tête / pied
% ============================================================
\pagestyle{fancy}
\fancyhf{}
\renewcommand{\headrulewidth}{0pt}
\renewcommand{\footrulewidth}{0pt}
\lhead{\raisebox{-0.15ex}{\poplogo\fontsize{13}{13}\selectfont\color{popink}OnBuch\color{poporange}+}}
\rhead{\poppill{\DOCMATIERE\ \textbullet\ \DOCNIVEAU\ \textbullet\ Leçon \DOCLECON}{white}{popink}}
\lfoot{\popxbold\fontsize{7.6}{7.6}\selectfont\color{popmuted}\MakeUppercase{Cours signé OnBuch+}\ \textbullet\ {\popsemi\DOCTITRE}}
\rfoot{\tikz[baseline=-0.6ex]\node[rounded corners=8.5pt,fill=popink,minimum height=17pt,minimum width=17pt,inner xsep=5pt,inner sep=0]{\popxbold\fontsize{8}{8}\selectfont\color{popcream}\thepage\,/\,\pageref*{LastPage}};}

% ============================================================
% Titres
% ============================================================
\newcommand{\chaptitre}[2]{%
  \begin{tcolorbox}[enhanced,colback=poporange,colframe=popink,boxrule=2.6pt,arc=11pt,
      drop shadow=popink,left=15pt,right=15pt,top=12pt,bottom=11pt,before skip=3pt,after skip=18pt]
    \poppill{#2}{popink}{popcream}\par\vspace{9pt}
    {\raggedright\hyphenpenalty=10000\exhyphenpenalty=10000\popdisplay\fontsize{22}{24}\selectfont\color{popcream}\MakeUppercase{#1}\par}
  \end{tcolorbox}
}
% Section numérotée : pastille orange avec le numéro + titre Archivo
\newcounter{popsec}
\newcommand{\coursec}[1]{%
  \stepcounter{popsec}\setcounter{popsub}{0}%
  \par\vspace{16pt}\noindent
  \begin{minipage}{\linewidth}
  \tikz[baseline=-0.7ex]\node[draw=popink,line width=1.6pt,fill=poporange,rounded corners=4pt,minimum size=22pt,inner sep=0]{\popdisplay\fontsize{12}{12}\selectfont\color{popcream}\thepopsec};%
  \hspace{8pt}{\popdisplay\fontsize{15}{17}\selectfont\color{popink}\MakeUppercase{#1}}\par
  \vspace{4pt}\noindent{\color{poporange}\rule{36pt}{3.6pt}}
  \end{minipage}\par\nopagebreak\vspace{8pt}
}
\newcounter{popsub}
\newcommand{\courssub}[1]{%
  \stepcounter{popsub}%
  \par\vspace{9pt}\noindent{\popxbold\fontsize{11.5}{13}\selectfont\color{popink}{\color{poporange}\thepopsec.\thepopsub}\hspace{6pt}#1}\par\nopagebreak\vspace{4pt}
}

% ============================================================
% Boîtes pédagogiques — même recette : fond blanc, bordure épaisse colorée,
% ombre dure encre, badge plein. Titre optionnel [#1].
% ============================================================
\tcbset{popbase/.style={breakable,enhanced,colback=white,boxrule=2.3pt,arc=9pt,
  shadow={3pt}{-3pt}{0pt}{popink},left=14pt,right=14pt,top=11pt,bottom=11pt,before skip=11pt,after skip=15pt,
  fontupper=\popsemi\fontsize{10.6}{14.5}\selectfont\color{popink},
  fontlower=\popsemi\fontsize{10.6}{14.5}\selectfont\color{popink}}}
% Coche dessinée (le glyphe ✓ n'existe pas dans les polices du document)
\newcommand{\popcheck}[1]{\tikz[baseline=-0.5ex]\draw[#1,line width=1.6pt,line cap=round,line join=round](-0.13,0.0)--(-0.04,-0.09)--(0.13,0.1);}
\providecommand{\coche}{\popcheck{popgreen}}
\providecommand{\resultat}[1]{\par\smallskip\noindent\fcolorbox{popgreen}{popgreenL}{\parbox{\dimexpr\linewidth-2\fboxsep-2\fboxrule\relax}{\textbf{Résultat.} #1}}\par\smallskip}
\newcommand{\popboxhead}[3]{\poppill{#1}{#2}{white}\ifx\relax#3\relax\else\hspace{6pt}{\popxbold\fontsize{10.6}{12}\selectfont\color{popink}#3}\fi\par\vspace{7pt}}

\newtcolorbox{aretenir}[1][]{popbase,colframe=popgreen,before upper={\popboxhead{À retenir}{popgreen}{#1}}}
% Texte en anglais (document de lecture, dialogue, extrait) : cadre
% d'étude. Un titre optionnel (source, auteur) s'affiche dans l'en-tête.
\newtcolorbox{textebox}[1][]{popbase,colframe=popblue,before upper={\popboxhead{Text}{popblue}{#1}\begin{otherlanguage}{english}},after upper={\end{otherlanguage}}}
% Marques ✓ / ✗ (forme correcte / fautive) souvent utilisées par les modèles
\providecommand{\cross}{\textcolor{poppink}{\ensuremath{\times}}}
\providecommand{\xmark}{\cross}\providecommand{\crossmark}{\cross}\providecommand{\cmark}{\ensuremath{\checkmark}}
% Ligne de réponse / justification à compléter (inventée par les modèles)
\providecommand{\justification}[1]{\par\noindent\textit{Justification~:} \dotfill\par}
% \textipa (package tipa non chargé) : la police ne couvre pas l'API -> simple passage
\providecommand{\textipa}[1]{#1}
% Soulignement (ulem non chargé) : \uline, \ul
\providecommand{\uline}[1]{\underline{#1}}\providecommand{\ul}[1]{\underline{#1}}
% \glossaire{\item ...} inventée par les modèles : liste de glossaire
\providecommand{\glossaire}[1]{\begin{itemize}#1\end{itemize}}
% \alert (beamer) inventée par les modèles : texte mis en évidence
\providecommand{\alert}[1]{\textbf{\textcolor{poppink}{#1}}}
% variantes ASCII d'ordinaux écrites en macro
\providecommand{\iere}{\textsuperscript{re}}\providecommand{\ieme}{\textsuperscript{e}}\providecommand{\ere}{\textsuperscript{re}}\providecommand{\eme}{\textsuperscript{e}}\providecommand{\er}{\textsuperscript{er}}\providecommand{\ie}{\textsuperscript{e}}
% Ordinaux français écrits en macro par les modèles : 1\ière, 3\ième
\newcommand{\ière}{\textsuperscript{re}}\newcommand{\ième}{\textsuperscript{e}}
% \exemple (inventée par les modèles) : étiquette « Exemple : »
\providecommand{\exemple}{\textbf{Exemple~:}\ }
% \square utilisable aussi hors mode math (cases à cocher)
\AtBeginDocument{\let\squareMath\square\renewcommand{\square}{\ensuremath{\squareMath}}}
% Mot ou phrase en anglais à l'intérieur d'un paragraphe français
\newcommand{\eng}[1]{\ifmmode\text{\textit{#1}}\else\begingroup\selectlanguage{english}\itshape #1\endgroup\fi}
\let\esp\eng % alias toléré
% flèches d'intonation inventées par les modèles
\providecommand{\textsearrow}{}\renewcommand{\textsearrow}{\ensuremath{\searrow}}\providecommand{\textnearrow}{}\renewcommand{\textnearrow}{\ensuremath{\nearrow}}
% \fr{..} : mot français (inventé par les modèles)
\providecommand{\fr}[1]{\textit{#1}}
\newtcolorbox{exemplebox}[1][]{popbase,colframe=poporange,before upper={\popboxhead{Exemple}{poporange}{#1}}}
\newtcolorbox{definition}[1][]{popbase,colframe=popblue,before upper={\popboxhead{Définition}{popblue}{#1}}}
\newtcolorbox{propriete}[1][]{popbase,colframe=poppurple,before upper={\popboxhead{Propriété}{poppurple}{#1}}}
\newtcolorbox{methode}[1][]{popbase,colframe=popgold,before upper={\popboxhead{Méthode}{popgold}{#1}}}
\newtcolorbox{attention}[1][]{popbase,colframe=poppink,before upper={\popboxhead{Attention}{poppink}{#1}}}
\newtcolorbox{experience}[1][]{popbase,colframe=popblue,colback=popblueL,before upper={\popboxhead{Expérience}{popblue}{#1}}}
% « Le savais-tu ? » : carte violette pleine (contexte, culture, Cameroun)
\newtcolorbox{savaistu}[1][]{popbase,colback=poppurple,colframe=popink,
  fontupper=\popsemi\fontsize{10.4}{14.2}\selectfont\color{white},
  before upper={\poppill{Le savais-tu\,?}{white}{poppurple}\ifx\relax#1\relax\else\hspace{6pt}{\popxbold\color{white}#1}\fi\par\vspace{7pt}}}
% Exercice résolu : énoncé en haut, \tcblower puis la solution sur fond crème
\newcounter{popexo}\newcounter{popexr}
\newtcolorbox{exoresolu}[1][]{popbase,colframe=poporange,colbacklower=popcream,
  segmentation style={popink,line width=1.2pt,dashed},
  before upper={\stepcounter{popexr}\popboxhead{Exercice résolu \thepopexr}{poporange}{#1}},
  before lower={\poppill{Solution}{popink}{popcream}\par\vspace{6pt}}}
% Exercice (non résolu en place) — la correction est dans la partie Corrigés
\newtcolorbox{exercice}[1][]{popbase,colframe=popink,
  before upper={\stepcounter{popexo}\popboxhead{Exercice \thepopexo}{popink}{#1}}}
% Corrigé
\newtcolorbox{corrige}[1][]{popbase,colframe=popgreen,colback=popgreenL,
  before upper={\popboxhead{Corrigé}{popgreen}{#1}}}
% Objectifs / prérequis / bilan
\newtcolorbox{objectifs}{popbase,colframe=popink,colback=poporangeL,
  before upper={\popboxhead{Objectifs de la leçon}{popink}{}}}
\newtcolorbox{prerequis}{popbase,colframe=popink,colback=popblueL,
  before upper={\popboxhead{Prérequis}{popblue}{}}}
\newtcolorbox{bilan}{popbase,colframe=popink,colback=white,boxrule=2.8pt,
  before upper={\popboxhead{Fiche bilan}{poporange}{L'essentiel à connaître pour l'examen}}}
% Figure encadrée avec légende
\newtcolorbox{popfigure}[1][]{enhanced,breakable=false,colback=white,colframe=popline,boxrule=1.4pt,arc=8pt,
  left=10pt,right=10pt,top=10pt,bottom=8pt,before skip=10pt,after skip=12pt,halign=center,
  lower separated=false,halign lower=center,
  fontlower=\popsemi\fontsize{9}{11.5}\selectfont\color{popmuted},
  #1}
\newcommand{\legende}[1]{\par\vspace{6pt}{\popsemi\fontsize{9}{11.5}\selectfont\color{popmuted}#1}}

% ============================================================
% Signature OnBuch+
% ============================================================
\newcommand{\poplogoinline}[1]{{\poplogo\fontsize{#1}{#1}\selectfont\color{popink}OnBuch\color{poporange}+}}
\newcommand{\signatureonbuch}{%
  \begin{tcolorbox}[enhanced,colback=popink,colframe=popink,boxrule=2.2pt,arc=10pt,
      drop shadow=poporange,left=14pt,right=14pt,top=10pt,bottom=10pt,before skip=0pt,after skip=0pt]
    \tikz[baseline=-0.7ex]\node[circle,fill=popgreen,draw=popcream,line width=1.3pt,minimum size=22pt,inner sep=0]{\popcheck{white}};%
    \hspace{9pt}\begin{minipage}[c]{0.62\linewidth}
      {\popxbold\fontsize{12}{13}\selectfont\color{popcream}Signé\ \textbullet\ {\poplogo OnBuch\color{poporange}+}}\par\vspace{2pt}
      {\raggedright\popsemi\fontsize{8}{10.5}\selectfont\color{popline}\MakeUppercase{Équipe pédagogique OnBuch+}\\\MakeUppercase{Conforme au programme officiel MINESEC}\par}
    \end{minipage}\hfill
    {\poplogo\fontsize{22}{22}\selectfont\color{popcream}OnBuch\color{poporange}+}
  \end{tcolorbox}%
}

% ============================================================
% Page de garde
% #1 titre · #2 matière · #3 niveau · #4 module · #5 n° leçon · #6 durée ·
% #7 description / objectifs (texte ou itemize)
% ============================================================
\newcommand{\pagegarde}[7]{%
  \thispagestyle{empty}
  \newgeometry{a4paper,margin=1.7cm,top=1.6cm,bottom=1.4cm}%
  \begin{tikzpicture}[remember picture,overlay]
    \fill[poporange!18] ([xshift=-3.2cm,yshift=-3.6cm]current page.north east) circle (4.6cm);
    \fill[poppurple!14] ([xshift=2.4cm,yshift=7.6cm]current page.south west) circle (3.8cm);
  \end{tikzpicture}%
  {\poplogo\fontsize{30}{30}\selectfont\color{popink}OnBuch\color{poporange}+}\hfill
  \raisebox{0.6ex}{\poppill{Cours signé \textbullet\ Premium}{popink}{popcream}}\par
  \vspace{1pt}\popsquiggle{poppurple}\par
  \vspace{8pt}
  \begin{tcolorbox}[enhanced,colback=poporange,colframe=popink,boxrule=2.8pt,arc=13pt,
      drop shadow=popink,left=18pt,right=18pt,top=12pt,bottom=13pt,after skip=10pt]
    \poppill{#2 \textbullet\ #3}{popink}{popcream}\hspace{5pt}\poppill{#4}{white}{popink}\par\vspace{8pt}
    {\popdisplay\fontsize{14}{14}\selectfont\color{popink}LEÇON #5}\par\vspace{4pt}
    {\raggedright\hyphenpenalty=10000\exhyphenpenalty=10000\popdisplay\fontsize{25}{27}\selectfont\color{popcream}\MakeUppercase{#1}\par}
    \vspace{8pt}
    {\popxbold\fontsize{9.5}{9.5}\selectfont\color{popink}\MakeUppercase{Durée conseillée : #6}}
  \end{tcolorbox}
  \begin{tcolorbox}[enhanced,colback=white,colframe=popink,boxrule=2.2pt,arc=10pt,
      drop shadow=popink,left=16pt,right=16pt,top=10pt,bottom=10pt,after skip=8pt]
    \poppill{Dans cette leçon}{poppurple}{white}\par\vspace{5pt}
    \popsemi\fontsize{9.3}{12.6}\selectfont\color{popink}#7
  \end{tcolorbox}
  \noindent\begin{minipage}[t]{0.487\linewidth}
    \begin{tcolorbox}[enhanced,colback=popgreen,colframe=popink,boxrule=2.2pt,arc=10pt,
        drop shadow=popink,left=11pt,right=11pt,top=10pt,bottom=10pt,height=3.1cm,valign=center]
      \tcbox[colback=white,colframe=popink,boxrule=1.3pt,arc=5pt,boxsep=0pt,nobeforeafter,
        left=3pt,right=3pt,top=3pt,bottom=3pt,valign=center]{\qrcode[height=1.9cm]{https://play.google.com/store/apps/details?id=com.luvvix.onbuch&hl=fr}}%
      \hspace{8pt}\begin{minipage}[c]{0.5\linewidth}
        {\popdisplay\fontsize{11}{12.5}\selectfont\color{white}\MakeUppercase{Télécharge OnBuch}}\par\vspace{4pt}
        {\raggedright\popsemi\fontsize{8.4}{11}\selectfont\color{white}Gratuit \textbullet\ Google Play\\Tuteur IA, annales, concours, résultats\par}
      \end{minipage}
    \end{tcolorbox}
  \end{minipage}\hfill
  \begin{minipage}[t]{0.487\linewidth}
    \begin{tcolorbox}[enhanced,colback=poppurple,colframe=popink,boxrule=2.2pt,arc=10pt,
        drop shadow=popink,left=11pt,right=11pt,top=10pt,bottom=10pt,height=3.1cm,valign=center]
      \tikz[baseline=-0.7ex]\node[circle,fill=white,draw=popink,line width=1.3pt,minimum size=1.6cm,inner sep=0]{\popdisplay\fontsize{24}{24}\selectfont\color{poppurple}+};%
      \hspace{8pt}\begin{minipage}[c]{0.6\linewidth}
        {\popdisplay\fontsize{11}{12.5}\selectfont\color{white}\MakeUppercase{Passe à OnBuch+}}\par\vspace{4pt}
        {\raggedright\popsemi\fontsize{8.4}{11}\selectfont\color{white}Tous les cours signés, Tuteur IA illimité, fiches PDF — onglet OnBuch+ de l'app\par}
      \end{minipage}
    \end{tcolorbox}
  \end{minipage}
  \vspace{8pt}
  \popcontenu
  \vfill
  \signatureonbuch
  \restoregeometry
  \newpage
}

% Rangée « Ce cours contient » (page de garde)
\newcommand{\poptile}[3]{% #1 couleur #2 gros texte #3 légende
  \begin{tcolorbox}[enhanced,colback=#1,colframe=popink,boxrule=2pt,arc=9pt,shadow={2.5pt}{-2.5pt}{0pt}{popink},
      left=6pt,right=6pt,top=9pt,bottom=9pt,halign=center,valign=center,height=1.75cm,width=\linewidth,nobeforeafter]
    {\popdisplay\fontsize{12.5}{13}\selectfont\color{white}\MakeUppercase{#2}}\par\vspace{3pt}
    {\centering\popsemi\fontsize{7.8}{9.5}\selectfont\color{white}#3\par}
  \end{tcolorbox}}
\newcommand{\popcontenu}{%
  \par\noindent{\popxboldls\fontsize{7.5}{7.5}\selectfont\color{popmuted}CE COURS SIGNÉ CONTIENT}\par\vspace{7pt}
  \noindent\begin{minipage}[t]{0.235\linewidth}\poptile{popblue}{Cours}{complet, illustré, pas à pas}\end{minipage}\hfill
  \begin{minipage}[t]{0.235\linewidth}\poptile{poporange}{Méthodes}{et exercices résolus}\end{minipage}\hfill
  \begin{minipage}[t]{0.235\linewidth}\poptile{poppink}{Exercices}{type examen + corrigés}\end{minipage}\hfill
  \begin{minipage}[t]{0.235\linewidth}\poptile{popgreen}{Fiche bilan}{l'essentiel à retenir}\end{minipage}\par}

% Sommaire
\newtcolorbox{sommaire}{enhanced,colback=white,colframe=popink,boxrule=2.2pt,arc=10pt,
  drop shadow=popink,left=18pt,right=18pt,top=13pt,bottom=13pt,before skip=6pt,after skip=20pt,
  before upper={\poppill{Sommaire}{poppurple}{white}\par\vspace{9pt}\popsemi\fontsize{10.6}{14.5}\selectfont\color{popink}}}

% Signature de fin de cours — poussée tout en bas de la dernière page
\newcommand{\findecours}{%
  \par\vfill\nopagebreak
  \begin{center}\popsquiggle{poporange}\end{center}\vspace{4pt}
  \signatureonbuch
}
\AtBeginDocument{\def\llbracket{\mathopen{[\mkern-3.5mu[}}\def\rrbracket{\mathclose{]\mkern-3.5mu]}}} % intervalles d'entiers (glyphe ⟦ absent de la police)
% Glyphes absents de Fira Math : redéfinis à partir de glyphes disponibles
\AtBeginDocument{%
  \def\setminus{\mathbin{\backslash}}\let\smallsetminus\setminus
  \def\vdots{\mathord{\vbox{\baselineskip3.2pt\lineskiplimit0pt\kern4pt\hbox{.}\hbox{.}\hbox{.}}}}%
  \def\ddots{\mathinner{\mkern1mu\raise6.5pt\hbox{.}\mkern2mu\raise3.6pt\hbox{.}\mkern2mu\raise0.7pt\hbox{.}\mkern1mu}}%
  \def\triangle{\mathord{\scalebox{0.9}{$\symup{\Delta}$}}}%
  \def\nabla{\mathord{\raisebox{\depth}{\scalebox{1}[-1]{$\symup{\Delta}$}}}}%
  \def\checkmark{\popcheck{popdark}}%
  \def\star{\mathbin{\ast}}\def\bigstar{\mathord{\ast}}%
  \def\diamond{\mathbin{\scalebox{0.6}{$\Diamond$}}}%
  \def\bigcup{\mathop{\mathchoice{\raisebox{-0.3em}{\scalebox{1.7}{$\cup$}}}{\scalebox{1.25}{$\cup$}}{\cup}{\cup}}\displaylimits}%
  \def\bigcap{\mathop{\mathchoice{\raisebox{-0.3em}{\scalebox{1.7}{$\cap$}}}{\scalebox{1.25}{$\cap$}}{\cap}{\cap}}\displaylimits}%
  \def\lesseqqgtr{\mathrel{\lessgtr}}\def\bigoplus{\mathop{\oplus}\displaylimits}%
}
\providecommand{\ssi}{si et seulement si} % abréviation fréquente
\AtBeginDocument{\providecommand{\wideparen}{\overparen}} % arc de cercle

%% FIGFIX-BEGIN v5 — correctifs globaux des figures (docs/onbuch-generation/tools/figfix)
\usepackage{adjustbox}\usepackage{etoolbox}\tcbuselibrary{raster}
% toute figure TikZ/pgfplots plus large que la ligne est réduite à la largeur disponible
\BeforeBeginEnvironment{tikzpicture}{\begin{adjustbox}{max width=\linewidth}}
\AfterEndEnvironment{tikzpicture}{\end{adjustbox}}
% légendes pgfplots lisibles par-dessus les courbes
\pgfplotsset{every axis/.append style={legend style={fill=white,fill opacity=0.92,text opacity=1,draw=black!15,inner sep=3pt}}}
% étiquettes d'arêtes (syntaxe edge["..."]) avec fond blanc
\tikzset{every edge quotes/.append style={fill=white,inner sep=1.2pt}}
%% FIGFIX-END
\begin{document}

\pagegarde{Lesson 53 — General revision: Revise previous aspects of grammar}{Anglais}{Tle A}{Chapter 5 : Media and Communication}{EN53}{2 h}{Cette leçon de révision générale synthétise l'ensemble des points de grammaire et de vocabulaire étudiés en Terminale A (temps verbaux, modaux, voix passive, discours rapporté, conditionnels, relatifs) et propose un entraînement progressif de type Bac : compréhension, grammaire, vocabulaire et composition.
\vspace{4pt}
\begin{multicols}{2}\small
\begin{itemize}
\item À la fin de cette leçon, je sais choisir le temps verbal correct selon le contexte (present, past, future, perfect, continuous).
\item Je sais employer les auxiliaires modaux (must, should, could, might...) avec leurs nuances de sens.
\item Je sais transformer une phrase active en phrase passive et inversement.
\item Je sais rapporter des paroles et des questions au discours indirect avec les changements nécessaires.
\item Je sais construire les trois types de phrases conditionnelles et les souhaits avec 'wish'.
\item Je sais relier des idées avec les pronoms relatifs et les connecteurs logiques.
\end{itemize}\end{multicols}}

\begin{sommaire}
\begin{multicols}{2}
\begin{enumerate}
\item Vue d'ensemble : la carte des temps verbaux
\item Modaux et semi-auxiliaires
\item Voix passive et discours rapporté
\item Conditionnels, souhaits et phrases complexes
\item Pièges des francophones : le grand nettoyage
\item Entraînement type Bac : méthode et stratégies
\item Activité d'intégration
\item Méthodes et exercices résolus
\item Exercices
\item Corrigés des exercices
\item Fiche bilan
\end{enumerate}
\end{multicols}
\end{sommaire}

\begin{objectifs}
\begin{itemize}
\item À la fin de cette leçon, je sais choisir le temps verbal correct selon le contexte (present, past, future, perfect, continuous).
\item Je sais employer les auxiliaires modaux (must, should, could, might...) avec leurs nuances de sens.
\item Je sais transformer une phrase active en phrase passive et inversement.
\item Je sais rapporter des paroles et des questions au discours indirect avec les changements nécessaires.
\item Je sais construire les trois types de phrases conditionnelles et les souhaits avec 'wish'.
\item Je sais relier des idées avec les pronoms relatifs et les connecteurs logiques.
\item Je sais repérer et corriger les erreurs typiques des francophones dans une phrase.
\item Je sais gérer mon temps et ma méthode face à une épreuve complète de type Bac.
\end{itemize}
\end{objectifs}

\begin{prerequis}
\begin{itemize}
\item Les temps de base étudiés depuis la classe de Première (present simple/continuous, past simple, present perfect, futurs)
\item Les modaux et semi-auxiliaires vus dans les chapitres 1 à 4
\item La voix passive et le discours rapporté (leçons de l'année)
\item Le vocabulaire thématique des chapitres 1 à 5 (health, environment, technology, media...)
\end{itemize}
\end{prerequis}

\newpage

\coursec{Vue d'ensemble : la carte des temps verbaux}

\begin{textebox}[Exam fever in Buea]
The Baccalauréat is only a few weeks away. In a classroom in Buea, a student opens her English paper and panics: should she write “I have seen him yesterday” or “I saw him yesterday”? Like her, many candidates lose easy marks not because they ignore the rules, but because they have never organised them. This final lesson is your revision toolbox: it gathers all the grammar points of the year — tenses, modals, passive voice, reported speech, conditionals, relative clauses — into one clear system, with exam-type training.
\end{textebox}

\textbf{Problématique.} Comment choisir le bon temps verbal en quelques secondes, le jour de l'examen ? Réponse de cette section : il faut d'abord posséder une \cle{carte mentale} complète des temps anglais — trois zones temporelles (present, past, future) croisées avec quatre aspects (simple, continuous, perfect, perfect continuous). C'est cette carte que nous allons construire ensemble, temps par temps, avec le contexte du chapitre : \emph{Media and Communication}.

\courssub{Les trois zones temporelles et les quatre aspects}

\begin{definition}[Tense et aspect]
En anglais, tout verbe conjugué porte deux informations : le \cle{temps} (où se situe l'action : present, past, future) et l'\cle{aspect} (comment l'action est vue : simple = fait global ou habitude ; continuous = en cours ; perfect = bilan ou antériorité ; perfect continuous = durée jusqu'à un point de repère).
\end{definition}

\begin{propriete}[Les trois questions pour choisir un temps]
Avant de conjuguer, pose-toi trois questions :
\begin{enumerate}
\item \textbf{QUAND ?} L'action se situe-t-elle dans le présent, le passé ou le futur ?
\item \textbf{FINIE ou EN COURS ?} Action terminée / habitude → aspect \emph{simple} ; action en train de se dérouler → aspect \emph{continuous}.
\item \textbf{LIEN avec le présent ou ANTÉRIORITÉ ?} Bilan, expérience, résultat visible maintenant → \emph{perfect} ; action antérieure à une autre action passée → \emph{past perfect}.
\end{enumerate}
\end{propriete}

\begin{center}
\begin{tabularx}{\linewidth}{|l|X|X|X|}
\hline
\rowcolor{popblueL} Aspect & Present & Past & Future \\
\hline
Simple & She reports the news. & She reported the news. & She will report the news. \\
\hline
Continuous & She is reporting live. & She was reporting live. & She will be reporting live. \\
\hline
Perfect & She has reported the news. & She had reported the news. & She will have reported the news. \\
\hline
Perfect continuous & She has been reporting since 6 a.m. & She had been reporting for hours. & She will have been reporting for hours. \\
\hline
\end{tabularx}
\end{center}

\begin{popfigure}
\begin{tikzpicture}[scale=1.0]
\draw[-{Stealth}, thick, popdark] (-5.5,0) -- (5.5,0);
\node[above, popdark, font=\bfseries] at (-4.5,0) {PAST};
\node[above, popdark, font=\bfseries] at (0,0) {NOW};
\node[above, popdark, font=\bfseries] at (4.5,0) {FUTURE};
\fill[popdark] (0,0) circle (0.08);
\draw[-{Stealth}, very thick, poppurple] (-4.6,-0.9) -- (-3.4,-0.9);
\node[below, poppurple, align=center] at (-4.0,-1.0) {past perfect\\ (had started)};
\draw[-{Stealth}, very thick, popblue] (-3.2,-0.9) -- (-1.6,-0.9);
\node[below, popblue, align=center] at (-2.4,-1.0) {past simple\\ (started, yesterday)};
\draw[-{Stealth}, very thick, popgreen] (-1.5,-0.9) -- (0,-0.9);
\node[below, popgreen, align=center] at (-0.8,-1.0) {present perfect\\ (has just announced)};
\draw[-{Stealth}, very thick, poporange] (0.3,-0.9) -- (2.2,-0.9);
\node[below, poporange, align=center] at (1.2,-1.0) {will / going to\\ (will broadcast)};
\draw[-{Stealth}, very thick, poppink] (2.5,-0.9) -- (4.6,-0.9);
\node[below, poppink, align=center] at (3.5,-1.0) {future perfect\\ (will have ended)};
\end{tikzpicture}
\legende{La frise des temps : le past perfect se situe en arrière du past simple ; le present perfect aboutit à NOW ; will et going to se projettent après NOW.}
\end{popfigure}

\begin{popfigure}
\begin{center}\tcbox[colback=poporange!70,colframe=poporange!70,coltext=white,arc=9pt,boxsep=4pt,left=6pt,right=6pt]{\bfseries\small TENSES}\end{center}
\vspace{-2pt}
\begin{tcbraster}[raster columns=2,raster equal height=rows,raster column skip=6pt,raster row skip=6pt,enhanced,blankest]
\begin{tcolorbox}[enhanced,colback=white,colframe=popgreen!60,boxrule=0.9pt,arc=3pt,title={present},fonttitle=\bfseries\scriptsize,coltitle=white,colbacktitle=popgreen!60,left=4pt,right=4pt,top=2pt,bottom=2pt,boxsep=1pt,toptitle=1.5pt,bottomtitle=1.5pt,halign=left]
\begin{itemize}[nosep,leftmargin=*,itemsep=1pt,font=\scriptsize]
\item simple
\item continuous
\item perfect
\end{itemize}
\end{tcolorbox}
\begin{tcolorbox}[enhanced,colback=white,colframe=popblue!50,boxrule=0.9pt,arc=3pt,title={past},fonttitle=\bfseries\scriptsize,coltitle=white,colbacktitle=popblue!50,left=4pt,right=4pt,top=2pt,bottom=2pt,boxsep=1pt,toptitle=1.5pt,bottomtitle=1.5pt,halign=left]
\begin{itemize}[nosep,leftmargin=*,itemsep=1pt,font=\scriptsize]
\item simple
\item continuous
\item perfect
\end{itemize}
\end{tcolorbox}
\begin{tcolorbox}[enhanced,colback=white,colframe=poppurple!50,boxrule=0.9pt,arc=3pt,title={future},fonttitle=\bfseries\scriptsize,coltitle=white,colbacktitle=poppurple!50,left=4pt,right=4pt,top=2pt,bottom=2pt,boxsep=1pt,toptitle=1.5pt,bottomtitle=1.5pt,halign=left]
\begin{itemize}[nosep,leftmargin=*,itemsep=1pt,font=\scriptsize]
\item will
\item going to
\item present cont.
\end{itemize}
\end{tcolorbox}
\end{tcbraster}

\legende{Carte mentale des temps : trois branches temporelles, chacune subdivisée selon l'aspect.}
\end{popfigure}

\courssub{Le présent : present simple et present continuous}

\begin{exemplebox}[Observer avant de conclure]
\eng{Journalists report facts every day.} « Les journalistes rapportent les faits chaque jour. »

\eng{The BBC is broadcasting live tonight.} « La BBC diffuse en direct ce soir. »

Observe : dans la première phrase, \eng{every day} signale une habitude → \emph{present simple}. Dans la seconde, \eng{tonight} signale un événement ponctuel programmé → \emph{present continuous}.
\end{exemplebox}

\begin{propriete}[Present simple — forme et emplois]
\textbf{Forme :} base verbale ; à la 3\textsuperscript{e} personne du singulier, on ajoute \textbf{-s} ou \textbf{-es} (\eng{he watches}, \eng{she reports}). Négation et question avec \textbf{do/does} : \eng{She does not work on Sundays.} / \eng{Does he read the newspaper?}

\textbf{Emplois :}
\begin{itemize}
\item habitudes et routines : \eng{My father listens to CRTV radio every morning.} « Mon père écoute la radio CRTV chaque matin. »
\item vérités générales et faits permanents : \eng{The sun rises in the east.} « Le soleil se lève à l'est. »
\item horaires officiels (programmes, transports) : \eng{The news bulletin starts at 8 p.m.} « Le journal télévisé commence à 20 h. »
\end{itemize}
\textbf{Marqueurs typiques :} always, usually, often, every day, on Mondays, never.
\end{propriete}

\begin{propriete}[Present continuous — forme et emplois]
\textbf{Forme :} \textbf{be} (am/is/are) + verbe en \textbf{-ing} : \eng{I am reading}, \eng{she is writing}, \eng{they are filming}.

\textbf{Emplois :}
\begin{itemize}
\item action en cours au moment où l'on parle : \eng{Listen! The presenter is speaking.} « Écoute ! Le présentateur est en train de parler. »
\item action temporaire : \eng{She is staying in Douala this month.} « Elle séjourne à Douala ce mois-ci. »
\item projet futur organisé : \eng{We are interviewing the mayor tomorrow.} « Nous interviewons le maire demain. »
\end{itemize}
\textbf{Exception majeure :} les verbes d'état (know, believe, like, want, understand, own) ne prennent presque jamais la forme continue : \eng{I know the answer.} et non \emph{I am knowing}.
\end{propriete}

\begin{attention}[Le présent continu n'existe pas en français]
Le français dit « je lis » pour les deux situations. L'anglais distingue : \eng{I read the newspaper every morning} (habitude) et \eng{I am reading the newspaper now} (en ce moment). Ne traduis jamais « je suis en train de » mot à mot : \emph{I am in train of reading} est impossible.
\end{attention}

\courssub{Le passé : past simple, past continuous, past perfect}

\begin{exemplebox}[Trois phrases, trois temps]
\eng{The reporter interviewed the minister yesterday.} « Le journaliste a interviewé le ministre hier. »

\eng{I was watching the news when the phone rang.} « Je regardais les informations quand le téléphone a sonné. »

\eng{The programme had started before we arrived.} « L'émission avait commencé avant que nous arrivions. »
\end{exemplebox}

\begin{propriete}[Past simple — action passée, datée, terminée]
\textbf{Forme :} verbe + \textbf{-ed} (réguliers) ou 2\textsuperscript{e} colonne des verbes irréguliers (go → went, see → saw, write → wrote). Négation et question avec \textbf{did} + base verbale : \eng{Did you watch the debate?} / \eng{She did not watch it.}

\textbf{Emploi :} action achevée dans un passé \emph{coupé du présent}, souvent avec un repère de temps : yesterday, last week, in 2019, two days ago.

\eng{The station broadcast the match last Saturday.} « La chaîne a diffusé le match samedi dernier. »
\end{propriete}

\begin{propriete}[Past continuous — action en cours dans le passé]
\textbf{Forme :} \textbf{was/were} + verbe en \textbf{-ing}.

\textbf{Emplois :}
\begin{itemize}
\item action longue interrompue par une action courte (past simple) : \eng{While I was listening to the radio, the electricity went off.} « Pendant que j'écoutais la radio, l'électricité a été coupée. »
\item toile de fond, description : \eng{The crowd was cheering and the cameras were filming.} « La foule acclamait et les caméras filmaient. »
\end{itemize}
\textbf{Marqueurs :} while, when, as, at 8 o'clock yesterday.
\end{propriete}

\begin{propriete}[Past perfect — l'antériorité]
\textbf{Forme :} \textbf{had} + participe passé.

\textbf{Emploi :} une action \emph{antérieure} à une autre action passée. C'est le « passé du passé ».

\eng{When we arrived at the studio, the interview had already started.} « Quand nous sommes arrivés au studio, l'interview avait déjà commencé. »

\textbf{Marqueurs :} before, after, by the time, already, when.
\end{propriete}

\begin{center}
\begin{tabularx}{\linewidth}{|X|X|X|}
\hline
\rowcolor{popblueL} Français & English & Remarque \\
\hline
Je regardais la télé quand il est arrivé. & I was watching TV when he arrived. & longue action : continuous ; courte : simple \\
\hline
Elle avait déjà publié l'article. & She had already published the article. & antériorité : past perfect \\
\hline
Il a quitté Bamenda en 2020. & He left Bamenda in 2020. & date précise : past simple \\
\hline
\end{tabularx}
\end{center}

\courssub{Le present perfect : le pont entre passé et présent}

\begin{propriete}[Present perfect — forme et emplois]
\textbf{Forme :} \textbf{have/has} + participe passé : \eng{I have read}, \eng{she has written}.

\textbf{Emplois :}
\begin{itemize}
\item bilan, résultat visible maintenant : \eng{The government has just announced new media laws.} « Le gouvernement vient d'annoncer de nouvelles lois sur les médias. »
\item expérience de vie (sans date) : \eng{Have you ever visited a TV studio?} « As-tu déjà visité un studio de télévision ? »
\item action commencée dans le passé et qui continue : \eng{She has worked for CRTV for ten years.} « Elle travaille à la CRTV depuis dix ans. »
\end{itemize}
\textbf{Marqueurs :} just, already, yet, ever, never, since, for, recently, so far.
\end{propriete}

\begin{exemplebox}[Present perfect ou past simple ? La différence de sens]
\eng{She has worked for CRTV for ten years.} « Elle travaille à la CRTV depuis dix ans. » → elle y travaille \emph{encore}.

\eng{She worked for CRTV for ten years, then retired.} « Elle a travaillé à la CRTV pendant dix ans, puis a pris sa retraite. » → c'est \emph{terminé}.

Le temps choisi change le sens : le present perfect relie au présent, le past simple coupe du présent.
\end{exemplebox}

\begin{attention}[L'erreur n° 1 des francophones : since et for]
Le français utilise le présent : « Je suis ici depuis 2020. » L'anglais \textbf{exige} le present perfect :

\eng{I am here since 2020.} ✗ → \eng{I have been here since 2020.} ✓

De même : \eng{I have seen him yesterday.} ✗ → \eng{I saw him yesterday.} ✓ (un repère passé comme \emph{yesterday} interdit le present perfect). Retiens : \textbf{since + point de départ} (since 2020, since Monday) ; \textbf{for + durée} (for ten years, for a long time).
\end{attention}

\courssub{Le futur : will, going to, present continuous}

\begin{propriete}[Trois façons de parler du futur]
\begin{itemize}
\item \textbf{will + base verbale} : décision spontanée, promesse, prédiction d'opinion. \eng{I think the new channel will be very popular.} « Je pense que la nouvelle chaîne sera très populaire. » / \eng{— The phone is ringing! — I will answer it.} « Le téléphone sonne ! Je vais répondre. »
\item \textbf{be going to + base verbale} : projet déjà décidé, ou signe visible. \eng{We are going to launch a school newspaper.} « Nous allons lancer un journal scolaire. » / \eng{Look at those clouds! It is going to rain.} « Regarde ces nuages ! Il va pleuvoir. »
\item \textbf{present continuous} : rendez-vous fixé, agenda précis. \eng{The BBC is broadcasting live tonight.} « La BBC diffuse en direct ce soir. » / \eng{I am meeting the editor at 3 p.m.} « Je rencontre le rédacteur en chef à 15 h. »
\end{itemize}
\end{propriete}

\begin{attention}[Pas de will après when, if, as soon as]
Dans les subordonnées de temps et de condition, le futur se rend par le présent : \eng{When the programme ends, I will call you.} « Quand l'émission se terminera, je t'appellerai. » — et non \emph{when the programme will end}.
\end{attention}

\courssub{Méthode et entraînement}

\begin{methode}[Choisir le bon temps à l'examen, en quatre étapes]
\begin{enumerate}
\item \textbf{Souligne les indices de temps} dans la phrase : yesterday, last year → past simple ; since, for, already, just, yet → present perfect ; while, when (action longue) → continuous ; by the time, before → past perfect ; tomorrow, next week → futur.
\item \textbf{Pose les trois questions} : quand ? fini ou en cours ? lien avec le présent ou antériorité ?
\item \textbf{Vérifie les verbes d'état} : know, believe, want, like → jamais de forme continue.
\item \textbf{Relis} : auxiliaire correct (do/does/did/have/has/had/be) + forme du verbe (base, -ing, participe passé).
\end{enumerate}
\end{methode}

\begin{exoresolu}[Exercice résolu : mettre le verbe au bon temps]
Énoncé. Complete the sentences with the correct tense of the verb in brackets.

1. The reporter \dots{} (interview) the minister yesterday.

2. She \dots{} (work) for CRTV since 2015.

3. I \dots{} (watch) the news when the lights went off.

4. By the time we arrived, the debate \dots{} (start).

5. Look! The cameraman \dots{} (film) the crowd.

6. We \dots{} (visit) the new media centre next Friday. (agenda fixé)

\tcblower
Solution détaillée.

1. \eng{The reporter interviewed the minister yesterday.} — l'indice \emph{yesterday} impose le past simple (interview → interviewed).

2. \eng{She has worked for CRTV since 2015.} — \emph{since + point de départ} et action qui continue → present perfect (has + worked).

3. \eng{I was watching the news when the lights went off.} — action longue interrompue par une action courte → past continuous (was watching) puis past simple (went off).

4. \eng{By the time we arrived, the debate had started.} — \emph{by the time} signale l'antériorité → past perfect (had started).

5. \eng{Look! The cameraman is filming the crowd.} — \emph{Look!} annonce une action en cours maintenant → present continuous (is filming).

6. \eng{We are visiting the new media centre next Friday.} — rendez-vous fixé à l'agenda → present continuous à valeur de futur (are visiting).
\end{exoresolu}

\begin{exercice}[À toi de jouer]
Put the verbs in brackets into the correct tense.

1. Journalists \dots{} (check) their sources before publishing.

2. He \dots{} (never / see) a live broadcast before.

3. While she \dots{} (type) her article, her computer crashed.

4. They \dots{} (already / announce) the results when we switched on the TV.

5. I think social media \dots{} (change) journalism completely.

6. We \dots{} (live) in Yaoundé for five years now.
\end{exercice}

\begin{corrige}[Exercice « À toi de jouer »]
1. \eng{Journalists check their sources before publishing.} — vérité générale → present simple.

2. \eng{He has never seen a live broadcast before.} — expérience sans date, \emph{never/before} → present perfect (see, saw, seen).

3. \eng{While she was typing her article, her computer crashed.} — action longue interrompue → past continuous + past simple.

4. \eng{They had already announced the results when we switched on the TV.} — antériorité (\emph{already} + \emph{when}) → past perfect.

5. \eng{I think social media will change journalism completely.} — prédiction d'opinion après \emph{I think} → will + base verbale.

6. \eng{We have lived in Yaoundé for five years now.} — \emph{for + durée}, action qui continue → present perfect. Attention : pas de présent simple avec \emph{for five years}.
\end{corrige}

\begin{aretenir}[L'essentiel de la carte des temps]
\begin{itemize}
\item Trois zones : present, past, future ; quatre aspects : simple, continuous, perfect, perfect continuous.
\item Present simple = habitude, vérité ; present continuous = en cours ou projet fixé.
\item Past simple = passé daté et terminé ; past continuous = action longue interrompue ; past perfect = avant une autre action passée.
\item Present perfect = bilan, expérience, action qui continue (since/for) ; jamais avec \emph{yesterday}.
\item Futur : will (décision, prédiction), going to (projet, signe visible), present continuous (agenda).
\item À l'examen : souligne les indices, pose les trois questions, vérifie l'auxiliaire.
\end{itemize}
\end{aretenir}

\coursec{Modaux et semi-auxiliaires}

Dans la section précédente, nous avons dressé la carte complète des temps verbaux : chaque temps situe l'action dans le temps. Mais un verbe ne dit pas seulement \emph{quand} : il dit aussi si l'action est \emph{obligatoire}, \emph{interdite}, \emph{possible}, \emph{conseillée} ou \emph{probable}. C'est précisément le rôle des \cle{modaux} (\eng{modal verbs}) : de petits auxiliaires — \eng{must, can, should, may, might} — qui colorent le sens du verbe principal. Dans le domaine des médias, ils sont omniprésents : \eng{Journalists must check their sources. Readers should compare several newspapers. The news might be fake.} Observons d'abord un texte authentique, relevons-y tous les modaux, puis tirons-en les règles.

\begin{textebox}[A letter to the editor — The Daily Post, Douala]
Dear Editor,

I am writing about the way news spreads on social media today. Every morning, thousands of Cameroonians receive videos and messages on WhatsApp, and many of them believe everything they read. This must stop. We cannot simply forward every shocking story we see: we must check the source first. It could be a rumour, or even deliberate disinformation.

Journalists should protect their sources, but they mustn't invent facts. Last month, a popular blog published a story about a school in Bamenda that had supposedly closed. The editor should have verified the information before publishing it, because the story was completely false. The school was open, and the students were sitting their exams!

Readers don't have to buy a newspaper to be informed — much reliable news is free online — but they ought to choose serious websites. May I suggest one simple rule? Before you share a story, ask yourself three questions: Who wrote it? When? Can I find it elsewhere? If you can't answer, you'd better not forward it.

The truth might be slower than the rumour, but it is always stronger.

Yours faithfully,

Marcel E., Yaoundé
\end{textebox}

\textbf{Glossaire.} \emph{to spread} (v.) : se répandre ; \emph{to forward} (v.) : transférer ; \emph{a rumour} (n.) : une rumeur ; \emph{disinformation} (n.) : la désinformation ; \emph{a source} (n.) : une source ; \emph{supposedly} (adv.) : soi-disant ; \emph{to verify} (v.) : vérifier ; \emph{reliable} (adj.) : fiable ; \emph{to share} (v.) : partager.

\textbf{Questions de compréhension.} 1) What must readers do before forwarding a story? 2) What mistake did the blog editor make? 3) According to Marcel, do readers have to pay to be well informed? 4) What three questions should you ask yourself before sharing? 5) Find in the text one modal expressing prohibition and one expressing regret about the past.

\courssub{La structure commune à tous les modaux}

Avant d'étudier chaque valeur, retenons la grammaire commune : tous les modaux obéissent aux mêmes règles mécaniques. Ce sont des \emph{auxiliaires}, pas des verbes ordinaires.

\begin{propriete}[Les quatre lois des modaux]
\begin{enumerate}
\item Tout modal est suivi de la \cle{base verbale} (l'infinitif sans \eng{to}) : \eng{He must go.} et jamais \eng{He musts go} ni \eng{He must to go}.
\item Le modal est \cle{invariable} : pas de \eng{-s} à la troisième personne du singulier. \eng{She can swim.} (jamais \eng{she cans}).
\item La négation se fait \cle{directement}, sans \eng{do} : \eng{must not (mustn't), cannot (can't), should not (shouldn't), may not, might not}.
\item L'interrogation se fait par \cle{inversion}, sans \eng{do} : \eng{Must I come? Can you help me? Should we wait?}
\end{enumerate}
\end{propriete}

\begin{attention}[Erreurs typiques des francophones]
Le français dit « il doit partir », « elle peut nager » : la tentation est grande de conjuguer le modal comme un verbe ordinaire. Dites \eng{He must leave.} et non \eng{He musts leave.} ; \eng{She can swim.} et non \eng{She can to swim.} De même, l'interrogation \eng{Do you can swim?} est impossible : dites \eng{Can you swim?} Enfin, \eng{cannot} s'écrit en \textbf{un seul mot} (forme contractée : \eng{can't}, prononcé « kânte » en anglais britannique).
\end{attention}

\courssub{Obligation, interdiction et absence d'obligation}

\begin{definition}[must, have to, mustn't, don't have to]
\eng{Must} exprime une obligation ressentie par le locuteur (ou une règle forte) ; \eng{have to} exprime une obligation extérieure (loi, règlement, circonstances). \eng{Mustn't} exprime l'\cle{interdiction} ; \eng{don't have to} exprime l'\cle{absence d'obligation}.
\end{definition}

\begin{exemplebox}[Observer et comparer]
\eng{You must check your sources before publishing.} « Tu dois vérifier tes sources avant de publier. » (principe moral du journaliste)\par
\eng{Students have to wear a uniform at our school.} « Les élèves doivent porter l'uniforme dans notre lycée. » (règlement extérieur)\par
\eng{You mustn't smoke here.} « Il est interdit de fumer ici. »\par
\eng{You don't have to buy the paper; it's free online.} « Tu n'es pas obligé d'acheter le journal ; il est gratuit en ligne. »
\end{exemplebox}

\begin{aretenir}[La nuance capitale : mustn't $\neq$ don't have to]
\eng{You mustn't smoke here.} = il est \textbf{interdit} de fumer ici (ne fume pas !).\par
\eng{You don't have to smoke here.} = tu n'es pas \textbf{obligé} de fumer ici (c'est libre).\par
En français, « il ne faut pas » traduit \eng{mustn't}, tandis que « tu n'es pas obligé de » traduit \eng{don't have to}. Confondre les deux change complètement le sens ! À l'oral, \eng{mustn't} se prononce « meu-seunt » (le premier \eng{t} est muet).
\end{aretenir}

\begin{propriete}[Conjugaison de have to]
\eng{Have to} est un \cle{semi-auxiliaire} : il se conjugue comme un verbe ordinaire, avec \eng{do} à la négation et à l'interrogation.
\begin{center}
\begin{tabular}{|l|l|l|}
\hline
\rowcolor{popblueL} Forme & Exemple & Traduction \\
\hline
Affirmative & She has to leave early. & Elle doit partir tôt. \\
\hline
Négative & She doesn't have to leave. & Elle n'est pas obligée de partir. \\
\hline
Interrogative & Does she have to leave? & Est-ce qu'elle doit partir ? \\
\hline
Passé & She had to leave. & Elle a dû partir. \\
\hline
Futur & She will have to leave. & Elle devra partir. \\
\hline
\end{tabular}
\end{center}
Remarque : \eng{must} n'a ni passé ni futur ; on utilise alors \eng{had to} et \eng{will have to}. \eng{Yesterday, I had to finish my article before six o'clock.} « Hier, j'ai dû finir mon article avant six heures. »
\end{propriete}

\begin{definition}[needn't : une autre façon de dire « pas obligé »]
\eng{Needn't} + base verbale équivaut à \eng{don't have to} : \eng{You needn't buy the paper; it's free online.} « Tu n'es pas obligé d'acheter le journal ; il est gratuit en ligne. » Attention : \eng{needn't} est un vrai modal (invariable, sans \eng{do}), alors que \eng{need} verbe ordinaire se construit avec \eng{to} : \eng{You need to check your sources.} « Tu dois vérifier tes sources. »
\end{definition}

\courssub{Le conseil : should, ought to, had better}

\begin{definition}[should, ought to, had better]
\eng{Should} et \eng{ought to} expriment le \cle{conseil}, la recommandation (« tu devrais »). \eng{Had better} exprime un conseil plus \cle{pressant}, avec une idée d'avertissement (« tu ferais bien de »).
\end{definition}

\begin{exemplebox}[Exemples traduits]
\eng{Journalists should protect their sources.} « Les journalistes devraient protéger leurs sources. »\par
\eng{You ought to read more than one newspaper.} « Tu devrais lire plus d'un journal. »\par
\eng{You'd better not forward that message.} « Tu ferais bien de ne pas transférer ce message. »
\end{exemplebox}

\begin{propriete}[Formes]
\eng{Should} et \eng{ought to} sont suivis de la base verbale ; \eng{ought to} garde son \eng{to} : \eng{You ought to go.} Négations : \eng{shouldn't}, \eng{oughtn't to} (rare), \eng{had better not}. Attention : après \eng{had better}, la base verbale \emph{sans} \eng{to} : \eng{You'd better hurry.} « Tu ferais bien de te dépêcher. »
\end{propriete}

\begin{savaistu}[You'd better hurry!]
À l'oral britannique, \eng{had better} se contracte presque toujours : \eng{You'd better hurry!} « Tu ferais bien de te dépêcher ! » La contraction \eng{'d} représente ici \eng{had} et non \eng{would} — le contexte (la base verbale qui suit) permet de les distinguer : \eng{You'd better go} (\eng{had better} + \eng{go}) contre \eng{You'd like to go} (\eng{would like}).
\end{savaistu}

\courssub{Capacité, possibilité et permission}

\begin{definition}[can, could, be able to, may, might]
\eng{Can} exprime la \cle{capacité} (savoir-faire) et la possibilité générale ; \eng{could} est son passé (ou sa forme atténuée) ; \eng{be able to} remplace \eng{can} aux temps qui lui manquent ; \eng{may} et \eng{might} expriment la possibilité incertaine (« il se peut que »).
\end{definition}

\begin{exemplebox}[Exemples traduits]
\eng{She can speak three languages.} « Elle sait parler trois langues. » (capacité)\par
\eng{When I was ten, I could read English comics.} « À dix ans, je savais lire les bandes dessinées anglaises. » (capacité passée)\par
\eng{The news might be fake.} « Cette nouvelle est peut-être fausse. » (possibilité)\par
\eng{It may rain during the match in Buea.} « Il se peut qu'il pleuve pendant le match à Buea. »
\end{exemplebox}

\begin{attention}[can n'a ni futur ni participe passé]
\eng{Can} est un verbe \cle{défectif} : il n'existe qu'au présent (\eng{can}) et au passé (\eng{could}). Pour le futur et les temps composés, utilisez \eng{be able to} : \eng{One day, you will be able to read Shakespeare.} « Un jour, tu pourras lire Shakespeare. » ; \eng{I have been able to contact the journalist.} « J'ai pu contacter le journaliste. » Ne dites jamais \eng{I will can}.
\end{attention}

\begin{definition}[La permission : can, could, may]
Pour demander la permission : \eng{Can I...?} (familier), \eng{Could I...?} (poli), \eng{May I...?} (très poli, formel). \eng{May I ask a question, Sir?} « Puis-je poser une question, Monsieur ? » Pour accorder la permission : \eng{You can / may use my phone.} ; pour la refuser : \eng{No, you can't.} ou, plus ferme, \eng{No, you may not.}
\end{definition}

\courssub{La déduction logique : l'échelle de la certitude}

Les modaux servent aussi à exprimer ce que le locuteur \emph{déduit} d'une situation, à partir d'indices. C'est l'échelle de la certitude, capitale à l'examen.

\begin{propriete}[Déduction : must, may/might, can't]
\begin{itemize}
\item \eng{Must} = certitude affirmative (« il doit être », c'est sûr) : \eng{He must be tired; he worked all night.} « Il doit être fatigué ; il a travaillé toute la nuit. »
\item \eng{May / might} = probabilité moyenne (« il se peut ») : \eng{She may be at the studio.} « Elle est peut-être au studio. »
\item \eng{Can't} = impossibilité, certitude négative (« ce n'est pas possible que ») : \eng{He can't be tired; he slept ten hours!} « Il ne peut pas être fatigué ; il a dormi dix heures ! »
\end{itemize}
\end{propriete}

\begin{attention}[Le contraire de « must » en déduction]
En déduction, le contraire de \eng{must} n'est pas \eng{mustn't} mais \eng{can't} ! \eng{He must be ill} « il doit être malade » se nie par \eng{He can't be ill} « il ne peut pas être malade ». \eng{He mustn't be ill} voudrait dire « il n'a pas le droit d'être malade » — absurde. Repérez toujours l'indice qui justifie la déduction : \eng{I saw him at the market}, \eng{he worked all night}.
\end{attention}

\begin{popfigure}
\begin{tikzpicture}
\draw[-{Stealth}, line width=1.2pt, popmuted] (0,0) -- (12.5,0);
\node[below, popmuted] at (6.25,-0.1) {\small degré de certitude croissant};
\foreach \x/\val/\col in {0.6/0\%/poppink, 4.0/50\%/popgold, 7.4/80\%/poporange, 11.4/100\%/popgreen}{
  \fill[\col] (\x,0) circle (0.09);
  \node[above, \col, font=\small\bfseries] at (\x,0.05) {\val};
}
\node[above=0.55cm, align=center, poppink, font=\small] at (0.6,0) {\textbf{can't}\\ \eng{He can't be at work;}\\ \eng{I saw him at the market.}};
\node[above=0.55cm, align=center, popgold, font=\small] at (4.0,0) {\textbf{may / might}\\ \eng{The news might be fake.}};
\node[above=0.55cm, align=center, poporange, font=\small] at (7.4,0) {\textbf{should}\\ \eng{She should be home by now.}};
\node[above=0.55cm, align=center, popgreen, font=\small] at (11.4,0) {\textbf{must}\\ \eng{He must be tired;}\\ \eng{he worked all night.}};
\end{tikzpicture}
\legende{L'échelle de la certitude : de l'impossibilité (\eng{can't}) à la certitude (\eng{must}).}
\end{popfigure}

\courssub{Le regret et la critique du passé : should have + participe passé}

\begin{propriete}[should have + participe passé]
Pour exprimer le \cle{regret} ou la \cle{critique} d'une action passée qui n'a pas eu lieu, on utilise \eng{should have} + participe passé (négation : \eng{shouldn't have}). Cela correspond au conditionnel passé français « aurait dû ».\par
\eng{The editor should have verified the information.} « Le rédacteur aurait dû vérifier l'information. » (il ne l'a pas fait)\par
\eng{You shouldn't have believed that rumour.} « Tu n'aurais pas dû croire cette rumeur. » (tu l'as crue)
\end{propriete}

Même logique avec \eng{could have} (« aurait pu ») : \eng{We could have checked the date.} « Nous aurions pu vérifier la date. » Et avec \eng{might have} (« aurait peut-être ») : \eng{He might have missed the bus.} « Il a peut-être raté le bus. »

\begin{attention}[should have, pas should had]
Après \eng{have}, c'est toujours le \cle{participe passé} : \eng{should have gone}, \eng{should have seen}, \eng{should have written}. À l'oral, \eng{have} se réduit (« choude-ève »), ce qui pousse certains élèves à écrire \eng{should of} — faute grave à l'examen ! Écrivez toujours \eng{should have}.
\end{attention}

\courssub{Tableau récapitulatif}

\begin{center}
\begin{tabularx}{\linewidth}{|l|l|X|X|}
\hline
\rowcolor{popblueL} Modal & Valeur & Exemple & Français \\
\hline
must & obligation / certitude & You must check your sources. & Tu dois vérifier tes sources. \\
\hline
have to & obligation extérieure & Reporters have to respect the law. & Les journalistes doivent respecter la loi. \\
\hline
mustn't & interdiction & You mustn't invent facts. & Il est interdit d'inventer des faits. \\
\hline
don't have to / needn't & absence d'obligation & You don't have to pay; it's free. & Tu n'es pas obligé de payer ; c'est gratuit. \\
\hline
should / ought to & conseil & You should compare sources. & Tu devrais comparer les sources. \\
\hline
had better & conseil pressant & You'd better not share it. & Tu ferais bien de ne pas le partager. \\
\hline
can / could & capacité, permission & Can I use your phone? & Je peux utiliser ton téléphone ? \\
\hline
may / might & possibilité, permission polie & It may be true. / May I come in? & C'est peut-être vrai. / Puis-je entrer ? \\
\hline
can't & impossibilité (déduction) & That can't be true! & Ce n'est pas possible que ce soit vrai ! \\
\hline
should have + p.p. & regret du passé & He should have checked. & Il aurait dû vérifier. \\
\hline
\end{tabularx}
\end{center}

\courssub{Exercice résolu et entraînement}

\begin{exoresolu}[Compléter avec le modal approprié]
Énoncé. Complétez avec \eng{must}, \eng{mustn't}, \eng{don't have to}, \eng{should}, \eng{can't}, \eng{might} ou \eng{should have}.\par
1. You \ldots\ cross the road when the light is red. It's dangerous.\par
2. You \ldots\ buy a television licence to watch YouTube videos; they are free.\par
3. Listen! Someone is knocking at the door. It \ldots\ be the postman; he always comes at this time.\par
4. That story \ldots\ be true — no serious newspaper has reported it.\par
5. The journalist \ldots\ verified his facts before publishing. Now everyone thinks the minister has resigned!
\tcblower
Solution détaillée.\par
1. \eng{mustn't} : il s'agit d'une interdiction (danger). « Il est interdit de traverser au feu rouge. »\par
2. \eng{don't have to} : absence d'obligation (« tu n'es pas obligé d'acheter »), confirmée par \eng{they are free}.\par
3. \eng{must} : déduction, certitude fondée sur un indice (\eng{he always comes at this time}). « Ce doit être le facteur. »\par
4. \eng{can't} : déduction négative, impossibilité. « Cette histoire ne peut pas être vraie. »\par
5. \eng{should have} + participe passé : critique d'une action passée non accomplie. « Le journaliste aurait dû vérifier ses faits. »
\end{exoresolu}

\begin{exercice}[À vous de jouer]
1. Traduisez : a) « Tu n'es pas obligé d'acheter le journal ; il est gratuit en ligne. » b) « Les journalistes devraient protéger leurs sources. » c) « Il doit être fatigué ; il a travaillé toute la nuit. » d) « Le rédacteur aurait dû vérifier l'information. »\par
2. Corrigez les erreurs : a) \eng{He musts go now.} b) \eng{Do you can speak English?} c) \eng{I will can help you tomorrow.} d) \eng{She must to see a doctor.}\par
3. Réécrivez avec le modal demandé : a) \eng{It is forbidden to take photos in the studio.} (\eng{mustn't}) b) \eng{It is not necessary to book a seat; the debate is free.} (\eng{needn't}) c) \eng{I am sure she is a journalist; I saw her on television.} (\eng{must})
\end{exercice}

\begin{corrige}[Exercice : À vous de jouer]
1. a) \eng{You don't have to buy the paper; it's free online.} b) \eng{Journalists should protect their sources.} c) \eng{He must be tired; he worked all night.} d) \eng{The editor should have verified the information.}\par
2. a) \eng{He must go now.} (le modal est invariable) b) \eng{Can you speak English?} (inversion directe, sans \eng{do}) c) \eng{I will be able to help you tomorrow.} (\eng{can} n'a pas de futur) d) \eng{She must see a doctor.} (base verbale sans \eng{to} après \eng{must}).\par
3. a) \eng{You mustn't take photos in the studio.} b) \eng{You needn't book a seat; the debate is free.} c) \eng{She must be a journalist; I saw her on television.}
\end{corrige}

\begin{experience}[Débat guidé : regulating social media]
Par groupes de quatre, débattez de la question : \eng{Should the government control social media in Cameroon?} Chaque élève doit produire au moins trois phrases avec des modaux différents : une obligation (\eng{must / have to}), un conseil (\eng{should / ought to / had better}) et une possibilité (\eng{may / might / can't}). Exemple : \eng{The government should fight fake news, but it mustn't censor the press.} Un rapporteur résume les arguments du groupe devant la classe.
\end{experience}

\begin{aretenir}[L'essentiel de la section]
Modal + base verbale, jamais de \eng{-s}, jamais de \eng{do}. \eng{Mustn't} = interdit, \eng{don't have to} / \eng{needn't} = pas obligé. En déduction, le contraire de \eng{must} est \eng{can't}. \eng{Can} se remplace par \eng{be able to} au futur et aux temps composés. \eng{Should have} + participe passé = « aurait dû ». Gardez en tête l'échelle de la certitude : \eng{can't} (0 \%) $\rightarrow$ \eng{may/might} (50 \%) $\rightarrow$ \eng{should} (80 \%) $\rightarrow$ \eng{must} (100 \%).
\end{aretenir}

\coursec{Voix passive et discours rapporté}

Dans la section précédente, nous avons révisé les modaux et semi-auxiliaires (\eng{must, should, could, have to...}), qui permettent d'exprimer l'obligation, la possibilité ou le conseil. Ces mêmes modaux réapparaîtront d'ailleurs dans cette section : au passif (\eng{The law must be respected.}) et au discours rapporté (\eng{He said he could help.}). Nous abordons maintenant deux transformations fondamentales de la phrase anglaise, systématiquement présentes à l'examen : la \cle{voix passive} et le \cle{discours rapporté}. Toutes deux obéissent à des mécanismes précis qu'il suffit de maîtriser pour gagner des points faciles.

\courssub{La voix passive : observer avant de comprendre}

\begin{textebox}[From a Cameroonian news report]
Two journalists were arrested yesterday in Douala after a report on corruption was broadcast on a private television channel. The station was visited by police officers and several computers were seized. According to witnesses, no explanation was given to the staff. The programme had been watched by thousands of viewers the previous evening. A press release will be published tomorrow by the journalists' union, and a demonstration is being organised in Yaoundé. In many countries, such events are condemned by international organisations, because press freedom must be protected everywhere.
\end{textebox}

\textbf{Glossaire.} \emph{to seize} (v.) : saisir ; \emph{staff} (n.) : le personnel ; \emph{viewers} (n.) : téléspectateurs ; \emph{press release} (n.) : communiqué de presse ; \emph{to condemn} (v.) : condamner.

Repérez mentalement tous les groupes verbaux du texte : \eng{were arrested}, \eng{was broadcast}, \eng{was visited}, \eng{were seized}, \eng{was given}, \eng{had been watched}, \eng{will be published}, \eng{is being organised}, \eng{are condemned}, \eng{must be protected}. Tous suivent le même moule : un auxiliaire \cle{be} conjugué + un \cle{participe passé}. Pourquoi le journaliste a-t-il choisi cette forme ? Parce que l'important, c'est l'événement (les arrestations, les saisies), pas celui qui l'a provoqué.

\begin{propriete}[Actif → passif : la transformation]
Pour passer de la voix active à la voix passive :
\begin{enumerate}
\item Le \textbf{complément d'objet direct} de la phrase active devient le \textbf{sujet} de la phrase passive.
\item Le verbe se met à la forme \textbf{be + participe passé}, et \textbf{be prend le temps} de la phrase active.
\item L'ancien sujet devient le \textbf{complément d'agent}, introduit par \eng{by} — ou disparaît.
\end{enumerate}
\eng{They broadcast the news.} → \eng{The news is broadcast (by them).} « On diffuse les informations. / Les informations sont diffusées. »
\end{propriete}

\begin{attention}[Le temps ne change jamais !]
C'est l'auxiliaire \eng{be} qui porte le temps, pas le participe. Erreur classique : employer \eng{The news was broadcast} pour un présent. Non : actif au présent (\eng{they broadcast}) → passif au présent (\eng{is broadcast}). Actif au prétérit (\eng{they broadcast}) → passif au prétérit (\eng{was broadcast}).
\end{attention}

\begin{center}
\begin{tabularx}{\linewidth}{|l|X|X|}
\hline
\rowcolor{popblueL} Temps & Actif & Passif \\
\hline
Present simple & They write articles. & Articles \textbf{are written}. \\
\hline
Present continuous & They are filming the match. & The match \textbf{is being filmed}. \\
\hline
Past simple & The police arrested him. & He \textbf{was arrested}. \\
\hline
Present perfect & They have sold the house. & The house \textbf{has been sold}. \\
\hline
Past perfect & They had cancelled the show. & The show \textbf{had been cancelled}. \\
\hline
Future & They will publish the results. & The results \textbf{will be published}. \\
\hline
Modal & We must protect children. & Children \textbf{must be protected}. \\
\hline
\end{tabularx}
\end{center}

\courssub{Emplois du passif et complément d'agent}

\begin{definition}[Quand utilise-t-on le passif ?]
\begin{itemize}
\item Quand on \textbf{ignore} qui fait l'action : \eng{My phone has been stolen.} « Mon téléphone a été volé. »
\item Quand l'agent est \textbf{évident} ou sans importance : \eng{The thief was arrested.} (évidemment par la police).
\item Quand on veut \textbf{minimiser} l'agent ou rester impersonnel — style journalistique et scientifique : \eng{Two journalists were arrested yesterday.} ; \eng{The samples were heated to 100 degrees.}
\end{itemize}
Le complément d'agent en \eng{by} n'est gardé que s'il apporte une information réelle : \eng{The song was written by a student from Buea.} « La chanson a été écrite par un élève de Buea. »
\end{definition}

\begin{attention}[Faux amis et pièges du français]
\begin{itemize}
\item Le français « on » se rend souvent par le passif anglais : « On parle anglais ici. » → \eng{English is spoken here.} (et non \eng{One speaks...}).
\item Jamais de passif avec les verbes intransitifs (sans COD) : \eng{The accident was happened.} ✗ → \eng{The accident happened.} ✓
\item Attention aux participes passés irréguliers : \emph{broadcast--broadcast--broadcast}, \emph{write--wrote--written}, \emph{steal--stole--stolen}.
\end{itemize}
\end{attention}

\begin{exoresolu}[Mettez les phrases au passif]
Énoncé. 1) \eng{The students cleaned the classroom.} 2) \eng{They will build a new bridge in Douala.} 3) \eng{Someone has broken the window.}
\tcblower
Solution. 1) Le COD \eng{the classroom} devient sujet ; \eng{cleaned} est au prétérit → \eng{was} : \eng{The classroom was cleaned by the students.} 2) Futur → \eng{will be} : \eng{A new bridge will be built in Douala.} 3) Present perfect → \eng{has been} ; l'agent \eng{someone} est omis : \eng{The window has been broken.}
\end{exoresolu}

\courssub{Le discours rapporté : le recul des temps}

Passons à une situation de la vie scolaire. Le matin, Aïcha dit : \eng{“I am tired.”} L'après-midi, tu racontes : \eng{Aïcha said she was tired.} En anglais comme en français, rapporter des paroles provoque un \cle{recul des temps} : le présent recule au passé.

\begin{exemplebox}[Observer la transformation]
\begin{itemize}
\item \eng{“I am tired,” she said.} → \eng{She said (that) she was tired.} « Je suis fatiguée », a-t-elle dit. → Elle a dit qu'elle était fatiguée.
\item \eng{“Where do you live?” he asked.} → \eng{He asked me where I lived.} « Où habites-tu ? » → Il m'a demandé où j'habitais.
\item \eng{“I will call you tomorrow,” he said.} → \eng{He said he would call me the next day.} « Je t'appellerai demain. » → Il a dit qu'il m'appellerait le lendemain.
\end{itemize}
\end{exemplebox}

\begin{propriete}[Les règles du recul des temps]
Quand le verbe introducteur est au passé (\eng{said, told, asked}), les temps rapportés reculent :
\begin{itemize}
\item present simple → past simple : \eng{“I work hard.”} → \eng{He said he worked hard.}
\item present continuous → past continuous : \eng{“I am watching TV.”} → \eng{She said she was watching TV.}
\item past simple / present perfect → past perfect : \eng{“I saw / have seen the film.”} → \eng{He said he had seen the film.}
\item \eng{will} → \eng{would} ; \eng{can} → \eng{could} ; \eng{must} → \eng{had to}.
\end{itemize}
\textbf{Exception :} si le verbe introducteur est au présent (\eng{He says...}) ou si la vérité rapportée est permanente, pas de recul : \eng{The teacher said that water boils at 100 degrees.}
\end{propriete}

\begin{popfigure}
\begin{center}
\begin{tikzpicture}[
box/.style={draw=popblue, fill=popblueL, rounded corners=2pt, align=center, font=\small, inner sep=4pt},
box2/.style={draw=poporange, fill=poporangeL, rounded corners=2pt, align=center, font=\small, inner sep=4pt},
fleche/.style={-{Stealth[length=3mm]}, thick, popdark}]
\node[box, align=center] (a) at (0,2.4){present simple\\ \eng{I work}};
\node[box2, align=center] (b) at (5.6,2.4){past simple\\ \eng{he worked}};
\draw[fleche] (a) -- (b);
\node[box, align=center] (c) at (0,1.2){\eng{will}\\ \eng{I will come}};
\node[box2, align=center] (d) at (5.6,1.2){\eng{would}\\ \eng{he would come}};
\draw[fleche] (c) -- (d);
\node[box, align=center] (e) at (0,0){\eng{can}\\ \eng{I can swim}};
\node[box2, align=center] (f) at (5.6,0){\eng{could}\\ \eng{she could swim}};
\draw[fleche] (e) -- (f);
\node[box, align=center] (g) at (0,-1.2){past simple / present perfect\\ \eng{I saw / have seen}};
\node[box2, align=center] (h) at (5.6,-1.2){past perfect\\ \eng{he had seen}};
\draw[fleche] (g) -- (h);
\node[font=\small\itshape, popmuted] at (0,3.1) {discours direct};
\node[font=\small\itshape, popmuted] at (5.6,3.1) {discours rapporté};
\end{tikzpicture}
\end{center}
\legende{Le recul des temps au discours rapporté : chaque temps « recule » d'un cran vers le passé.}
\end{popfigure}

\begin{propriete}[Les marqueurs de temps et de lieu changent aussi]
\begin{center}
\begin{tabularx}{\linewidth}{|X|X|}
\hline
\rowcolor{popblueL} Discours direct & Discours rapporté \\
\hline
\eng{now} & \eng{then} \\
\hline
\eng{today} & \eng{that day} \\
\hline
\eng{yesterday} & \eng{the day before} \\
\hline
\eng{tomorrow} & \eng{the next day / the following day} \\
\hline
\eng{ago} (\eng{two days ago}) & \eng{before} (\eng{two days before}) \\
\hline
\eng{here} & \eng{there} \\
\hline
\eng{this / these} & \eng{that / those} \\
\hline
\end{tabularx}
\end{center}
\eng{“I saw this film here yesterday.”} → \eng{She said she had seen that film there the day before.}
\end{propriete}

\courssub{Questions, ordres et conseils rapportés}

\begin{methode}[Rapporter une question en 4 étapes]
\begin{enumerate}
\item Choisir le verbe introducteur : \eng{asked} (+ éventuellement un complément : \eng{asked me}).
\item Pour une question \textbf{fermée} (réponse oui/non) : utiliser \eng{if} ou \eng{whether}. Pour une question \textbf{ouverte} : garder le mot interrogatif (\eng{where, when, why, how...}).
\item Remettre la phrase à l'\textbf{ordre affirmatif} : sujet + verbe, \textbf{sans inversion} et sans \eng{do/does/did}.
\item Faire reculer le temps.
\end{enumerate}
\eng{“Do you like the programme?”} → \eng{He asked if I liked the programme.} « Il a demandé si j'aimais l'émission. »
\end{methode}

\begin{attention}[L'erreur n° 1 : l'inversion]
Au discours rapporté, il n'y a \textbf{jamais d'inversion} : \eng{He asked me where did I live.} ✗ → \eng{He asked me where I lived.} ✓ Le point d'interrogation disparaît aussi : la phrase rapportée se termine par un point.
\end{attention}

\begin{propriete}[Ordres et conseils rapportés]
Structure : \textbf{tell / order / ask / advise + complément + to + base verbale}. À la forme négative : \eng{not to + base verbale}.
\begin{itemize}
\item \eng{“Switch off your phone,” she said.} → \eng{She told me to switch off my phone.} « Elle m'a dit d'éteindre mon téléphone. »
\item \eng{“Don't be late,” the teacher said.} → \eng{The teacher told us not to be late.} « Le professeur nous a dit de ne pas être en retard. »
\item \eng{“You should revise regularly.”} → \eng{He advised me to revise regularly.} « Il m'a conseillé de réviser régulièrement. »
\end{itemize}
\end{propriete}

\begin{attention}[Say ou tell ? Le grand classique]
\eng{say something (to somebody)} mais \eng{tell somebody something}. On ne dit jamais \eng{He said me...} ✗ → \eng{He told me...} ✓ ou \eng{He said to me...} (plus rare). De même : \eng{tell the truth, tell a lie, tell a story}, mais \eng{say hello, say goodbye, say sorry}.
\end{attention}

\begin{definition}[Varier les verbes introducteurs]
Pour enrichir vos productions écrites, remplacez \eng{said} par un verbe plus précis : \eng{explain} (expliquer), \eng{add} (ajouter), \eng{claim} (affirmer, prétendre), \eng{admit} (admettre), \eng{deny} (nier), \eng{promise} (promettre), \eng{warn} (avertir), \eng{suggest} (suggérer), \eng{reply} (répondre).
\eng{He denied having taken the money.} « Il a nié avoir pris l'argent. » ; \eng{She promised to help me.} « Elle a promis de m'aider. » ; \eng{The captain warned us that the road was dangerous.} « Le capitaine nous a avertis que la route était dangereuse. »
\end{definition}

\begin{exoresolu}[Rapportez les paroles au discours indirect]
Énoncé. 1) \eng{“I have finished my homework,” Mary said.} 2) \eng{“Did you watch the match?” Peter asked me.} 3) \eng{“Don't touch that wire!” the electrician said to us.}
\tcblower
Solution. 1) Present perfect → past perfect : \eng{Mary said (that) she had finished her homework.} 2) Question fermée → \eng{if} + ordre affirmatif + recul : \eng{Peter asked me if I had watched the match.} 3) Ordre négatif → \eng{told + complément + not to} : \eng{The electrician told us not to touch that wire.}
\end{exoresolu}

\begin{exercice}[À vous de jouer]
1) Mettez au passif : \eng{The company employs 200 workers.} / \eng{They have cancelled the concert.}
2) Rapportez : \eng{“I will meet you here tomorrow,” John said.} / \eng{“Why are you late?” the teacher asked the boy.}
3) Corrigez l'erreur : \eng{She said me that she was busy.}
\end{exercice}

\begin{corrige}[Exercice « À vous de jouer »]
1) \eng{200 workers are employed by the company.} ; \eng{The concert has been cancelled.} 2) \eng{John said he would meet me there the next day.} ; \eng{The teacher asked the boy why he was late.} 3) \eng{She told me that she was busy.} (ou \eng{She said that she was busy.})
\end{corrige}

\begin{aretenir}[L'essentiel de la section]
\begin{itemize}
\item Passif : COD → sujet ; \eng{be} au temps de la phrase active + participe passé ; agent en \eng{by} souvent omis.
\item Discours rapporté : recul des temps (present → past, will → would, can → could, past/present perfect → past perfect) et changement des marqueurs (\eng{today} → \eng{that day}...).
\item Questions rapportées : \eng{if/whether} ou mot interrogatif + ordre affirmatif, sans inversion.
\item Ordres : \eng{tell/ask/advise + complément + (not) to + base verbale}.
\item \eng{say something} / \eng{tell somebody} : ne jamais confondre.
\end{itemize}
\end{aretenir}

\coursec{Conditionnels, souhaits et phrases complexes}

Dans la section précédente, nous avons révisé la voix passive et le discours rapporté, deux mécanismes qui permettent de \emph{transformer} une phrase sans changer son sens. Nous abordons maintenant un autre pilier de l'examen : les phrases \cle{complexes}, c'est-à-dire celles qui combinent plusieurs propositions grâce à \eng{if}, \eng{wish}, \eng{unless}, aux pronoms relatifs et aux connecteurs logiques. C'est dans ces structures que le correcteur du Baccalauréat repère immédiatement un candidat solide.

\courssub{Les quatre conditionnels : observer avant de conclure}

\begin{textebox}[Reading text — “The Editorial Meeting”]
Every morning, the editorial team of \emph{The Coastal Herald} meets in Douala. “If the Internet connection works today, we will watch the press conference live,” says the editor. Last month, the connection failed during an important debate, and the journalists missed the story. “If we had checked the equipment the day before, we would not have lost that scoop,” admits the technician. He wishes he had been more careful.

The editor then turns to a young trainee. “If I were you, I would verify every fact twice before publishing. If you publish false information, readers lose trust — and they never come back.” She adds that if a newspaper lies once, its reputation suffers for years. It is a general rule of the profession: if you heat ice, it melts; if you cheat your readers, they leave.

The trainee wishes she had more experience, but she is determined to learn. “Unless we work as a team, we cannot produce a good paper,” the editor concludes. “A newspaper is like an orchestra: if one musician plays the wrong note, the whole concert is ruined.”

\textbf{Glossary}: scoop (n.) = exclusivité journalistique ; trainee (n.) = stagiaire ; trust (n.) = confiance ; to cheat (v.) = tromper.
\end{textebox}

\textbf{Questions de compréhension}
\begin{enumerate}
\item What problem did the team face last month?
\item What advice does the editor give to the trainee?
\item Find in the text one example of each conditional type (0, 1, 2, 3).
\end{enumerate}

Ce court texte contient les quatre types de conditionnels. Observons-les avant d'énoncer les règles.

\begin{popfigure}
\begin{tikzpicture}[font=\small]
\draw[fill=popblueL, draw=popblue, thick] (0,0) rectangle (12.5,5.4);
\node[font=\bfseries\small, text=popdark] at (6.25,5.0) {The four conditionals};
\node[font=\bfseries\footnotesize, text=popblue] at (1.1,4.4) {Type};
\node[font=\bfseries\footnotesize, text=popblue] at (4.6,4.4) {Structure};
\node[font=\bfseries\footnotesize, text=popblue] at (8.2,4.4) {Meaning};
\node[font=\bfseries\footnotesize, text=popblue] at (11.0,4.4) {Example};
\draw[popline] (0.3,4.1) -- (12.2,4.1);
\node[font=\footnotesize\bfseries, text=popgreen] at (1.1,3.5) {0};
\node[font=\footnotesize, align=center] at (4.6,3.5) {If + present,\\present};
\node[font=\footnotesize, align=center] at (8.2,3.5) {general truth,\\scientific fact};
\node[font=\footnotesize, align=center] at (11.0,3.5) {If you heat ice,\\it melts.};
\draw[popline] (0.3,2.95) -- (12.2,2.95);
\node[font=\footnotesize\bfseries, text=poporange] at (1.1,2.35) {1};
\node[font=\footnotesize, align=center] at (4.6,2.35) {If + present,\\will + base verb};
\node[font=\footnotesize, align=center] at (8.2,2.35) {real, probable\\future};
\node[font=\footnotesize, align=center] at (11.0,2.35) {If the connection works,\\we will watch the debate.};
\draw[popline] (0.3,1.8) -- (12.2,1.8);
\node[font=\footnotesize\bfseries, text=poppurple] at (1.1,1.2) {2};
\node[font=\footnotesize, align=center] at (4.6,1.2) {If + past simple,\\would + base verb};
\node[font=\footnotesize, align=center] at (8.2,1.2) {imaginary,\\unreal present};
\node[font=\footnotesize, align=center] at (11.0,1.2) {If I were a journalist,\\I would investigate.};
\draw[popline] (0.3,0.65) -- (12.2,0.65);
\node[font=\footnotesize\bfseries, text=poppink] at (1.1,0.15) {3};
\node[font=\footnotesize, align=center] at (4.6,0.15) {If + past perfect,\\would have + past participle};
\node[font=\footnotesize, align=center] at (8.2,0.15) {unreal past,\\regret};
\node[font=\footnotesize, align=center] at (11.0,0.15) {If they had checked the facts,\\the scandal would not have happened.};
\end{tikzpicture}
\legende{Tableau récapitulatif des quatre conditionnels}
\end{popfigure}

\begin{propriete}[Règle d'or : jamais de WILL après IF]
Dans la subordonnée introduite par \eng{if}, on n'emploie \textbf{jamais} \eng{will}, même si le sens est futur. Le futur (\eng{will}) se place dans la proposition principale.\par
\eng{If it will rain, we will stay home.} ✗\par
\eng{If it rains, we will stay home.} ✓ « S'il pleut, nous resterons à la maison. »\par
Exception rare : \eng{will} après \eng{if} exprime alors la \emph{volonté} : \eng{If you will insist on smoking, do it outside.} « Si tu tiens absolument à fumer… »
\end{propriete}

\begin{popfigure}
\begin{tikzpicture}[font=\small, node distance=0.9cm]
\node[draw=popblue, thick, fill=popblueL, rounded corners, align=center, font=\bfseries] (if) {IF};
\node[draw=popgreen, thick, fill=popgreenL, rounded corners, align=center, right=2.6cm of if, yshift=1.8cm] (t1) {probable, real future\\\eng{If it rains, we will stay.}};
\node[draw=poppurple, thick, fill=poppurpleL, rounded corners, align=center, right=2.6cm of if] (t2) {imaginary, unreal now\\\eng{If I were rich, I would travel.}};
\node[draw=poppink, thick, fill=poppinkL, rounded corners, align=center, right=2.6cm of if, yshift=-1.8cm] (t3) {unreal past, regret\\\eng{If I had known, I would have come.}};
\draw[-{Stealth}, thick, popgreen] (if) -- (t1);
\draw[-{Stealth}, thick, poppurple] (if) -- (t2);
\draw[-{Stealth}, thick, poppink] (if) -- (t3);
\end{tikzpicture}
\legende{Choisir son conditionnel : probable, imaginaire ou passé ?}
\end{popfigure}

\begin{exemplebox}[Exemples traduits]
\eng{If I were you, I would apologise.} « Si j'étais toi, je m'excuserais. »\par
\eng{If the government invests in rural radio, more farmers will get information.} « Si le gouvernement investit dans la radio rurale, davantage d'agriculteurs seront informés. »\par
\eng{If she had read the whole article, she would have understood the journalist's point.} « Si elle avait lu l'article en entier, elle aurait compris le point de vue du journaliste. »\par
\eng{If you drop your phone, the screen breaks.} « Si tu laisses tomber ton téléphone, l'écran se casse. » (vérité générale, type 0)
\end{exemplebox}

\begin{attention}[Were pour toutes les personnes]
Au conditionnel de type 2, la forme soutenue — et \textbf{recommandée à l'examen} — emploie \eng{were} pour toutes les personnes : \eng{If I were}, \eng{If she were}, \eng{If it were}. La forme \eng{If I was} existe à l'oral familier, mais elle est pénalisée dans une copie. Retenez aussi l'expression figée \eng{If I were you…} « Si j'étais toi/vous… », idéale pour donner un conseil.
\end{attention}

\courssub{Wish : exprimer le regret}

\begin{propriete}[Les constructions de WISH]
\begin{itemize}
\item \textbf{Regret sur le présent} : \eng{wish} + \cle{past simple}. \eng{I wish I had more time to read.} « J'aimerais avoir plus de temps pour lire. » (mais je n'en ai pas)
\item \textbf{Regret sur le passé} : \eng{wish} + \cle{past perfect}. \eng{I wish I had revised earlier.} « Si seulement j'avais révisé plus tôt. » (mais je ne l'ai pas fait)
\item \textbf{Souhait de changement} : \eng{wish} + \eng{would} + base verbale (pour se plaindre d'une situation agaçante). \eng{I wish the neighbours would stop that noise.} « J'aimerais que les voisins arrêtent ce bruit. »
\end{itemize}
\end{propriete}

\begin{exemplebox}[Exemple traduit]
\eng{I wish I had listened to the teacher.} « Si seulement j'avais écouté le professeur. »\par
\eng{She wishes she lived near the school.} « Elle aimerait habiter près de l'école. »
\end{exemplebox}

\begin{attention}[La logique du recul temporel]
Avec \eng{wish} comme avec les conditionnels, l'anglais fait \emph{reculer} le temps d'un cran pour marquer l'irréel : présent $\rightarrow$ past simple ; passé $\rightarrow$ past perfect. Ne traduisez jamais mot à mot : \eng{I wish I have} ✗. Dites \eng{I wish I had} ✓.
\end{attention}

\courssub{Unless : la condition négative}

\begin{definition}[Unless]
\eng{Unless} signifie \eng{if\ldots not} « à moins que… (ne) ». Attention : la phrase qui suit \eng{unless} est \textbf{affirmative} en anglais, alors que le français utilise souvent la négation.\par
\eng{Unless you subscribe, you cannot read the full article.} « À moins que vous ne vous abonniez, vous ne pouvez pas lire l'article en entier. » (= \eng{If you do not subscribe, you cannot read the full article.})
\end{definition}

\begin{attention}[Ne pas doubler la négation]
\eng{Unless you don't pay, you can't watch the match.} ✗ — Dites : \eng{Unless you pay, you can't watch the match.} ✓ « À moins de payer, tu ne peux pas regarder le match. »
\end{attention}

\courssub{Les pronoms relatifs : relier deux idées}

\begin{propriete}[Qui remplace quoi ?]
\begin{itemize}
\item \cle{who} : les personnes (sujet). \eng{The journalist who interviewed the minister is famous.}
\item \cle{which} : les choses et les animaux. \eng{The article which was published yesterday caused a debate.}
\item \cle{that} : personnes \emph{ou} choses, dans les propositions déterminatives uniquement. \eng{The radio that my grandfather bought still works.}
\item \cle{whose} : la possession (« dont », « duquel »). \eng{The reporter whose car was stolen called the police.}
\item \cle{where} : le lieu (« où »). \eng{This is the studio where the news is recorded.}
\end{itemize}
\end{propriete}

\begin{center}
\begin{tabularx}{\linewidth}{|X|X|X|}
\hline
\rowcolor{popblueL} Français & English & Remarque \\
\hline
L'homme qui parle & The man who is speaking & personne, sujet \\
\hline
Le livre que je lis & The book (that) I am reading & relatif objet : omissible \\
\hline
La ville où je suis né & The town where I was born & lieu \\
\hline
Le garçon dont le père est médecin & The boy whose father is a doctor & possession \\
\hline
\end{tabularx}
\end{center}

\begin{propriete}[Déterminative ou explicative ? La question des virgules]
\begin{itemize}
\item \textbf{Proposition déterminative} (indispensable au sens) : \emph{sans virgules}, \eng{that} possible, relatif objet omissible. \eng{The students who revised passed the exam.} « Les élèves qui ont révisé ont réussi. » (seulement ceux-là)
\item \textbf{Proposition explicative} (simple précision) : \emph{entre virgules}, jamais \eng{that}, jamais d'omission. \eng{My brother, who lives in Buea, works for a TV channel.} « Mon frère, qui habite à Buea, travaille pour une chaîne de télévision. »
\end{itemize}
\end{propriete}

\begin{savaistu}[Whose pour les choses ?]
Dans un style soutenu — très apprécié à l'examen — \eng{whose} peut s'employer pour des choses, et pas seulement pour des personnes : \eng{a country whose media are free} « un pays dont les médias sont libres » ; \eng{a newspaper whose readers trust it} « un journal dont les lecteurs lui font confiance ». Cela évite la lourdeur de \eng{of which}.
\end{savaistu}

\courssub{Les connecteurs logiques : articuler sa pensée}

\begin{center}
\begin{tabularx}{\linewidth}{|l|X|X|}
\hline
\rowcolor{popblueL} Fonction & Connecteurs & Exemple \\
\hline
Addition & moreover, furthermore, also & \eng{The report is biased; moreover, it is outdated.} \\
\hline
Opposition & however, whereas, although & \eng{Although the signal was weak, we heard the speech.} \\
\hline
Cause & because, since, as & \eng{Since the facts were false, the editor apologised.} \\
\hline
Conséquence & therefore, consequently, so & \eng{The rumour spread fast; therefore, panic followed.} \\
\hline
Ordre & first, then, finally & \eng{First, check the source; then, publish.} \\
\hline
\end{tabularx}
\end{center}

\begin{attention}[Although et but ne font pas bon ménage]
Le français dit « bien que… mais… » ; l'anglais interdit la double marque : \eng{Although it was late, but he continued.} ✗. Choisissez : \eng{Although it was late, he continued.} ✓ ou \eng{It was late, but he continued.} ✓. Même règle pour \eng{because} et \eng{so} : jamais ensemble dans la même phrase.
\end{attention}

\begin{exoresolu}[Exercice type Bac : phrases complexes]
Énoncé — Complete each sentence with the correct form of the verb in brackets, then link sentences 5 and 6 with a relative pronoun.\par
1. If the journalist (check) his sources, the article would not have been withdrawn.\par
2. Unless the government (act) quickly, fake news will spread.\par
3. I wish I (speak) English fluently.\par
4. If you mix blue and yellow, you (get) green.\par
5. The reporter won a prize. She investigated corruption in the port of Douala.
\tcblower
Solution détaillée\par
1. \eng{If the journalist \textbf{had checked} his sources, the article would not have been withdrawn.} — Type 3 (regret sur le passé) : \eng{if} + past perfect, \eng{would have} + participe passé.\par
2. \eng{Unless the government \textbf{acts} quickly, fake news will spread.} — \eng{unless} = \eng{if\ldots not} ; après \eng{unless}, présent simple (jamais \eng{will}), et phrase affirmative.\par
3. \eng{I wish I \textbf{spoke} English fluently.} — Regret sur le présent : \eng{wish} + past simple (« J'aimerais parler anglais couramment »).\par
4. \eng{If you mix blue and yellow, you \textbf{get} green.} — Type 0 : vérité générale, présent dans les deux propositions.\par
5. \eng{The reporter \textbf{who} investigated corruption in the port of Douala won a prize.} — Proposition déterminative (sans virgules) ; \eng{who} car il s'agit d'une personne sujet. « La journaliste qui a enquêté sur la corruption dans le port de Douala a remporté un prix. »
\end{exoresolu}

\begin{exercice}[À vous de jouer]
1. Put the verbs in the correct tense: (a) If I (be) you, I (apologise) to the editor. (b) If they had listened to the weather forecast, they (not/travel). (c) I wish I (revise) the conditionals before the test.\par
2. Join the sentences with \eng{who}, \eng{which}, \eng{whose} or \eng{where}: (a) The studio is in Yaoundé. The programme is recorded there. (b) The student has won a scholarship. Her essay was excellent.\par
3. Rewrite with \eng{unless}: If you do not switch off your phone, the teacher will take it.
\end{exercice}

\begin{corrige}[Exercice « À vous de jouer »]
1. (a) \eng{If I \textbf{were} you, I \textbf{would apologise} to the editor.} (type 2, \eng{were} pour toutes les personnes). (b) \eng{If they had listened to the weather forecast, they \textbf{would not have travelled}.} (type 3). (c) \eng{I wish I \textbf{had revised} the conditionals before the test.} (regret passé : past perfect).\par
2. (a) \eng{The studio \textbf{where} the programme is recorded is in Yaoundé.} (b) \eng{The student \textbf{whose} essay was excellent has won a scholarship.}\par
3. \eng{\textbf{Unless} you switch off your phone, the teacher \textbf{will take} it.} (phrase affirmative après \eng{unless}, \eng{will} dans la principale).
\end{corrige}

\begin{aretenir}[L'essentiel de la section]
\begin{itemize}
\item Quatre conditionnels : type 0 (présent + présent), type 1 (présent + \eng{will}), type 2 (past simple + \eng{would}), type 3 (past perfect + \eng{would have} + participe).
\item Jamais de \eng{will} après \eng{if} ; \eng{were} pour toutes les personnes au type 2.
\item \eng{Wish} + past simple (regret présent), \eng{wish} + past perfect (regret passé).
\item \eng{Unless} = \eng{if\ldots not}, suivi d'une forme affirmative.
\item Relatifs : \eng{who} (personnes), \eng{which} (choses), \eng{that} (déterminatives), \eng{whose} (possession), \eng{where} (lieu) ; virgules = explicative, jamais \eng{that}.
\item Connecteurs : un seul marqueur par relation — \eng{although} \emph{ou} \eng{but}, \eng{because} \emph{ou} \eng{so}.
\end{itemize}
\end{aretenir}

\coursec{Pièges des francophones : le grand nettoyage}

Après avoir révisé les conditionnels, les souhaits et les phrases complexes, il nous reste un dernier chantier, peut-être le plus rentable pour l'examen : le \cle{grand nettoyage} des erreurs typiques des francophones. Ce sont des fautes que l'on retrouve chaque année dans les copies du baccalauréat : faux amis, pluriels impossibles, prépositions oubliées, doubles négations. Les corriger, c'est gagner des points immédiatement, car le correcteur les repère en une seconde.

\courssub{Les faux amis récurrents}

Observez d'abord cette phrase extraite d'une copie d'élève :

\eng{I am actually in Yaoundé, so I will assist the meeting in the library.}

Cette phrase contient \textbf{trois faux amis}. Un \cle{faux ami} (\eng{false friend}) est un mot anglais qui ressemble à un mot français mais n'a pas le même sens. Le cerveau francophone traduit automatiquement, et c'est le piège. La phrase ci-dessus voulait dire : « Je suis actuellement à Yaoundé, donc j'assisterai à la réunion à la bibliothèque. »

\begin{definition}[Les cinq faux amis à connaître par cœur]
\begin{itemize}
\item \eng{actually} = \textbf{en fait, en réalité} (et non « actuellement »). Pour « actuellement », dites \eng{currently} ou \eng{at present}.
\item \eng{to assist} = \textbf{aider} (et non « assister à »). Pour « assister à », dites \eng{to attend}.
\item \eng{a library} = \textbf{une bibliothèque} (et non « une librairie »). Une librairie se dit \eng{a bookshop} (britannique) ou \eng{a bookstore} (américain).
\item \eng{sensible} = \textbf{raisonnable, sensé} (et non « sensible »). « Sensible » se dit \eng{sensitive}.
\item \eng{eventually} = \textbf{finalement, à la fin} (et non « éventuellement »). « Éventuellement » se dit \eng{possibly} ou \eng{if necessary}.
\end{itemize}
\end{definition}

\begin{exemplebox}[Exemples traduits]
\eng{Actually, I don't like football.} = En fait, je n'aime pas le football.\par
\eng{She is currently working in Douala.} = Elle travaille actuellement à Douala.\par
\eng{The nurse assisted the doctor during the operation.} = L'infirmière a aidé le médecin pendant l'opération.\par
\eng{Thousands of fans attended the match in Buea.} = Des milliers de supporters ont assisté au match à Buea.\par
\eng{I borrowed this novel from the library.} = J'ai emprunté ce roman à la bibliothèque.\par
\eng{Be sensible: don't spend all your money at the market.} = Sois raisonnable : ne dépense pas tout ton argent au marché.\par
\eng{He failed many times, but he eventually passed the exam.} = Il a échoué plusieurs fois, mais il a finalement réussi l'examen.
\end{exemplebox}

\begin{attention}[Faux amis : le piège de l'examen]
Dans les sujets de compréhension, un faux ami peut vous faire comprendre le texte à l'envers. Si vous lisez \eng{The event was eventually cancelled}, cela signifie « L'événement a \textbf{finalement} été annulé », et non « éventuellement annulé ». Soulignez mentalement chaque mot qui ressemble au français et vérifiez son sens dans le contexte. Autres pièges fréquents : \eng{a lecture} = une conférence (et non « une lecture ») ; \eng{to demand} = exiger (et non « demander », qui se dit \eng{to ask}) ; \eng{a journey} = un voyage (et non « une journée », qui se dit \eng{a day}).
\end{attention}

\courssub{Les indénombrables : des mots toujours au singulier}

En français, on dit « des informations », « des conseils », « les nouvelles sont bonnes ». En anglais, certains noms sont \cle{indénombrables} (\eng{uncountable nouns}) : ils n'ont \textbf{jamais} de pluriel et prennent toujours un verbe au singulier.

\begin{propriete}[Les cinq indénombrables piégeux]
\eng{information}, \eng{news}, \eng{advice}, \eng{furniture}, \eng{homework} sont \textbf{toujours au singulier} : pas de \eng{-s}, pas de \eng{a/an} devant, verbe au singulier.\par
\eng{The news is good.} = Les nouvelles sont bonnes. (Attention : \eng{news} finit par un \eng{-s} mais reste singulier.)\par
Pour exprimer la quantité, on utilise \eng{some}, \eng{a piece of}, \eng{a bit of} : \eng{a piece of advice} = un conseil ; \eng{a piece of furniture} = un meuble ; \eng{some homework} = des devoirs.
\end{propriete}

\begin{attention}['informations' ✗ → 'information' ✓]
Ne dites jamais \eng{informations}, \eng{advices} ou \eng{furnitures} : ces formes n'existent pas.\par
\eng{I need some information about the exam.} = J'ai besoin de renseignements sur l'examen.\par
\eng{She gave me two pieces of advice.} = Elle m'a donné deux conseils.
\end{attention}

\begin{center}
\begin{tabularx}{\linewidth}{|X|X|X|}
\hline
\rowcolor{popblueL} Français & English & Remarque \\
\hline
des informations & (some) information & jamais de -s \\
\hline
les nouvelles sont bonnes & the news is good & verbe au singulier \\
\hline
un conseil & a piece of advice & « advice » sans -s \\
\hline
des meubles & furniture & un meuble = a piece of furniture \\
\hline
beaucoup de devoirs & a lot of homework & jamais « homeworks » \\
\hline
\end{tabularx}
\end{center}

\courssub{Les prépositions piégeuses}

Le français et l'anglais n'utilisent pas les mêmes prépositions, et parfois le français n'en met aucune là où l'anglais en exige une — ou l'inverse. C'est l'une des sources d'erreurs les plus fréquentes.

\begin{propriete}[Verbes et adjectifs à préposition obligatoire]
\begin{itemize}
\item \eng{to listen \textbf{to}} : \eng{Listen to the teacher.} = Écoute le professeur. (le français « écouter » est direct, l'anglais exige \eng{to})
\item \eng{to wait \textbf{for}} : \eng{I am waiting for the bendskin.} = J'attends le bendskin.
\item \eng{to depend \textbf{on}} : \eng{It depends on the weather.} = Cela dépend du temps.
\item \eng{to be interested \textbf{in}} : \eng{She is interested in journalism.} = Elle s'intéresse au journalisme.
\item \eng{to be good \textbf{at}} : \eng{He is good at mathematics.} = Il est bon en mathématiques.
\item \eng{to consist \textbf{of}} : \eng{The team consists of eleven players.} = L'équipe se compose de onze joueurs.
\end{itemize}
\end{propriete}

\begin{exemplebox}[Deux erreurs classiques de construction]
\eng{I am agree} ✗ → \eng{I agree} ✓ = Je suis d'accord. (\eng{agree} est un verbe, pas un adjectif : pas de \eng{be} devant.)\par
\eng{He married with her} ✗ → \eng{He married her} ✓ = Il l'a épousée. (\eng{to marry} est un verbe transitif direct : pas de préposition. En revanche, avec l'adjectif : \eng{She is married to a teacher.} = Elle est mariée à un professeur.)
\end{exemplebox}

\courssub{L'accord sujet-verbe : les cas qui trompent}

\begin{propriete}[Singulier ou pluriel ?]
\begin{itemize}
\item \eng{Everybody is...}, \eng{everyone is...}, \eng{somebody is...} : toujours le \textbf{singulier}. \eng{Everybody is ready.} = Tout le monde est prêt.
\item \eng{The police are...} : toujours le \textbf{pluriel} en anglais. \eng{The police are investigating.} = La police enquête. (Pour un seul agent : \eng{a police officer}.)
\item \eng{Neither of them is...} : \textbf{singulier}. \eng{Neither of them is coming.} = Aucun des deux ne vient.
\item \eng{There is} + singulier, \eng{there are} + pluriel : \eng{There is a problem.} / \eng{There are two problems.} Ne dites jamais « there is two problems ».
\end{itemize}
\end{propriete}

\courssub{Les adjectifs : invariables et ordonnés}

\begin{propriete}[Deux règles d'or des adjectifs anglais]
\textbf{1. L'adjectif anglais est invariable} : jamais de \eng{-s}, jamais d'accord. \eng{tall girls} = des filles grandes (et non « talls girls »). \eng{a black bag, two black bags}.\par
\textbf{2. L'ordre des adjectifs} devant le nom suit une hiérarchie fixe : \textbf{opinion + taille + âge + couleur} (+ origine + matière).\par
\eng{a beautiful small old black bag} = un beau petit vieux sac noir.\par
En français, l'ordre est différent et l'adjectif se place souvent après le nom : ne transposez jamais l'ordre français.
\end{propriete}

\courssub{La double négation : interdite en anglais}

\begin{propriete}[Une seule négation par phrase]
En français on dit « je ne sais rien » avec deux éléments négatifs. En anglais standard, \textbf{une seule négation} par phrase :\par
\eng{I don't know anything.} ✓ (et non « I don't know nothing » ✗) = Je ne sais rien.\par
\eng{Nobody came.} ✓ (et non « Nobody didn't come » ✗) = Personne n'est venu.\par
Les mots \eng{anything, anyone, anybody, anywhere, ever} accompagnent la négation ; \eng{nothing, nobody, no one, nowhere, never} la remplacent.
\end{propriete}

\courssub{Much / many, few / little : dénombrable ou non ?}

\begin{propriete}[Le choix dépend de la nature du nom]
\begin{itemize}
\item \eng{much} + \textbf{indénombrable} : \eng{much money, much time, much information}.
\item \eng{many} + \textbf{dénombrable pluriel} : \eng{many books, many students, many ideas}.
\item \eng{little} (peu de) + indénombrable ; \eng{few} (peu de) + dénombrable : \eng{little water} = peu d'eau ; \eng{few friends} = peu d'amis.
\item \eng{a little} / \eng{a few} = un peu de / quelques (sens positif) : \eng{I have a few friends in Bamenda.} = J'ai quelques amis à Bamenda.
\end{itemize}
\end{propriete}

\courssub{La prononciation piégeuse}

\begin{propriete}[Trois difficultés phonétiques majeures]
\textbf{1. Le -ed final} des verbes réguliers au prétérit a trois prononciations :\par
— « \textbf{t} » après un son sourd (\eng{worked} « weurkt », \eng{stopped} « stopte ») ;\par
— « \textbf{d} » après un son sonore ou une voyelle (\eng{played} « plèïd », \eng{listened} « li-seund ») ;\par
— « \textbf{id} » seulement après \eng{t} ou \eng{d} (\eng{wanted} « won-tid », \eng{decided} « di-saï-did »).\par
\textbf{2. Le th} se prononce avec la langue entre les dents : \eng{think}, \eng{three} (sourd, comme un « s » sifflé avec la langue avancée) ; \eng{this}, \eng{mother} (sonore, comme un « z » doux). Ne dites jamais « sink » pour \eng{think}.\par
\textbf{3. Le h anglais est aspiré}, contrairement au français : \eng{house} « haousse », \eng{happy} « hapi », \eng{homework} « home-weurk ». Soufflez légèrement avant la voyelle.
\end{propriete}

\begin{popfigure}
\begin{tikzpicture}
\node[draw=popdark, fill=popblueL, rounded corners, font=\bfseries, minimum width=13.5cm, minimum height=0.8cm] at (0,0) {Erreurs fréquentes des francophones : le grand nettoyage};
\node[draw=popdark, fill=poppinkL, rounded corners, align=left, text width=6.2cm, anchor=north west] at (-6.75,-0.7) {\textbf{Erreur fréquente ✗}\\[2pt]
1. I am actually in Douala.\\
2. I need informations.\\
3. I am agree with you.\\
4. I don't know nothing.\\
5. She gave me many advices.\\
6. He married with her.};
\node[draw=popdark, fill=popgreenL, rounded corners, align=left, text width=6.2cm, anchor=north west] at (0.15,-0.7) {\textbf{Correction ✓}\\[2pt]
1. I am currently in Douala.\\
2. I need (some) information.\\
3. I agree with you.\\
4. I don't know anything.\\
5. She gave me a lot of advice.\\
6. He married her.};
\draw[-{Stealth}, very thick, poporange] (-0.4,-2.6) -- (0.3,-2.6);
\end{tikzpicture}
\legende{Six erreurs emblématiques des francophones et leurs corrections}
\end{popfigure}

\begin{exoresolu}[Corriger une phrase de copie]
Énoncé — La phrase suivante, extraite d'une copie de Terminale, contient cinq erreurs. Trouvez-les et corrigez-les :\par
\eng{Actually, the police is looking for informations because nobody didn't see nothing after the accident.}
\tcblower
Solution détaillée :\par
1. \eng{Actually} ✗ → \eng{Currently} ✓ : « actuellement » ne se dit pas \eng{actually} (faux ami).\par
2. \eng{the police is} ✗ → \eng{the police are} ✓ : \eng{police} est toujours pluriel en anglais.\par
3. \eng{informations} ✗ → \eng{information} ✓ : nom indénombrable, jamais de \eng{-s}.\par
4. \eng{nobody didn't see} ✗ → \eng{nobody saw} ✓ : double négation interdite ; \eng{nobody} porte déjà la négation, le verbe redevient affirmatif.\par
5. \eng{nothing} ✗ → \eng{anything} ✓ : après une phrase déjà négative, on utilise \eng{anything}.\par
Phrase corrigée : \eng{Currently, the police are looking for information because nobody saw anything after the accident.} = Actuellement, la police recherche des renseignements parce que personne n'a rien vu après l'accident.
\end{exoresolu}

\begin{methode}[La relecture en quatre questions pour l'épreuve de grammaire]
Avant de rendre votre copie, relisez chaque phrase en vous posant quatre questions, dans l'ordre :
\begin{enumerate}
\item \textbf{Le temps est-il correct ?} (cohérence avec le contexte : prétérit pour un récit passé, present perfect avec \eng{for/since}...)
\item \textbf{L'accord sujet-verbe est-il correct ?} (troisième personne avec \eng{-s}, \eng{everybody is}, \eng{the police are}...)
\item \textbf{La préposition est-elle correcte ?} (\eng{listen to}, \eng{wait for}, \eng{interested in}, \eng{good at}...)
\item \textbf{Le nom est-il dénombrable ou non ?} (pas de pluriel à \eng{information, advice, news} ; \eng{much} ou \eng{many} ?)
\end{enumerate}
Cette grille de quatre questions permet de repérer l'immense majorité des fautes avant le correcteur.
\end{methode}

\begin{savaistu}[L'article « the » et les noms abstraits]
Contrairement au français, l'anglais n'emploie \textbf{pas} l'article \eng{the} devant les noms abstraits pris dans un sens général. Le français dit « \textbf{La} vie est courte », l'anglais dit \eng{Life is short.} (et non « The life is short »). De même : \eng{Freedom is precious.} = La liberté est précieuse. \eng{Time is money.} = Le temps, c'est de l'argent. Mais dès que le nom est précisé, \eng{the} revient : \eng{The life of farmers is hard.} = La vie des agriculteurs est dure. C'est l'inverse exact du réflexe français : méfiez-vous de votre « la » automatique !
\end{savaistu}

\begin{exercice}[À vous de jouer]
Corrigez les phrases suivantes (une erreur par phrase) :
\begin{enumerate}
\item \eng{I will assist the concert at the stadium on Saturday.}
\item \eng{She bought new furnitures for her house in Buea.}
\item \eng{Everybody are invited to the meeting.}
\item \eng{He is listening the radio in his room.}
\item \eng{I don't have no money for the taxi.}
\item \eng{There are too much students in this class.}
\end{enumerate}
\end{exercice}

\begin{corrige}[Exercice « À vous de jouer »]
1. \eng{I will \textbf{attend} the concert...} = J'assisterai au concert. (\eng{assist} = aider.)\par
2. \eng{She bought new \textbf{furniture}...} = Elle a acheté de nouveaux meubles. (indénombrable.)\par
3. \eng{Everybody \textbf{is} invited...} = Tout le monde est invité. (singulier.)\par
4. \eng{He is listening \textbf{to} the radio...} = Il écoute la radio. (préposition \eng{to} obligatoire.)\par
5. \eng{I don't have \textbf{any} money...} = Je n'ai pas d'argent. (double négation interdite.)\par
6. \eng{There are too \textbf{many} students...} = Il y a trop d'élèves. (\eng{students} est dénombrable.)
\end{corrige}

\coursec{Entraînement type Bac : méthode et stratégies}

Vous savez maintenant repérer et corriger les pièges classiques des francophones : faux amis, prépositions rebelles, ordre des mots, temps mal choisis. Il reste une dernière étape, et non la moindre : transformer toutes ces connaissances en \cle{points} le jour de l'examen. Car au Baccalauréat A, la réussite ne dépend pas seulement de ce que vous savez, mais de \emph{la manière dont vous gérez l'épreuve}. Un candidat moyen mais méthodique obtient souvent une meilleure note qu'un bon élève désorganisé. Cette section vous donne la stratégie complète, minute par minute, puis un entraînement chronométré grandeur nature.

\courssub{Connaître l'épreuve : la structure de l'examen d'anglais au Bac A}

\begin{definition}[L'épreuve d'anglais au Baccalauréat A]
L'épreuve écrite d'anglais comporte généralement trois grandes parties :
\begin{enumerate}
\item \cle{Reading comprehension} : un texte en anglais (souvent sur un sujet de société : médias, éducation, environnement, jeunesse) suivi de questions de compréhension.
\item \cle{Grammar and vocabulary} : des exercices de langue (phrases à compléter, transformations, mise à la forme passive ou au discours indirect, vocabulaire à retrouver).
\item \cle{Composition} : une production écrite (essay ou lettre) d'environ 150 à 250 mots sur un sujet lié au thème du texte.
\end{enumerate}
\end{definition}

\begin{attention}[Ce que l'examinateur attend vraiment]
L'examinateur ne cherche pas à vous piéger : il vérifie trois choses — vous \emph{comprenez} un texte authentique, vous \emph{maîtrisez} les structures de la langue, et vous savez \emph{exprimer une opinion} en anglais correct. Chaque partie se prépare avec une méthode précise. Improviser, c'est perdre des points inutilement.
\end{attention}

\courssub{Gérer son temps : le plan de bataille en 90 minutes}

La gestion du temps est la compétence invisible de l'examen. Voici le découpage recommandé pour une épreuve de 1 h 30 :

\begin{center}
\begin{tabularx}{\linewidth}{|l|X|X|}
\hline
\rowcolor{popblueL} \textbf{Temps} & \textbf{Étape} & \textbf{Objectif} \\
\hline
10 min & Lecture du texte & Comprendre le sens global, repérer le thème et la structure \\
\hline
20 min & Questions de compréhension & Répondre avec ses propres mots, phrases complètes \\
\hline
20 min & Grammaire et vocabulaire & Identifier la structure testée, appliquer la règle, vérifier \\
\hline
30 min & Composition & Plan rapide (5 min), rédaction (20 min), relecture (5 min) \\
\hline
10 min & Relecture générale & Orthographe, accords, ponctuation, majuscules \\
\hline
\end{tabularx}
\end{center}

\begin{aretenir}[Les trois règles d'or du temps]
\begin{itemize}
\item \textbf{Ne bloquez jamais} sur une question difficile : passez et revenez-y à la fin.
\item \textbf{Réservez toujours} les 10 dernières minutes à la relecture : une copie relue gagne en moyenne 1 à 2 points.
\item \textbf{Respectez la hiérarchie des points} : la composition rapporte beaucoup ; ne la sacrifiez jamais faute de temps.
\end{itemize}
\end{aretenir}

\courssub{Méthode de compréhension : lire intelligemment}

\begin{methode}[Comprendre un texte en 4 étapes]
\begin{enumerate}
\item \textbf{Lisez d'abord les questions}, avant le texte. Votre cerveau saura ainsi ce qu'il cherche.
\item \textbf{Repérez les mots-clés} des questions (noms propres, dates, verbes d'action, adjectifs d'opinion).
\item \textbf{Localisez les passages} du texte qui contiennent ces mots-clés ou leurs synonymes, puis relisez ces passages attentivement.
\item \textbf{Rédigez la réponse avec vos propres mots}, en phrase complète, en réutilisant la structure de la question : \eng{Why do young people prefer social media?} → \eng{Young people prefer social media because they provide fast and free news.}
\end{enumerate}
\end{methode}

\begin{attention}[Le piège du copier-coller]
Ne recopiez \textbf{jamais} de longues phrases du texte comme réponse de compréhension. L'examinateur attend une \cle{reformulation} : recopier prouve que vous avez trouvé le passage, mais pas que vous l'avez compris. Changez les mots, changez la structure, gardez le sens. Mauvais : \eng{“Fake news spreads faster than verified information.”} (recopié). Bon : \eng{False information circulates more quickly than checked facts.} (reformulé).
\end{attention}

\courssub{Méthode de grammaire : diagnostiquer avant de soigner}

Face à un exercice de grammaire, le réflexe du bon candidat est de se poser une question : \emph{quelle structure est testée ici ?}

\begin{methode}[Résoudre un exercice de grammaire en 3 temps]
\begin{enumerate}
\item \textbf{Identifier} la structure testée : est-ce un temps verbal (indices : \eng{yesterday}, \eng{since}, \eng{already}) ? Un modal (nuance d'obligation, de probabilité) ? Le passif (l'objet devient sujet) ? Le discours rapporté (recul des temps) ? Un conditionnel (\eng{if}) ?
\item \textbf{Appliquer la règle} correspondante, sans se laisser influencer par le français.
\item \textbf{Vérifier l'accord} : sujet-verbe (\eng{he goes}, pas \eng{he go}), temps cohérents, forme du participe passé.
\end{enumerate}
\end{methode}

\begin{exemplebox}[Diagnostic express]
\eng{If I had money, I would buy a smartphone.} → structure testée : \cle{conditionnel de type 2} (if + prétérit, would + base verbale) ; « Si j'avais de l'argent, j'achèterais un smartphone. »\par
\eng{The match has been won by our team.} → structure testée : \cle{passif au present perfect} ; « Le match a été gagné par notre équipe. »\par
\eng{She said she was tired.} → structure testée : \cle{discours rapporté} avec recul du temps (is → was) ; « Elle a dit qu'elle était fatiguée. »
\end{exemplebox}

\courssub{Méthode de composition : construire avant d'écrire}

\begin{methode}[Réussir sa composition en 4 étapes]
\begin{enumerate}
\item \textbf{Reformuler le sujet} dans l'introduction : montrez que vous l'avez compris, sans le recopier mot à mot.
\item \textbf{Annoncer le plan} en une phrase : \eng{This essay will first examine the advantages of social media, then discuss their dangers, and finally give a personal opinion.}
\item \textbf{Développer} : un paragraphe = une idée + un exemple. Deux ou trois paragraphes suffisent.
\item \textbf{Conclure} en donnant votre opinion \emph{nuancée} : \eng{In conclusion, although social media have undeniable benefits, I believe that...}
\end{enumerate}
\end{methode}

\begin{exemplebox}[Un sujet type analysé]
Sujet : \eng{Do you think social media do more harm than good?} — « Pensez-vous que les réseaux sociaux font plus de mal que de bien ? »\par
Plan gagnant en 3 parties :
\begin{itemize}
\item \textbf{Avantages} : information rapide et gratuite, lien avec la famille et les amis (\eng{instant news, keeping in touch}) ;
\item \textbf{Dangers} : fake news, addiction, cyberbullying (\eng{false information, wasting time}) ;
\item \textbf{Opinion personnelle} : les réseaux sont un outil ; tout dépend de l'usage qu'on en fait (\eng{it depends on how we use them}).
\end{itemize}
\end{exemplebox}

\begin{aretenir}[Les connecteurs qui font gagner des points]
\begin{center}
\begin{tabularx}{\linewidth}{|X|X|X|}
\hline
\rowcolor{popblueL} \textbf{Fonction} & \textbf{English} & \textbf{Français} \\
\hline
Ajout & \eng{moreover, furthermore, in addition} & de plus, en outre \\
\hline
Opposition & \eng{however, nevertheless, whereas} & cependant, néanmoins, tandis que \\
\hline
Cause & \eng{because, since, as} & parce que, puisque \\
\hline
Conséquence & \eng{therefore, consequently, as a result} & donc, par conséquent \\
\hline
Exemple & \eng{for example, for instance} & par exemple \\
\hline
Conclusion & \eng{in conclusion, to sum up, on the whole} & en conclusion, en résumé \\
\hline
\end{tabularx}
\end{center}
\end{aretenir}

\begin{savaistu}[Le secret des correcteurs]
Les correcteurs du Bac valorisent les connecteurs variés (\eng{however, moreover, consequently}) et les structures complexes \emph{correctes} (conditionnels, passif, discours rapporté). Mais attention : \textbf{mieux vaut une phrase simple juste qu'une phrase complexe fausse}. \eng{Social media are useful, but they can be dangerous.} vaut mieux qu'une phrase alambiquée pleine de fautes. La correction prime toujours sur l'ambition.
\end{savaistu}

\courssub{La vérification finale : les 10 minutes qui sauvent des points}

\begin{attention}[Check-list de relecture]
\begin{itemize}
\item \textbf{Homophones} : \eng{their} (leur) / \eng{there} (là) / \eng{they're} (ils sont) ; \eng{its} (son, sa) / \eng{it's} (c'est, il est).
\item \textbf{Majuscules} : début de phrase, \eng{I} toujours en majuscule, langues et nationalités (\eng{English, Cameroonian}).
\item \textbf{Ponctuation} : pas d'espace avant le point ou la virgule en anglais ; apostrophe du génitif (\eng{my father's shop}).
\item \textbf{Accords} : -s à la 3\textsuperscript{e} personne du singulier (\eng{she works}), pluriels réguliers.
\item \textbf{Faux amis} : \eng{actually} = en fait (pas « actuellement »), \eng{sensible} = raisonnable, \eng{a library} = une bibliothèque (pas une librairie).
\end{itemize}
\end{attention}

\courssub{Entraînement chronométré : annale sur les médias}

Lisez le texte suivant en 10 minutes maximum, puis répondez aux questions. Chronométrez-vous !

\begin{textebox}[Reading text — “The New Face of News in Cameroon”]
Social media have transformed communication in Cameroon. Young people in Yaoundé and Bamenda now get their news from Facebook and WhatsApp rather than from traditional newspapers. A student in Buea can learn about a football result, a political event or a road accident within seconds, simply by checking her phone. For many families, the smartphone has replaced the radio as the main source of information.

However, this revolution has a dark side: fake news spreads faster than verified information. During elections, false rumours circulate widely and create tension between communities. Some messages are shared thousands of times before anybody checks whether they are true. Media experts say that citizens should be trained to check sources before sharing content. Schools, they argue, must teach young people how to recognise reliable websites and how to doubt sensational headlines.

The government and several private organisations have started awareness campaigns. Posters in cybercafés remind users that sharing a false message can have serious consequences. Journalists also insist that professional media must remain the reference: their work is checked, corrected and verified before publication.

If nothing is done, misinformation could threaten social peace. But if young Cameroonians learn to think critically, social media could become a powerful tool for democracy rather than a danger. The choice, experts conclude, is in the hands of the users themselves.
\end{textebox}

\textbf{Glossaire}\par
\begin{itemize}
\item \eng{fake news} (n.) : fausses informations ; \eng{to spread} (v.) : se propager ; \eng{verified} (adj.) : vérifié
\item \eng{a rumour} (n.) : une rumeur ; \eng{reliable} (adj.) : fiable ; \eng{a headline} (n.) : un titre (de presse)
\item \eng{awareness campaign} (n.) : campagne de sensibilisation ; \eng{to threaten} (v.) : menacer
\end{itemize}

\begin{exercice}[Annale — 10 questions (30 minutes)]
\textbf{A. Comprehension.} Answer in your own words.
\begin{enumerate}
\item \eng{Where do young Cameroonians get their news nowadays?}
\item \eng{What has replaced the radio in many families?}
\item \eng{What is the “dark side” of the social media revolution?}
\item \eng{According to media experts, what should citizens do before sharing content?}
\item \eng{What could happen if nothing is done about misinformation?}
\end{enumerate}
\textbf{B. Grammar and vocabulary.}
\begin{enumerate}
\setcounter{enumi}{5}
\item \eng{Put into the passive voice: Young people share false messages.}
\item \eng{Report the speech: The expert said: “Citizens must check sources.”}
\item \eng{Complete with the right tense: If schools taught media literacy, young people (be) better prepared.}
\item \eng{Find a synonym in the text for “false information”.}
\item \eng{Find the opposite of “to hide information” (use a verb from the text).}
\end{enumerate}
\end{exercice}

\begin{corrige}[Annale — 10 questions]
\textbf{A.} 1. \eng{They get their news from social media such as Facebook and WhatsApp.} 2. \eng{The smartphone has replaced the radio.} 3. \eng{The dark side is that false information circulates more quickly than checked facts.} 4. \eng{They should verify the sources of the content they share.} 5. \eng{Misinformation could endanger peace in society.}\par
\textbf{B.} 6. \eng{False messages are shared by young people.} 7. \eng{The expert said (that) citizens had to check sources.} (recul : must → had to). 8. \eng{If schools taught media literacy, young people would be better prepared.} (conditionnel type 2). 9. \eng{Fake news} / \eng{misinformation}. 10. \eng{To share} (partager, diffuser).
\end{corrige}

\begin{exoresolu}[Composition guidée — sujet type Bac]
Énoncé : \eng{“Social media do more harm than good.” Do you agree? Write an essay of about 150 words.}
\tcblower
\textbf{Solution — modèle de rédaction commenté :}\par
\eng{Nowadays, most young Cameroonians use social media every day. Some people believe these platforms are dangerous, while others see them as a useful tool. This essay will examine both sides before giving a personal opinion.}\par
\eng{On the one hand, social media provide fast and free information. For example, a student in Bamenda can follow the news or contact his family in Douala within seconds.}\par
\eng{On the other hand, fake news spreads quickly and can create tension. Moreover, many teenagers waste hours on their phones instead of studying.}\par
\eng{In conclusion, although social media have real dangers, I believe they do more good than harm, provided that users learn to check what they read.}\par
Points forts du modèle : connecteurs variés (\eng{nowadays, on the one hand, moreover, although, in conclusion}), un conditionnel implicite (\eng{provided that}), phrases courtes et correctes.
\end{exoresolu}

\courssub{Stratégie d'examen : la synthèse visuelle}

\begin{popfigure}
\begin{center}
\begin{tikzpicture}[
  node distance=0.55cm,
  etape/.style={rectangle, rounded corners=4pt, draw=popblue, fill=popblueL, text width=4.6cm, align=center, minimum height=0.9cm, font=\small},
  fleche/.style={-{Stealth[length=3mm]}, thick, poporange}
]
\node[etape] (n1) {\textbf{1. Lire les questions} avant le texte};
\node[etape, below=of n1] (n2) {\textbf{2. Repérer les mots-clés} (noms, dates, verbes)};
\node[etape, below=of n2] (n3) {\textbf{3. Localiser} les passages dans le texte};
\node[etape, below=of n3] (n4) {\textbf{4. Rédiger} avec ses propres mots};
\node[etape, below=of n4] (n5) {\textbf{5. Relire} : orthographe, accords, majuscules};
\draw[fleche] (n1) -- (n2);
\draw[fleche] (n2) -- (n3);
\draw[fleche] (n3) -- (n4);
\draw[fleche] (n4) -- (n5);
\end{tikzpicture}
\end{center}
\legende{Organigramme de la stratégie d'examen : cinq étapes, dans l'ordre, sans exception.}
\end{popfigure}

\courssub{Auto-évaluation : votre bilan avant le grand jour}

\begin{experience}[Ma grille d'auto-évaluation]
Après l'entraînement chronométré, remplissez honnêtement cette grille et notez un point d'action pour chaque faiblesse :
\begin{center}
\begin{tabularx}{\linewidth}{|X|c|c|X|}
\hline
\rowcolor{popblueL} \textbf{Compétence} & \textbf{OK} & \textbf{À revoir} & \textbf{Point d'action} \\
\hline
Comprendre le sens global d'un texte & & & \\
\hline
Répondre sans recopier le texte & & & \\
\hline
Identifier la structure grammaticale testée & & & \\
\hline
Temps verbaux (present perfect, conditionnels) & & & \\
\hline
Passif et discours rapporté & & & \\
\hline
Rédiger une composition en 3 parties & & & \\
\hline
Utiliser des connecteurs variés & & & \\
\hline
Éviter les faux amis et homophones & & & \\
\hline
Gérer mon temps (90 minutes) & & & \\
\hline
\end{tabularx}
\end{center}
Discutez de votre grille avec un camarade ou votre professeur : un point faible identifié aujourd'hui est un point gagné demain.
\end{experience}

\begin{aretenir}[Le mot de la fin]
Le jour de l'examen, rappelez-vous ces cinq vérités : lisez les questions \emph{avant} le texte ; reformulez toujours vos réponses ; une phrase simple et juste vaut mieux qu'une phrase complexe et fausse ; gardez 10 minutes pour relire ; et surtout, restez calme — vous avez révisé toute l'année, vous êtes prêts. \eng{Good luck!}
\end{aretenir}

\newpage

\coursec{Activité d'intégration}

\coursec{Lesson 53 — General Revision: Revise previous aspects of grammar}

\courssub{Objectifs de l'activité}
Cette activité de synthèse vise à :
\begin{itemize}
    \item \textbf{Mobiliser l'ensemble des acquis grammaticaux} de l'année (temps, modaux, passif, discours rapporté, conditionnels, phrases complexes) dans une situation d'examen blanc chronométrée.
    \item \textbf{Consolider les stratégies de compréhension} (lecture et écoute simulée) et de production écrite (composition argumentée) selon les attendus du Baccalauréat.
    \item \textbf{Développer l'auto-évaluation et la correction entre pairs} à l'aide d'une grille critériée, pour identifier les erreurs récurrentes et les corriger par l'explication de la règle.
    \item \textbf{Établir collectivement une « Top 10 des erreurs à ne plus commettre »} comme outil de révision final jusqu'à l'examen.
\end{itemize}

\courssub{Situation}
Vous êtes élèves en Terminale A au Lycée Bilingue de Buea. Dans le cadre de la préparation au Baccalauréat, votre professeur d'anglais organise un \textbf{Mock Exam (examen blanc)} final. Le sujet porte sur un thème d'actualité brûlant au Cameroun : \emph{la régulation des réseaux sociaux}. Vous disposez de \textbf{45 minutes} pour traiter l'ensemble du sujet (compréhension, grammaire, vocabulaire, composition). Le texte support est un article d'opinion publié dans un journal jeunesse camerounais anglophone, \textit{The Youth Voice}.

\begin{textebox}[Article : « Social Media in Cameroon: A Double-Edged Sword » — The Youth Voice, Yaoundé Edition]
Social media has revolutionised the way young Cameroonians communicate, learn and do business. In cities like Douala, Yaoundé and Buea, it is common to see students using WhatsApp groups to share lesson notes, or young entrepreneurs advertising their products on Facebook and TikTok. The internet has become a powerful tool for financial inclusion; many youths who were previously unemployed now earn a living as content creators or digital marketers.

However, this digital boom has a dark side. Cyberbullying, the spread of fake news and online scams (locally known as \eng{« 419 »}) have increased dramatically. Last year, a false rumour about a cholera outbreak in Bamenda spread like wildfire on social media, causing panic and disrupting schools. Moreover, addiction to screens affects academic performance; many candidates stay up late scrolling through videos instead of revising their lessons.

The government has proposed a new bill to regulate digital platforms. The law would require users to register with their real identities and impose heavy fines on those who publish hate speech or false information. While some citizens welcome this move as a necessary step to protect society, others fear it could restrict freedom of expression. A student from the University of Buea declared during a radio programme: “If the government controls what we post, we will lose our only space to criticise bad governance.”

The debate is far from over. Stakeholders — parents, teachers, telecom operators and civil society — must find a balance between security and liberty. As the Minister of Posts and Telecommunications recently stated, “We cannot ban the internet, but we must teach our youth how to use it responsibly.” The future of Cameroon's digital space depends on the choices made today.
\end{textebox}

\courssub{Consignes}
\emph{Durée totale : 45 minutes. Travail individuel dans un premier temps, puis correction en groupes de 4.}

\begin{enumerate}
    \item \textbf{Compréhension de l'écrit (10 marks)} \\
    Lisez le texte ci-dessus et répondez aux questions suivantes en anglais.
    \begin{enumerate}
        \item \textbf{True / False / Not Mentioned} — Cochez la bonne case et justifiez par une citation du texte.
        \begin{itemize}
            \item[a)] Social media has created job opportunities for young Cameroonians.
            \item[b)] The government has already passed the law regulating digital platforms.
            \item[c)] The student from the University of Buea supports the new bill.
        \end{itemize}
        \item \textbf{Multiple Choice} — Choisissez la bonne réponse.
        \begin{itemize}
            \item[d)] The phrase “spread like wildfire” (paragraph 2) means:
            \begin{enumerate}
                \item spread slowly
                \item spread very quickly
                \item was controlled by firefighters
            \end{enumerate}
            \item[e)] According to the Minister, the solution is to:
            \begin{enumerate}
                \item ban the internet completely
                \item educate the youth on responsible use
                \item arrest all content creators
            \end{enumerate}
        \end{itemize}
        \item \textbf{Short Answers} — Répondez en phrases complètes.
        \begin{itemize}
            \item[f)] What are two negative effects of social media mentioned in the text?
            \item[g)] What does the proposed bill require from users?
        \end{itemize}
    \end{enumerate}

    \item \textbf{Grammaire (8 marks)} \\
    Complétez les phrases suivantes en conjuguant le verbe entre parenthèses au temps approprié, en utilisant le modal ou la structure demandée, ou en transformant la phrase.
    \begin{enumerate}
        \item By the time the Minister \eng{(arrive)} tomorrow, the debate \eng{(already / start)}. \textit{(Future Perfect / Future Simple)}
        \item The government \eng{(discuss)} the bill for three months now, but no decision \eng{(take)} yet. \textit{(Present Perfect Continuous / Passive)}
        \item If the students \eng{(know)} the dangers of fake news earlier, they \eng{(not / share)} that rumour last year. \textit{(Third Conditional)}
        \item The teacher asked the students \eng{(not / use)} their phones during the exam. \textit{(Reported Speech / Infinitive)}
        \item Social media \eng{(must / regulate)} to protect users, but freedom of expression \eng{(must / preserve)}. \textit{(Passive Modals)}
        \item I wish the internet \eng{(be)} cheaper in rural areas. \textit{(Wish + Past Simple)}
        \item The rumour \eng{(spread)} so fast that schools \eng{(close)} for two days. \textit{(Past Simple / Past Perfect)}
        \item Not only \eng{(the bill / aim)} to fight cybercrime, but it \eng{(also / protect)} children online. \textit{(Inversion / Emphasis)}
    \end{enumerate}

    \item \textbf{Vocabulaire (5 marks)} \\
    Choisissez le mot qui convient parmi les propositions (attention aux faux amis et aux prépositions).
    \begin{enumerate}
        \item The government wants to \eng{(control / check / survey)} hate speech on the internet.
        \item Many youths are \eng{(addicted to / dependent of / sensible to)} their smartphones.
        \item The fake news caused a huge \eng{(scandal / damage / prejudice)} to the reputation of the hospital.
        \item Parents should \eng{(sensibilise / raise awareness / make conscious)} their children about online risks.
        \item The bill has been \eng{(discussed / argued / debated)} by MPs for weeks.
    \end{enumerate}

    \item \textbf{Expression écrite — Composition (12 marks)} \\
    \textbf{Sujet :} “Should social media be regulated in Cameroon?” \\
    Rédigez une composition argumentée d'environ \textbf{120 mots}. Structurez votre texte : Introduction (accroche + thèse), Développement (2 arguments avec exemples camerounais), Conclusion (ouverture). Respectez la ponctuation et les connecteurs logiques.
\end{enumerate}

\courssub{Ressources à mobiliser}
\emph{Rappel synthétique des règles clés pour réussir l'activité.}

\begin{propriete}[La carte des temps verbaux — Révision express]
\begin{center}
\begin{tabularx}{\linewidth}{|X|X|X|}
\hline
\rowcolor{popblueL} \textbf{Temps / Structure} & \textbf{Forme (Affirmative)} & \textbf{Emploi clé / Marqueurs} \\ \hline
Present Simple & S + V (s/es) & Habitudes, vérités générales (\eng{every day, usually}) \\ \hline
Present Perfect & S + have/has + V-en & Passé récent, expérience, lien présent (\eng{just, already, since, for}) \\ \hline
Past Simple & S + V-ed / irr. & Action terminée dans le passé (\eng{yesterday, last year, in 2023}) \\ \hline
Future (Will / Be going to) & S + will / be going to + V & Prédiction / Intention (\eng{tomorrow, next week}) \\ \hline
Conditionnel 1 (Real) & If + Present, \eng{will} + V & Possible / Probable (\eng{If it rains, I will stay}) \\ \hline
Conditionnel 2 (Unreal) & If + Preterit, \eng{would} + V & Hypothétique présent (\eng{If I were you, I would go}) \\ \hline
Conditionnel 3 (Past Unreal) & If + Past Perfect, \eng{would have} + V-en & Regret / Hypothèse passée (\eng{If I had known, I would have come}) \\ \hline
Wish / If only & \eng{Wish} + Past / Past Perfect & Regret présent / passé (\eng{I wish I knew / I wish I had known}) \\ \hline
Passif & S + be + V-en & Focus sur l'objet / agent inconnu (\eng{The bill was passed}) \\ \hline
Discours rapporté & Décalage des temps (Backshift) & \eng{He said he was tired} (au lieu de \eng{I am tired}) \\ \hline
Modaux (Passif) & Modal + be + V-en & Obligation / Conseil passif (\eng{It must be done}) \\ \hline
Inversion & Aux + S + V & Emphase / \eng{Never, Not only, Rarely} en tête \\ \hline
\end{tabularx}
\end{center}
\end{propriete}

\begin{attention}[Pièges « Spécial Francophones » — Le Top 10 à surveiller]
\begin{enumerate}
    \item \textbf{Since / For} : \eng{since} + point de départ (\eng{since 2020}) ; \eng{for} + durée (\eng{for three years}). \textit{Pas : « since three years »}.
    \item \textbf{Present Perfect vs Past Simple} : \eng{I have seen him yesterday} $\rightarrow$ \textbf{FAUX}. \eng{I saw him yesterday}.
    \item \textbf{Modaux + infinitif sans TO} : \eng{He must to go} $\rightarrow$ \textbf{FAUX}. \eng{He must go}. (Excepté \eng{ought to}, \eng{have to}, \eng{be able to}).
    \item \textbf{Passif : Oublier le BE} : \eng{The law passed} (actif) vs \eng{The law \textbf{was} passed} (passif).
    \item \textbf{Discours rapporté : Backshift} : \eng{will} $\rightarrow$ \eng{would} ; \eng{can} $\rightarrow$ \eng{could} ; \eng{Present} $\rightarrow$ \eng{Past} ; \eng{Past} $\rightarrow$ \eng{Past Perfect}.
    \item \textbf{Conditionnel 3} : \eng{If I would have known} $\rightarrow$ \textbf{FAUX}. \eng{If I \textbf{had} known}.
    \item \textbf{Faux amis} : \eng{Control} = contrôler (vérifier/diriger) ; \eng{Check} = vérifier ; \eng{Survey} = enquêter. \eng{Damage} (dommage, invariable) $\neq$ \eng{Prejudice} (préjugé / dommage juridique). \eng{Raise awareness} $\neq$ \eng{Sensibiliser} (n'existe pas en anglais).
    \item \textbf{Prépositions après adjectifs/verbes} : \eng{Addicted TO}, \eng{Dependent ON}, \eng{Interested IN}, \eng{Good AT}, \eng{Responsible FOR}.
    \item \textbf{Inversion} : \eng{Not only did he come} (auxiliaire \eng{did} + sujet + base verbale), pas \eng{Not only he came}.
    \item \textbf{Concordance des temps (Sequence of Tenses)} : Verbe principal au passé $\rightarrow$ subordonnée au passé (sauf vérité générale).
\end{enumerate}
\end{attention}

\courssub{Proposition de production et corrigé commenté}
\begin{exoresolu}[Corrigé complet du Mock Exam — The Final Mock Exam Challenge]
\textbf{PARTIE 1 : COMPRÉHENSION (10 marks)}

\begin{enumerate}
    \item[a)] \textbf{True.} \eng{« many youths who were previously unemployed now earn a living as content creators or digital marketers. »}
    \item[b)] \textbf{False.} \eng{« The government has proposed a new bill... »} (proposed $\neq$ passed).
    \item[c)] \textbf{False.} \eng{« If the government controls what we post, we will lose our only space to criticise bad governance. »} (Il s'y oppose).
    \item[d)] \textbf{2. spread very quickly.} (Idiome : \eng{spread like wildfire} = se propager comme une traînée de poudre).
    \item[e)] \textbf{2. educate the youth on responsible use.} \eng{« We must teach our youth how to use it responsibly. »}
    \item[f)] \textbf{Cyberbullying and the spread of fake news} (aussi acceptés : online scams / addiction affecting academic performance).
    \item[g)] \textbf{The bill requires users to register with their real identities and imposes heavy fines on those who publish hate speech or false information.}
\end{enumerate}

\textbf{PARTIE 2 : GRAMMAIRE (8 marks)}

\begin{enumerate}
    \item \eng{By the time the Minister \textbf{arrives} tomorrow, the debate \textbf{will already have started}.} \\
    \textit{Règle : Après \eng{by the time} + présent (valeur future), la principale est au Future Perfect (\eng{will have + V-en}). Le verbe \eng{arrive} prend \eng{s} (3e pers. sg).}
    \item \eng{The government \textbf{has been discussing} the bill for three months now, but no decision \textbf{has been taken} yet.} \\
    \textit{Règle : \eng{for three months} + action en cours $\rightarrow$ Present Perfect Continuous (\eng{have been + V-ing}). Passif : \eng{no decision} (sujet) + \eng{has been taken} (Present Perfect Passive).}
    \item \eng{If the students \textbf{had known} the dangers of fake news earlier, they \textbf{would not have shared} that rumour last year.} \\
    \textit{Règle : Third Conditional (irréel du passé). If + Past Perfect (\eng{had known}), \eng{would have + V-en} (\eng{would not have shared}).}
    \item \eng{The teacher asked the students \textbf{not to use} their phones during the exam.} \\
    \textit{Règle : Discours rapporté avec ordre/négation. \eng{ask/tell/order + objet + not to + BV}. Pas de \eng{that} + conjugaison.}
    \item \eng{Social media \textbf{must be regulated} to protect users, but freedom of expression \textbf{must be preserved}.} \\
    \textit{Règle : Modaux passifs. \eng{Modal + be + V-en}. \eng{Regulate} et \eng{preserve} sont des verbes d'action subis par les sujets.}
    \item \eng{I wish the internet \textbf{were} cheaper in rural areas.} \\
    \textit{Règle : \eng{Wish} + Past Simple (subjonctif) pour un regret présent. \eng{Were} pour toutes les personnes (forme formelle correcte au Bac).}
    \item \eng{The rumour \textbf{spread} so fast that schools \textbf{closed} for two days.} \\
    \textit{Règle : \eng{Spread} (irrégulier : spread-spread-spread) au Past Simple. La subordonnée de conséquence introduite par \eng{so... that} suit la chronologie : la rumeur s'est propagée \textbf{puis} les écoles ont fermé $\rightarrow$ Past Simple (\eng{closed}). \eng{Had closed} indiquerait une antériorité illogique (écoles fermées avant la rumeur).}
    \item \eng{\textbf{Not only does the bill aim} to fight cybercrime, but it \textbf{also protects} children online.} \\
    \textit{Règle : Inversion après \eng{Not only} en tête de phrase. Auxiliaire \eng{does} (Present Simple, 3e pers. sg) + sujet (\eng{the bill}) + BV (\eng{aim}). La seconde clause reste à l'ordre normal (\eng{it also protects}).}
\end{enumerate}

\textbf{PARTIE 3 : VOCABULAIRE (5 marks)}

\begin{enumerate}
    \item \textbf{Control} (dans le sens « réguler, maîtriser, exercer un contrôle sur »). \eng{Check} = vérifier ponctuellement ; \eng{Survey} = enquêter / sonder.
    \item \textbf{Addicted to} (préposition obligatoire après \eng{addicted}). \eng{Dependent ON} ; \eng{Sensitive TO}.
    \item \textbf{Damage} (invariable, signifie « dommage, préjudice matériel/moral »). \eng{Scandal} = scandale ; \eng{Prejudice} = préjugé ou dommage juridique spécifique.
    \item \textbf{Raise awareness} (locution verbale standard pour « sensibiliser »). \eng{Sensibiliser} n'existe pas ; \eng{Make conscious} ne s'utilise pas ainsi.
    \item \textbf{Debated} (verbe transitif direct : \eng{debate a bill}). \eng{Discuss} est aussi transitif direct (\eng{discuss the bill}), mais \eng{debate} est plus formel pour un texte de loi. \eng{Argued} implique une dispute et s'utilise avec \eng{about/over}.
\end{enumerate}

\textbf{PARTIE 4 : COMPOSITION — PROPOSITION MODÈLE (≈ 125 mots)}

\begin{textebox}[Modèle de composition : Should social media be regulated in Cameroon?]
Social media has become the new public square for Cameroonian youth. From WhatsApp groups in Douala to TikTok shops in Buea, it drives innovation and free speech. However, the rise of cyberbullying, fake news and the infamous “419” scams threatens our social fabric. Therefore, I firmly believe that regulation is necessary, provided it protects liberty.

Firstly, a legal framework would hold perpetrators accountable. Last year, a false cholera rumour in Bamenda caused panic because the author remained anonymous. If the proposed bill requiring real‑identity registration had been in place, the culprit would have been traced. Secondly, regulation can enforce digital literacy. The Minister rightly stated we must teach responsible use; schools should integrate media education to help students spot fake news.

Opponents fear censorship, arguing that the government might silence critics. This risk is real, but the solution is not anarchy. An independent body, not the state alone, should oversee the law.

In conclusion, smart regulation — not a ban — is the price of a safe digital Cameroon. Let us build a space where freedom walks hand in hand with responsibility.
\end{textebox}

\textbf{COMMENTAIRE DES CHOIX DE LANGUE (Pour la grille de correction collaborative)}

\begin{itemize}
    \item \textbf{Structure :} Introduction (accroche + thèse \eng{I firmly believe that...}), 2 arguments distincts (\eng{Firstly... Secondly...}), concession (\eng{Opponents fear... This risk is real, but...}), Conclusion (\eng{In conclusion...}) avec ouverture.
    \item \textbf{Grammaire mobilisée :}
    \begin{itemize}
        \item \eng{Present Perfect} (\eng{has become, has driven}) pour situation actuelle issue du passé.
        \item \eng{Third Conditional} (\eng{If the bill had been in place, the culprit would have been traced}) : hypothèse passée irréelle + passif modal parfait.
        \item \eng{Modaux} (\eng{should integrate, must teach, might silence}) : obligation, conseil, possibilité.
        \item \eng{Passif} (\eng{should be overseen, would have been traced}) : focus sur l'action/objet.
        \item \eng{Relative clauses} (\eng{where freedom walks...}) : fluidité.
    \end{itemize}
    \item \textbf{Vocabulaire précis :} \eng{Public square, social fabric, perpetrators, accountable, culprit, censorship, anarchy, independent body, oversight}.
    \item \textbf{Connecteurs :} \eng{However, Therefore, Firstly, Secondly, Moreover, Opponents argue, In conclusion}.
    \item \textbf{Contexte camerounais :} \eng{Douala, Buea, Bamenda, 419, Minister, WhatsApp, TikTok} — ancre le propos dans la réalité du candidat.
\end{itemize}
\end{exoresolu}

\courssub{Bilan de l'activité}
\begin{enumerate}
    \item \textbf{Échange et correction (15 min) :} Chaque groupe échange sa copie. Avec la grille fournie (barème + règles de référence), les correcteurs annotent : \eng{✓} (correct), \eng{✗ G} (erreur grammaire), \eng{✗ V} (vocabulaire), \eng{✗ S} (syntaxe/ordre), \eng{?} (incompréhensible). \textbf{Chaque erreur doit être justifiée par le numéro de la règle} (ex. \eng{✗ G3 : modal + to}).
    \item \textbf{Mise en commun (10 min) :} L'enseignant relève au tableau les 10 erreurs les plus fréquentes dans la classe. La classe vote pour le « Top 10 » final.
    \item \textbf{Affichage :} Le « Top 10 des erreurs à ne plus commettre » est rédigé sur une grande affiche (format A2) avec la règle, l'exemple fautif et la correction. Exemple :
    \begin{center}
    \begin{tabular}{|p{3cm}|p{5cm}|p{5cm}|}
    \hline
    \rowcolor{poporangeL} \textbf{N°} & \textbf{Erreur type (Francophone)} & \textbf{Correction + Règle} \\ \hline
    1 & \eng{He must to go} & \eng{He must go} — Modal sans TO (sauf ought to) \\ \hline
    2 & \eng{Since three years} & \eng{For three years} — \eng{For} + durée \\ \hline
    3 & \eng{If I would have known} & \eng{If I had known} — Pas de \eng{would} dans le \eng{if} \\ \hline
    ... & ... & ... \\ \hline
    \end{tabular}
    \end{center}
    \item \textbf{Engagement :} Chaque élève recopie son « Top 3 personnel » (ses 3 erreurs récurrentes) sur une fiche bristol à garder jusqu'au Bac.
\end{enumerate}

\begin{savaistu}[Le Bac anglais au Cameroun : Épreuve écrite]
L'épreuve écrite de l'anglais au Baccalauréat (Séries A) dure \textbf{3 heures}. Elle comprend généralement :
\begin{itemize}
    \item \textbf{Compréhension (20 pts) :} Un texte (300-400 mots) + 10 questions (V/F, QCM, Questions ouvertes, Vocabulaire en contexte, Référence).
    \item \textbf{Grammaire / Vocabulaire (20 pts) :} 10-15 items de transformation, mise à la forme, choix multiples, correction d'erreurs.
    \item \textbf{Expression écrite (20 pts) :} Une composition guidée (lettre, article, discours, essai argumentatif) de 150-200 mots. Les critères : Contenu/Idées (8), Organisation/Coherence (4), Grammaire/Structures (4), Vocabulaire/Orthographe (4).
\end{itemize}
\emph{Astuce :} Gérez votre temps ! 45 min compréhension, 45 min grammaire/vocab, 1h composition + 15 min relecture.
\end{savaistu}

\newpage

\coursec{Méthodes et exercices résolus}

\courssub{Savoir-faire 1 : Mobiliser ses acquis de l'année}

\begin{methode}[Réviser efficacement toute la grammaire de l'année]
\begin{enumerate}
\item \textbf{Faire l'inventaire.} Liste les grands chapitres étudiés : temps verbaux (present, past, future, perfect), modaux, voix passive, discours rapporté, conditionnels, comparatifs, relatives, gérondif et infinitif.
\item \textbf{Associer chaque notion à un signal.} \eng{since} $\rightarrow$ present perfect ; \eng{yesterday} $\rightarrow$ past simple ; \eng{by} + agent $\rightarrow$ passive ; \eng{if} + past $\rightarrow$ conditional type 2.
\item \textbf{Tester chaque règle par un exemple personnel.} Pour chaque point, écris une phrase sur ta vie : \eng{I have lived in Yaoundé since 2015.}
\item \textbf{Repérer tes erreurs récurrentes.} Relis tes copies : quelles fautes reviennent ? (3\textsuperscript{e} personne du singulier, \eng{informations}, prépositions...). Note-les sur une fiche.
\item \textbf{Travailler par transformations.} Prends une phrase simple et transforme-la : actif $\rightarrow$ passif, direct $\rightarrow$ indirect, affirmatif $\rightarrow$ négatif. C'est l'exercice-roi du Bac.
\item \textbf{Chronométrer.} Refais un exercice de grammaire en 10 minutes maximum pour simuler les conditions d'examen.
\end{enumerate}
\end{methode}

\begin{exoresolu}[Transformation de phrases — type Bac]
\textbf{Énoncé.} Rewrite each sentence as instructed without changing the meaning.
\begin{enumerate}
\item The journalist wrote the article yesterday. (Put into the passive voice.)
\item \eng{I will call you tomorrow,} the reporter told me. (Put into reported speech.)
\item The match was cancelled. It was raining heavily. (Join with \eng{because}.)
\item She speaks English fluently. (Turn into a question.)
\item I don't have a smartphone, so I can't follow the news online. (Rewrite beginning with \eng{If}.)
\end{enumerate}
\tcblower
\textbf{Solution détaillée.}
\begin{enumerate}
\item Le verbe \eng{wrote} est au past simple ; le passif se forme avec \eng{was/were} + participe passé. Le complément d'agent est introduit par \eng{by}.\\
\textbf{The article was written by the journalist yesterday.}
\item Discours rapporté au passé : \eng{will} recule à \eng{would}, \eng{tomorrow} devient \eng{the next day}.\\
\textbf{The reporter told me (that) he would call me the next day.}
\item La pluie est la cause de l'annulation ; \eng{because} introduit la subordonnée de cause.\\
\textbf{The match was cancelled because it was raining heavily.}
\item \eng{Speaks} est un present simple : la question exige l'auxiliaire \eng{does} et le verbe revient à la base verbale (plus de -s).\\
\textbf{Does she speak English fluently?}
\item Conditionnel de type 2 (irréel du présent) : \eng{If} + past simple, \eng{would} + base verbale.\\
\textbf{If I had a smartphone, I could follow the news online.} (ou \eng{I would be able to follow...})
\end{enumerate}
\textbf{Conclusion.} Chaque transformation obéit à un mécanisme précis : identifier le temps de départ, appliquer la règle de recul ou de structure, vérifier l'accord sujet-verbe. S'entraîner sur ces cinq types couvre l'essentiel de l'épreuve de grammaire.
\end{exoresolu}

\courssub{Savoir-faire 2 : S'entraîner à la compréhension écrite (reading)}

\begin{methode}[Réussir les questions de compréhension]
\begin{enumerate}
\item \textbf{Lire d'abord les questions}, puis le texte : tu sais ce que tu cherches.
\item \textbf{Repérer les mots-clés} de la question (noms propres, dates, verbes d'action) et les localiser dans le texte.
\item \textbf{Répondre avec des phrases complètes}, en reformulant : ne recopie jamais un long passage du texte.
\item \textbf{Respecter le temps de la question.} Question au past simple $\rightarrow$ réponse au past simple.
\item \textbf{Pour les questions de vocabulaire} (\eng{Find a word in the text which means...}), cherche un synonyme dans le paragraphe indiqué et vérifie la nature grammaticale (nom, verbe, adjectif).
\item \textbf{Pour \eng{True or False}}, justifie toujours par une courte citation du texte.
\end{enumerate}
\end{methode}

\begin{textebox}[Reading text: The rise of community radio in Cameroon]
In many towns and villages of Cameroon, community radio stations have become the main source of information for local people. Unlike national television, which broadcasts mostly in big cities, these small stations speak directly to their listeners in local languages as well as in English and French.

Journalists working for community radios often face serious difficulties. Equipment is expensive, electricity cuts are frequent, and reporters are rarely well paid. Nevertheless, they continue to inform the population about health campaigns, agricultural prices and local elections. During the last electoral campaign, for example, one station in the North-West region invited candidates to debate live on air. Thousands of listeners phoned in to ask questions.

According to media specialists, community radio plays a key role in democracy because it gives a voice to people who are usually ignored by the big national channels. However, they warn that journalists must be trained to check their facts, so that false information does not spread.
\end{textebox}

\begin{exoresolu}[Compréhension écrite — type Bac]
\textbf{Énoncé.} Read the text and answer the questions in complete sentences.
\begin{enumerate}
\item Why are community radio stations important for local people? (2 marks)
\item Mention two difficulties faced by community radio journalists. (2 marks)
\item What happened during the last electoral campaign in the North-West region? (2 marks)
\item Find in the text a word or expression which means: (a) \eng{to broadcast live} ; (b) \eng{verify}. (2 marks)
\item True or False? Justify: \eng{Community radio gives a voice to people ignored by national channels.} (2 marks)
\end{enumerate}
\tcblower
\textbf{Solution détaillée.}
\begin{enumerate}
\item La réponse se trouve au premier paragraphe (\eng{the main source of information}, \eng{local languages}).\\
\textbf{They are important because they are the main source of information for local people and they speak to listeners in local languages as well as in English and French.}
\item Deuxième paragraphe : choisir deux éléments parmi trois.\\
\textbf{They face expensive equipment and frequent electricity cuts.} (ou : \eng{reporters are rarely well paid}.)
\item Troisième phrase du deuxième paragraphe, au past simple.\\
\textbf{A radio station invited candidates to debate live on air, and thousands of listeners phoned in to ask questions.}
\item (a) \eng{live on air} = diffuser en direct ; (b) \eng{check} (their facts) = vérifier. Attention à la nature : on demande ici une expression et un verbe.\\
\textbf{(a) “live on air” ; (b) “to check”.}
\item La phrase du texte est presque identique : \eng{it gives a voice to people who are usually ignored by the big national channels}.\\
\textbf{True. The text says that community radio “gives a voice to people who are usually ignored by the big national channels”.}
\end{enumerate}
\textbf{Conclusion.} Une bonne réponse de compréhension est courte, complète, reformulée et au bon temps. La justification des \eng{True/False} doit citer le texte entre guillemets.
\end{exoresolu}

\courssub{Savoir-faire 3 : S'entraîner à la grammaire en contexte}

\begin{methode}[Réussir les exercices de grammaire (fill-in-the-gaps, multiple choice)]
\begin{enumerate}
\item \textbf{Lire toute la phrase} avant de choisir : le contexte donne le temps et le sens.
\item \textbf{Chercher les marqueurs temporels} : \eng{already, just, since, for} $\rightarrow$ present perfect ; \eng{ago, last week} $\rightarrow$ past simple ; \eng{while} + past continuous.
\item \textbf{Identifier la structure attendue} : après \eng{to} (infinitif), après une préposition (gérondif), après \eng{would rather} (base verbale).
\item \textbf{Vérifier l'accord sujet-verbe} : -s à la 3\textsuperscript{e} personne du singulier au present simple ; \eng{news, information} sont singuliers.
\item \textbf{Pour les modaux}, se demander ce qu'on exprime : obligation (\eng{must}), interdiction (\eng{mustn't}), absence d'obligation (\eng{needn't}), conseil (\eng{should}), capacité (\eng{can}), probabilité (\eng{may/might}).
\item \textbf{Relire la phrase complétée} à voix haute (mentalement) pour détecter une faute d'euphonie ou d'accord.
\end{enumerate}
\end{methode}

\begin{exoresolu}[Grammaire en contexte — fill in the gaps]
\textbf{Énoncé.} Complete the following passage with the correct form of the words in brackets.

Yesterday, while I \textbf{(1) (listen)} to the radio, the presenter announced that a new media centre \textbf{(2) (build)} in Buea next year. He said that journalists \textbf{(3) (must)} respect professional ethics if they wanted to keep the public's trust. I \textbf{(4) (never/hear)} such an interesting programme before. If all radio stations were as serious, people \textbf{(5) (be)} better informed.
\tcblower
\textbf{Solution détaillée.}
\begin{enumerate}
\item \eng{While} + action longue interrompue $\rightarrow$ past continuous : \textbf{was listening}.
\item \eng{Next year} + sujet qui subit l'action $\rightarrow$ futur passif (\eng{will be} + participe passé) : \textbf{will be built}.
\item Discours rapporté après \eng{said that} : le modal \eng{must} recule normalement à \eng{had to}. Réponse attendue : \textbf{had to} ; \eng{must} reste accepté si l'obligation est présentée comme toujours vraie.
\item \eng{Before} avec un repère passé (\eng{yesterday}) et \eng{never} $\rightarrow$ past perfect : \textbf{had never heard}.
\item Conditionnel type 2 : \eng{If} + past simple (\eng{were}), principale en \eng{would} + base verbale : \textbf{would be}.
\end{enumerate}
\textbf{Passage complété :} \eng{Yesterday, while I was listening to the radio, the presenter announced that a new media centre will be built in Buea next year. He said that journalists had to respect professional ethics if they wanted to keep the public's trust. I had never heard such an interesting programme before. If all radio stations were as serious, people would be better informed.}\\
\textbf{Conclusion.} Chaque trou se résout par un signal : marqueur temporel, structure passive, recul des temps, conditionnel. Surligne ces signaux avant de répondre.
\end{exoresolu}

\courssub{Savoir-faire 4 : S'entraîner au vocabulaire des médias}

\begin{methode}[Maîtriser le lexique thématique le jour de l'examen]
\begin{enumerate}
\item \textbf{Classer le vocabulaire par familles} : la presse (\eng{headline, article, editor}), la radio/TV (\eng{broadcast, channel, presenter, live}), Internet (\eng{post, share, fake news, social media}), la déontologie (\eng{source, censorship, fact-checking}).
\item \textbf{Apprendre les collocations}, pas les mots isolés : \eng{to broadcast news}, \eng{to check facts}, \eng{to spread rumours}, \eng{breaking news}.
\item \textbf{Attention aux faux amis} : \eng{actually} = en fait (pas « actuellement » = \eng{currently}) ; \eng{a report} = un reportage/un rapport ; \eng{to assist} = aider (pas « assister à » = \eng{attend}).
\item \textbf{Pour les questions de vocabulaire au Bac}, vérifie toujours la nature du mot demandé : si on cherche un adjectif, ne propose pas un nom.
\item \textbf{Réutiliser le lexique dans la composition} : un texte sur les médias gagne des points avec 4 ou 5 mots précis du champ lexical.
\end{enumerate}
\end{methode}

\begin{exoresolu}[Vocabulary — media and communication]
\textbf{Énoncé.} Choose the correct word to complete each sentence: \eng{headline / broadcast / source / fake news / censorship / presenter}.
\begin{enumerate}
\item The \dots\ of the evening news speaks five languages.
\item The radio station will \dots\ the football match live from Douala.
\item A journalist must never reveal a confidential \dots.
\item Many people believed the story, but it was only \dots.
\item The front-page \dots\ announced the election results in big letters.
\end{enumerate}
\tcblower
\textbf{Solution détaillée.}
\begin{enumerate}
\item On cherche une personne qui présente un journal télévisé : \textbf{presenter} (« pré-zan-teur »).
\item Verbe signifiant « diffuser » ; suivi de \eng{live} : \textbf{broadcast} (verbe irrégulier : broadcast-broadcast-broadcast).
\item Un journaliste protège ses informateurs : \textbf{source}.
\item Une fausse information qui circule : \textbf{fake news} (\eng{news} est indénombrable : \eng{it was fake news}, sans article ni pluriel).
\item Le gros titre de une d'un journal : \textbf{headline}.
\end{enumerate}
\textbf{Conclusion.} Pour choisir le bon mot, demande-toi : est-ce une personne, un verbe, une chose ? La nature grammaticale et la collocation (\eng{to broadcast live}, \eng{front-page headline}) lèvent presque toujours le doute.
\end{exoresolu}

\courssub{Savoir-faire 5 : S'entraîner à la composition (writing)}

\begin{methode}[Rédiger une composition type Bac en 5 étapes]
\begin{enumerate}
\item \textbf{Analyser le sujet} (2 min) : type de texte demandé (lettre, article, essai), nombre de mots, destinataire, registre (formel/informel).
\item \textbf{Faire un plan rapide} (3 min) : introduction (1 idée), développement (2 ou 3 paragraphes, une idée + un exemple chacun), conclusion (1 idée).
\item \textbf{Rédiger} (15 min) : phrases simples et correctes plutôt que phrases longues et fausses. Utilise des connecteurs : \eng{firstly, moreover, however, in addition, to conclude}.
\item \textbf{Mobiliser la grammaire révisée} : un passif, un conditionnel, un discours rapporté, un present perfect — cela montre ta maîtrise.
\item \textbf{Relire} (5 min) : -s de la 3\textsuperscript{e} personne, temps homogènes, orthographe (\eng{information} sans s, \eng{advice} sans s), ponctuation, nombre de mots respecté.
\end{enumerate}
\end{methode}

\begin{exoresolu}[Composition — sujet type Bac]
\textbf{Énoncé.} \eng{Write a composition of about 150 words on the following topic: “Social media does more harm than good to young people in Cameroon.” Do you agree? Give your opinion with examples.}
\tcblower
\textbf{Démarche.} Sujet argumentatif : donner son avis avec des arguments et des exemples. Plan : introduction (position), argument 1 + exemple, argument 2 + exemple, concession (\eng{however}), conclusion.

\textbf{Réponse rédigée (modèle).}

\eng{In my opinion, social media does more harm than good to young people in Cameroon, although it has some advantages.}

\eng{Firstly, many students spend hours on their phones instead of studying. Since Facebook and WhatsApp became popular, school results have dropped in many families, because young people chat late at night and are tired in class.}

\eng{Secondly, social media spreads fake news very quickly. During elections, for example, false information is shared thousands of times before it can be checked, and this creates fear and conflict.}

\eng{However, it must be admitted that social media can also be useful: it helps young people to find information, to learn English and to stay in contact with their families.}

\eng{To conclude, if young Cameroonians were taught to use social media wisely, the balance would be more positive. But for the moment, I believe the harm is greater than the good.}

(Environ 150 mots.)

\textbf{Commentaire grammatical.} Le modèle contient volontairement : un present perfect (\eng{have dropped}), un passif (\eng{is shared}, \eng{can be checked}), un conditionnel type 2 (\eng{if... were taught... would be}) et des connecteurs variés. C'est exactement ce que le correcteur récompense.
\end{exoresolu}

\begin{attention}[Les 5 erreurs qui coûtent des points au Bac]
\begin{enumerate}
\item \textbf{Oublier le -s de la 3\textsuperscript{e} personne du singulier.} \eng{He watch TV} $\times$ $\rightarrow$ \eng{He watches TV every evening.} C'est la faute la plus sanctionnée au present simple.
\item \textbf{Mettre un pluriel aux noms indénombrables.} \eng{informations, advices, news are} $\times$ $\rightarrow$ \eng{The information is useful. She gave me some advice. The news is good.}
\item \textbf{Calquer le français.} \eng{I have 16 years} $\times$ $\rightarrow$ \eng{I am sixteen years old.} ; \eng{actually} ne veut pas dire « actuellement » (\eng{currently}) ; \eng{to assist} ne veut pas dire « assister à » (\eng{to attend}).
\item \textbf{Mélanger les temps.} Avec \eng{since/for} et \eng{already/just}, le present perfect est obligatoire : \eng{I live here since 2020} $\times$ $\rightarrow$ \eng{I have lived here since 2020.} Au discours rapporté, n'oublie pas le recul : \eng{He said he will come} $\times$ $\rightarrow$ \eng{He said he would come.}
\item \textbf{Se tromper de préposition et d'ordre des mots.} \eng{to listen music} $\times$ $\rightarrow$ \eng{to listen to music} ; \eng{to depend of} $\times$ $\rightarrow$ \eng{to depend on} ; les adverbes de fréquence se placent avant le verbe : \eng{I watch often TV} $\times$ $\rightarrow$ \eng{I often watch TV.}
\end{enumerate}
\end{attention}

\newpage

\coursec{Exercices}

\courssub{Je vérifie mes connaissances}

\begin{exercice}[QCM : le bon choix]
Pour chaque phrase, entoure la bonne réponse.

\begin{enumerate}
\item The journalist said that he --- the story the day before.
\begin{itemize}
\item[a)] investigates \quad b)] investigated \quad c)] had investigated \quad d)] will investigate
\end{itemize}
\item The news --- every hour on this radio station.
\begin{itemize}
\item[a)] broadcast \quad b)] is broadcast \quad c)] broadcasts \quad d)] is broadcasting
\end{itemize}
\item You --- buy this newspaper; it is free online.
\begin{itemize}
\item[a)] mustn't \quad b)] don't have to \quad c)] can't \quad d)] shouldn't
\end{itemize}
\item If the reporter had checked his sources, he --- the mistake.
\begin{itemize}
\item[a)] would avoid \quad b)] will avoid \quad c)] would have avoided \quad d)] avoided
\end{itemize}
\item I have been working for this TV channel --- five years.
\begin{itemize}
\item[a)] since \quad b)] during \quad c)] for \quad d)] ago
\end{itemize}
\end{enumerate}
\end{exercice}

\begin{exercice}[Vrai ou faux — justifie ta réponse]
Dis si chaque phrase est correcte ou incorrecte. Si elle est incorrecte, corrige-la.

\begin{enumerate}
\item \eng{I am agree with your opinion about social media.}
\item \eng{She gave me many informations about the press conference.}
\item \eng{He has been living in Douala since 2020.}
\item \eng{The interview was being recorded when we arrived.}
\item \eng{The editor asked me if I had checked my sources.}
\end{enumerate}
\end{exercice}

\begin{exercice}[Vocabulaire : attention aux faux amis]
Complète chaque phrase avec le mot correct entre parenthèses.

\begin{enumerate}
\item The minister is --- in Buea for a media forum. (\emph{actually / currently})
\item Journalists must be --- to criticism from the public. (\emph{sensible / sensitive})
\item You can buy foreign magazines in this --- near the post office. (\emph{library / bookshop})
\item The reporter asked a very --- question about corruption. (\emph{delicate / delicious})
\item The channel will --- the football match live on Saturday. (\emph{assist / broadcast})
\end{enumerate}
\end{exercice}

\begin{exercice}[Conjugaison : mets les verbes au temps correct]
Mets les verbes entre parenthèses au temps qui convient.

\begin{enumerate}
\item When we arrived, the press conference (already / start).
\item The minister (speak) for an hour and the journalists (take) notes.
\item Every evening, my father (watch) the news on CRTV.
\item Look! The reporter (interview) a witness of the accident.
\item By next year, she (work) for this radio station for ten years.
\end{enumerate}
\end{exercice}

\courssub{Je m'entraîne}

\begin{exercice}[Traduction]
Traduis les phrases suivantes en anglais.

\begin{enumerate}
\item On m'a dit que le journal était publié chaque semaine.
\item Si j'avais le temps, je regarderais le journal télévisé tous les soirs.
\item Il faut vérifier ses sources avant de publier un article.
\item Le présentateur a demandé aux téléspectateurs de ne pas croire toutes les rumeurs.
\item Depuis quand travailles-tu pour cette chaîne de télévision ?
\end{enumerate}
\end{exercice}

\begin{exercice}[Transformations : voix passive et discours rapporté]
\textbf{A.} Mets les phrases suivantes à la voix passive.

\begin{enumerate}
\item \eng{The BBC broadcasts the news every hour.}
\item \eng{Someone has stolen the reporter's camera.}
\item \eng{The editor will publish the article tomorrow.}
\end{enumerate}

\textbf{B.} Rapporte les paroles au discours indirect.

\begin{enumerate}
\item \eng{“I will investigate this story,” the journalist said.}
\item \eng{“Did you check your sources?” the editor asked me.}
\item \eng{“Don't publish the photo,” the lawyer told the newspaper.}
\end{enumerate}
\end{exercice}

\begin{exercice}[Texte à trous : modaux et conditionnels]
Complète le texte avec les mots de la liste suivante (chaque mot ne sert qu'une fois) : \eng{must — don't have to — might — should — would — had checked}.

\begin{textebox}[Responsible journalism]
Young people today read most of their news online. You \ldots\ buy a newspaper any more — everything is free on the internet. But you \ldots\ always verify what you read, because anyone can publish false information. If readers shared every story without thinking, rumours \ldots\ spread very fast. If the reporter \ldots\ the facts, the scandal would never have happened. My advice? You \ldots\ compare at least two sources before believing a shocking headline, and remember: a serious journalist \ldots\ never invent quotations.
\end{textebox}
\end{exercice}

\begin{exercice}[Appariements : le vocabulaire des médias]
Associe chaque mot anglais (1--8) à sa définition (a--h).

\begin{center}
\begin{tabularx}{\linewidth}{|X|X|}
\hline
\rowcolor{popblueL} \textbf{Mot} & \textbf{Définition} \\
\hline
1. headline & a) a person who presents the news on television \\
\hline
2. news anchor & b) the title of an article, printed in large letters \\
\hline
3. source & c) a report sent by a journalist from another country \\
\hline
4. censorship & d) a person or document that provides information \\
\hline
5. dispatch & e) the control of what may be published or broadcast \\
\hline
6. advertisement & f) a false story spread deliberately to deceive people \\
\hline
7. fake news & g) a paid message that promotes a product \\
\hline
8. editor-in-chief & h) the person who directs a newspaper \\
\hline
\end{tabularx}
\end{center}
\end{exercice}

\courssub{Je me prépare au Bac}

\begin{exercice}[Compréhension écrite et langue — type Bac]
Lis le texte suivant, puis réponds aux questions en anglais, avec tes propres mots.

\begin{textebox}[The changing face of news in Cameroon]
Not long ago, most Cameroonians waited for the evening news on television or bought a newspaper at a kiosk in the morning. Today, the situation has changed dramatically. In Yaoundé and Douala, young people now receive information on their phones within seconds: a message in a WhatsApp group, a video on Facebook, a short post on X. News has never travelled so fast.

This speed, however, has a price. False stories circulate widely, and many readers share them without checking. “If people verified their sources, half of these rumours would disappear,” says a journalism teacher in Buea. Some traditional newspapers have closed because readers no longer buy printed copies, and advertising money has moved to the internet. Yet others have adapted: they publish online editions, produce podcasts and send news alerts by SMS.

Journalists themselves are adapting too. A good reporter must now write articles, record videos and answer readers' comments — sometimes all in the same day. “The tools have changed,” explains one editor, “but the mission has not: to inform the public honestly.” Whether news is printed on paper or displayed on a screen, the public still needs journalists who check their facts.
\end{textebox}

\textbf{A. Comprehension questions.} Answer in complete sentences.

\begin{enumerate}
\item How did most Cameroonians get the news in the past? (2 marks)
\item Give two ways in which young people receive information today. (2 marks)
\item According to the journalism teacher, what would happen if people verified their sources? (2 marks)
\item Why have some traditional newspapers closed? (2 marks)
\item What has not changed, according to the editor? (2 marks)
\end{enumerate}

\textbf{B. Language work.}

\begin{enumerate}
\item Find in the text the passive form used in the last paragraph and rewrite the sentence in the active voice. (2 marks)
\item Report the teacher's words: \eng{“If people verified their sources, half of these rumours would disappear.”} (2 marks)
\item Put the verbs in brackets in the correct tense: \eng{News never (travel) so fast before. Some newspapers (close) because readers no longer (buy) printed copies.} (3 marks)
\item Find in the text a word that means: (a) a small shop that sells newspapers; (b) untrue. (2 marks)
\end{enumerate}
\end{exercice}

\begin{exercice}[Traduction — type Bac]
Traduis le passage suivant en anglais.

\begin{textebox}[Texte à traduire]
Les médias jouent un rôle très important dans notre société. Chaque jour, des millions de Camerounais écoutent la radio ou regardent la télévision pour s'informer. Malheureusement, certaines informations publiées sur les réseaux sociaux sont fausses. Si les journalistes vérifiaient toujours leurs sources, le public aurait plus confiance en eux. Il faut donc apprendre aux jeunes à distinguer les vraies nouvelles des rumeurs.
\end{textebox}
\end{exercice}

\begin{exercice}[Composition — type Bac]
Write a composition of about \textbf{150 words} on the following topic:

\begin{center}
\emph{“Traditional newspapers are dying because of the internet. Do you agree?”}
\end{center}

Your composition must follow this plan:
\begin{enumerate}
\item \textbf{Introduction}: present the topic and the two possible points of view.
\item \textbf{Development}: give two arguments, with examples (you may refer to the situation in Cameroon).
\item \textbf{Conclusion}: give your personal opinion and justify it.
\end{enumerate}

Pay attention to grammar (tenses, modals, connectors such as \eng{however}, \eng{moreover}, \eng{in my opinion}) and to spelling. (20 marks)
\end{exercice}

\newpage

\coursec{Corrigés des exercices}

\begin{corrige}[Exercice 1 — QCM : le bon choix]

\begin{enumerate}
\item \textbf{c) had investigated.} Discours rapporté : le verbe au prétérit (\eng{investigated}) recule au plus-que-parfait (\eng{had investigated}) après \eng{said that}. L'expression \eng{the day before} confirme le recul temporel.
\item \textbf{b) is broadcast.} Voix passive au présent simple : \eng{the news} subit l'action. Attention : \eng{news} est singulier et le participe passé de \eng{broadcast} est \eng{broadcast} (invariable).
\item \textbf{b) don't have to.} Sens : « ce n'est pas nécessaire » (le journal est gratuit en ligne). \eng{Mustn't} signifierait « interdiction », ce qui est contraire au sens.
\item \textbf{c) would have avoided.} Conditionnel de type 3 : \eng{if + plus-que-parfait, would have + participe passé} (regret sur le passé).
\item \textbf{c) for.} \eng{For + durée} (\eng{for five years}) ; \eng{since} s'emploie avec un point de départ (\eng{since 2019}).
\end{enumerate}

\textbf{Barème indicatif :} 1 point par bonne réponse (5 points).
\end{corrige}

\begin{corrige}[Exercice 2 — Vrai ou faux]

\begin{enumerate}
\item \textbf{Incorrecte.} \eng{Agree} est un verbe : on ne peut pas le faire précéder de \eng{be}. Correction : \eng{I agree with your opinion about social media.}
\item \textbf{Incorrecte.} \eng{Information} est indénombrable en anglais : jamais de \eng{-s} ni de \eng{many}. Correction : \eng{She gave me a lot of information about the press conference.} (ou \eng{much information}).
\item \textbf{Correcte.} Present perfect continu + \eng{since} + point de départ : \eng{He has been living in Douala since 2020.}
\item \textbf{Correcte.} Voix passive au prétérit continu : \eng{was being recorded} — l'action était en cours quand nous sommes arrivés.
\item \textbf{Correcte.} Discours rapporté : question fermée introduite par \eng{if}, avec recul du prétérit au plus-que-parfait (\eng{had checked}).
\end{enumerate}

\textbf{Barème indicatif :} 1 point par réponse justifiée (5 points).
\end{corrige}

\begin{corrige}[Exercice 3 — Vocabulaire : les faux amis]

\begin{enumerate}
\item \eng{currently} — « actuellement ». \eng{Actually} signifie « en fait, en réalité ».
\item \eng{sensitive} — « sensible ». \eng{Sensible} signifie « raisonnable, sensé ».
\item \eng{bookshop} — « librairie » (commerce). Une \eng{library} est une bibliothèque (on y emprunte, on n'y achète pas).
\item \eng{delicate} — « délicate » (difficile, qui demande du tact). \eng{Delicious} signifie « délicieux » (nourriture).
\item \eng{broadcast} — « diffuser ». \eng{Assist} signifie « aider » ; « assister à » se dit \eng{attend} ou \eng{watch}.
\end{enumerate}

\textbf{Points de vigilance :} ces cinq faux amis reviennent très souvent au Bac. Apprends-les par paires de confusion : \eng{actually} (\eng{in fact}) \emph{vs} \eng{currently} (\eng{now}) ; \eng{sensible} (\eng{reasonable}) \emph{vs} \eng{sensitive} ; \eng{library} (\eng{bibliothèque}) \emph{vs} \eng{bookshop} (\eng{librairie}) ; \eng{assist} (\eng{help}) \emph{vs} \eng{attend} (\eng{assister à}) ; \eng{delicious} (\eng{délicieux}) \emph{vs} \eng{delicate} (\eng{délicat}).

\textbf{Barème indicatif :} 1 point par bonne réponse (5 points).
\end{corrige}

\begin{corrige}[Exercice 4 — Conjugaison]

\begin{enumerate}
\item \eng{When we arrived, the press conference \textbf{had already started}.} Plus-que-parfait : action antérieure à un moment passé (\eng{already} est un indice fort).
\item \eng{The minister \textbf{was speaking} for an hour and the journalists \textbf{were taking} notes.} Prétérit continu : deux actions longues simultanées dans le passé.
\item \eng{Every evening, my father \textbf{watches} the news on CRTV.} Présent simple : habitude (\eng{every evening}). Attention à la 3\textsuperscript{e} personne : \eng{watches}.
\item \eng{Look! The reporter \textbf{is interviewing} a witness of the accident.} Présent continu : action en cours au moment où l'on parle (\eng{Look!}).
\item \eng{By next year, she \textbf{will have been working} for this radio station for ten years.} Futur antérieur continu : bilan d'une durée à une date future (\eng{by next year} + \eng{for ten years}). La forme simple \eng{will have worked} est aussi acceptée.
\end{enumerate}

\textbf{Barème indicatif :} 1 point par verbe correctement conjugué (6 points).
\end{corrige}

\begin{corrige}[Exercice 5 — Traduction]

\begin{enumerate}
\item \eng{I was told that the newspaper was published every week.} Voix passive (\eng{I was told}) + recul du présent au prétérit dans le discours rapporté.
\item \eng{If I had time, I would watch the TV news every evening.} Conditionnel de type 2 : \eng{if + prétérit, would + base verbale}. Ne jamais mettre \eng{would} dans la subordonnée.
\item \eng{You must check your sources before publishing an article.} \eng{Il faut} = \eng{must} ; \eng{before} + gérondif (\eng{publishing}).
\item \eng{The news anchor asked the viewers not to believe all the rumours.} \eng{Ask somebody not to do something} ; \eng{présentateur} = \eng{news anchor} ou \eng{newsreader}.
\item \eng{How long have you been working for this TV channel?} \eng{Depuis quand} + durée = \eng{how long} avec present perfect continu. Ne jamais dire \eng{since when do you work}.
\end{enumerate}

\textbf{Points de vigilance :} méfie-toi du présent français après « depuis » (l'anglais exige le present perfect) et de « il faut » (jamais \eng{it must}).

\textbf{Barème indicatif :} 2 points par phrase (10 points) : 1 point pour la structure, 1 point pour le temps et le vocabulaire.
\end{corrige}

\begin{corrige}[Exercice 6 — Transformations]

\textbf{A. Voix passive}

\begin{enumerate}
\item \eng{The news is broadcast by the BBC every hour.} Présent simple passif : \eng{is + broadcast}.
\item \eng{The reporter's camera has been stolen.} Present perfect passif : \eng{has been + stolen}. Le complément d'agent \eng{someone} disparaît (agent inconnu ou sans importance).
\item \eng{The article will be published by the editor tomorrow.} Futur passif : \eng{will be + published}.
\end{enumerate}

\textbf{B. Discours rapporté}

\begin{enumerate}
\item \eng{The journalist said (that) he would investigate that story.} \eng{Will} devient \eng{would} ; \eng{this} devient \eng{that}.
\item \eng{The editor asked me if I had checked my sources.} Question fermée : \eng{if} + ordre sujet--verbe ; prétérit $\rightarrow$ plus-que-parfait.
\item \eng{The lawyer told the newspaper not to publish the photo.} Impératif négatif rapporté : \eng{told + complément + not to + base verbale}.
\end{enumerate}

\textbf{Points de vigilance :} au discours rapporté, on ne garde ni guillemets ni point d'interrogation, et l'ordre des mots redevient celui d'une phrase affirmative.

\textbf{Barème indicatif :} 1,5 point par transformation (9 points).
\end{corrige}

\begin{corrige}[Exercice 7 — Texte à trous : modaux et conditionnels]

\textbf{Texte corrigé :}

\par\textbf{Responsible journalism — corrigé}\par
\eng{Young people today read most of their news online. You \textbf{don't have to} buy a newspaper any more — everything is free on the internet. But you \textbf{must} always verify what you read, because anyone can publish false information. If readers shared every story without thinking, rumours \textbf{would} spread very fast. If the reporter \textbf{had checked} the facts, the scandal would never have happened. My advice? You \textbf{should} compare at least two sources before believing a shocking headline, and remember: a serious journalist \textbf{must} never invent quotations.}\par

\textbf{Justifications :}
\begin{itemize}
\item \eng{don't have to} : absence de nécessité (tout est gratuit).
\item \eng{must} (2\textsuperscript{e} trou) : forte nécessité, obligation morale de vérifier.
\item \eng{would} : conditionnel de type 2 (\eng{if + prétérit, would + base verbale}).
\item \eng{had checked} : plus-que-parfait dans un conditionnel de type 3 (regret sur le passé).
\item \eng{should} : conseil (\eng{my advice}).
\item \eng{must} (dernier trou) : interdiction forte formulée positivement — \eng{must never} = « ne doit jamais ».
\end{itemize}

\textbf{Barème indicatif :} 1 point par trou (6 points).
\end{corrige}

\begin{corrige}[Exercice 8 — Appariements : le vocabulaire des médias]

\begin{center}
\begin{tabularx}{\linewidth}{|X|X|}
\hline
\rowcolor{popblueL} \textbf{Mot} & \textbf{Réponse} \\
\hline
1. headline & b) the title of an article, printed in large letters \\
\hline
2. news anchor & a) a person who presents the news on television \\
\hline
3. source & d) a person or document that provides information \\
\hline
4. censorship & e) the control of what may be published or broadcast \\
\hline
5. dispatch & c) a report sent by a journalist from another country \\
\hline
6. advertisement & g) a paid message that promotes a product \\
\hline
7. fake news & f) a false story spread deliberately to deceive people \\
\hline
8. editor-in-chief & h) the person who directs a newspaper \\
\hline
\end{tabularx}
\end{center}

\textbf{Prononciation utile :} \eng{headline} « hède-laïne », \eng{anchor} « an-keur », \eng{censorship} « sen-seur-chip », \eng{dispatch} « dis-patche ».

\textbf{Barème indicatif :} 0,5 point par association correcte (4 points).
\end{corrige}

\begin{corrige}[Exercice 9 — Compréhension écrite et langue — type Bac]

\textbf{A. Comprehension questions — réponses modèles}

\begin{enumerate}
\item \eng{In the past, most Cameroonians watched the evening news on television or bought a newspaper at a kiosk in the morning.} (2 points)
\item \eng{Today, young people receive information on their phones, for example through messages in WhatsApp groups and videos on Facebook.} (2 points)
\item \eng{According to the teacher, if people verified their sources, half of the rumours would disappear.} (2 points)
\item \eng{Some traditional newspapers have closed because readers no longer buy printed copies and advertising money has moved to the internet.} (2 points)
\item \eng{According to the editor, the mission of journalists has not changed: it is still to inform the public honestly.} (2 points)
\end{enumerate}

\textbf{B. Language work}

\begin{enumerate}
\item La forme passive du dernier paragraphe : \eng{Whether news \textbf{is printed} on paper or \textbf{displayed} on a screen...} Voix active : \eng{Whether we print news on paper or display it on a screen...} (2 points)
\item \eng{The journalism teacher said that if people verified their sources, half of those rumours would disappear.} Le conditionnel de type 2 ne recule pas : \eng{verified} et \eng{would} restent inchangés. (2 points)
\item \eng{News \textbf{has never travelled} so fast before. Some newspapers \textbf{have closed} because readers no longer \textbf{buy} printed copies.} Present perfect (bilan au présent) puis présent simple (situation actuelle). (3 points)
\item (a) \eng{kiosk} ; (b) \eng{false}. (2 points)
\end{enumerate}

\textbf{Points de vigilance :} réponds toujours par des phrases complètes, avec tes propres mots ; ne recopie jamais un long passage du texte. Aux questions de vocabulaire, donne le mot exactement comme il apparaît dans le texte.

\textbf{Barème indicatif :} 19 points au total (A : 10 ; B : 9).
\end{corrige}

\begin{corrige}[Exercice 10 — Traduction — type Bac]

\textbf{Proposition de traduction :}

\par\textbf{Model translation}\par
\eng{The media play a very important role in our society. Every day, millions of Cameroonians listen to the radio or watch television to get information. Unfortunately, some of the information published on social media is false. If journalists always checked their sources, the public would trust them more. We must therefore teach young people to distinguish real news from rumours.}\par

\textbf{Points de vigilance :}
\begin{itemize}
\item \eng{The media play} : \eng{media} est pluriel en anglais soutenu (singulier : \eng{medium}).
\item \eng{to get information} ou \eng{to inform themselves} : \eng{information} reste indénombrable.
\item \eng{If journalists always checked..., the public would trust them more} : conditionnel de type 2 — pas de \eng{would} après \eng{if}.
\item \eng{to distinguish real news from rumours} : \eng{distinguish A from B} ; \eng{news} est singulier et indénombrable.
\item « Il faut » se rend ici par \eng{we must} ou \eng{it is necessary to}, jamais par \eng{it must}.
\end{itemize}

\textbf{Barème indicatif :} 10 points — 2 points par phrase (structure 1 point, temps et lexique 1 point).
\end{corrige}

\begin{corrige}[Exercice 11 — Composition — type Bac]

\textbf{Proposition de rédaction modèle (environ 150 mots) :}

\par\textbf{Model composition}\par
\eng{Are traditional newspapers really dying because of the internet? Some people believe that printed newspapers will soon disappear, while others think they still have an important role to play.}\par
\eng{On the one hand, it is true that the internet has changed our habits. In Yaoundé and Douala, most young people now read the news on their phones for free, so fewer people buy newspapers at kiosks. Moreover, advertising money has moved online, and several newspapers have closed. On the other hand, printed newspapers still offer serious, verified information, unlike many stories shared on social media. In rural areas of Cameroon, where the internet connection is weak, the radio and the newspaper remain essential.}\par
\eng{In my opinion, traditional newspapers are not dying; they are transforming themselves. If they publish online editions and keep checking their facts, they will survive, because the public will always need honest journalism.}\par

\textbf{Points de vigilance :}
\begin{itemize}
\item Respecte le plan en trois parties et la limite de mots (environ 150).
\item Emploie des connecteurs variés : \eng{on the one hand / on the other hand}, \eng{moreover}, \eng{however}, \eng{in my opinion}.
\item Vérifie les accords (3\textsuperscript{e} personne du singulier : \eng{the internet has changed}), les indénombrables (\eng{information}, \eng{news}) et les conditionnels (\eng{if they publish..., they will survive}).
\item Relis-toi : orthographe de \eng{believe}, \eng{necessary}, \eng{government}, \eng{which}.
\end{itemize}

\textbf{Barème indicatif (20 points) :} respect du sujet et du plan : 4 ; richesse du vocabulaire : 4 ; correction grammaticale : 6 ; cohérence et connecteurs : 3 ; orthographe et présentation : 3.
\end{corrige}

\newpage

\coursec{Fiche bilan}

\begin{popfigure}
\begin{center}\tcbox[colback=popdark,colframe=popdark,coltext=white,arc=9pt,boxsep=4pt,left=6pt,right=6pt]{\bfseries\small General Revision}\end{center}
\vspace{-2pt}
\begin{tcbraster}[raster columns=3,raster equal height=rows,raster column skip=6pt,raster row skip=6pt,enhanced,blankest]
\begin{tcolorbox}[enhanced,colback=popblue,colframe=popblue,coltext=white,boxrule=0.9pt,arc=3pt,left=4pt,right=4pt,top=3pt,bottom=3pt,boxsep=1pt,halign=left]\bfseries\scriptsize Verb Tenses Map\end{tcolorbox}
\begin{tcolorbox}[enhanced,colback=poporange,colframe=poporange,coltext=white,boxrule=0.9pt,arc=3pt,left=4pt,right=4pt,top=3pt,bottom=3pt,boxsep=1pt,halign=left]\bfseries\scriptsize Modals \& Semi-auxiliaries\end{tcolorbox}
\begin{tcolorbox}[enhanced,colback=popgreen,colframe=popgreen,coltext=white,boxrule=0.9pt,arc=3pt,left=4pt,right=4pt,top=3pt,bottom=3pt,boxsep=1pt,halign=left]\bfseries\scriptsize Passive Voice \& Reported Speech\end{tcolorbox}
\begin{tcolorbox}[enhanced,colback=poppurple,colframe=poppurple,coltext=white,boxrule=0.9pt,arc=3pt,left=4pt,right=4pt,top=3pt,bottom=3pt,boxsep=1pt,halign=left]\bfseries\scriptsize Conditionals, Wishes \& Complex Sent.\end{tcolorbox}
\begin{tcolorbox}[enhanced,colback=poppink,colframe=poppink,coltext=white,boxrule=0.9pt,arc=3pt,left=4pt,right=4pt,top=3pt,bottom=3pt,boxsep=1pt,halign=left]\bfseries\scriptsize Francophone Pitfalls\end{tcolorbox}
\begin{tcolorbox}[enhanced,colback=popgold,colframe=popgold,coltext=white,boxrule=0.9pt,arc=3pt,left=4pt,right=4pt,top=3pt,bottom=3pt,boxsep=1pt,halign=left]\bfseries\scriptsize Bac Strategies \& Practice\end{tcolorbox}
\end{tcbraster}

\legende{Carte mentale de la révision générale — Leçon 53}
\end{popfigure}

\begin{bilan}

\textbf{1. Lexique clé thématique (Media \& Communication + General)}
\begin{center}
\begin{tabularx}{\linewidth}{|X|X|X|}
\hline
\rowcolor{popblueL} \textbf{English} & \textbf{Français} & \textbf{Remarque} \\
\hline
broadcast / to broadcast & diffusion / diffuser & \emph{irrégulier : broadcast - broadcast - broadcast} \\
coverage & couverture (médiatique) & \eng{news coverage} \\
headline & gros titre, manchette & \eng{to make headlines} \\
outlet & organe de presse, média & \eng{media outlets} \\
reliable / unreliable & fiable / peu fiable & \emph{faux ami : \textbf{reliable} $\neq$ relais} \\
bias / biased & parti pris / partial & \eng{a biased report} \\
to spread (spread, spread) & propager, diffuser & \eng{spread fake news} \\
audience & public, auditoire & \emph{pas "audience" = hearing (juridique)} \\
deadline & date limite, échéance & \eng{meet a deadline} \\
freelancer & pigiste, travailleur indépendant & \\
\hline
\end{tabularx}
\end{center}

\textbf{2. Règles de grammaire essentielles (à connaître par cœur)}
\begin{center}
\begin{tabularx}{\linewidth}{|X|X|X|}
\hline
\rowcolor{popblueL} \textbf{Point} & \textbf{Forme / Structure} & \textbf{Emploi clé / Piège} \\
\hline
\textbf{Temps du récit} & \eng{Past Simple} vs \eng{Past Perfect} & \eng{Past Perfect} = antériorité (flash-back). \textbf{Pas} \eng{Past Perfect} si chronologie claire avec \eng{then/after}. \\
\textbf{Present Perfect} & \eng{have/has + V-en} & Lien présent : résultat, expérience, durée (\eng{since/for}). \textbf{Pas} avec date passée (\eng{yesterday}). \\
\textbf{Futurs} & \eng{will} (prédiction, décision instantanée) vs \eng{be going to} (intention, preuve) vs \eng{Pres. Cont.} (plan fixé) & \eng{Shall I/we?} = proposition. \\
\textbf{Modaux passé} & \eng{modal + have + V-en} & \eng{must have done} (certitude), \eng{can't have done} (impossibilité), \eng{should have done} (regret/conseil non suivi). \\
\textbf{Voix passive} & \eng{be + V-en} (tous temps) & Agent (\eng{by}) souvent omis. \eng{get} passif = oral, souvent négatif (\eng{get fired}). \\
\textbf{Discours rapporté} & \eng{Backshifting} des temps + pronoms + marqueurs spatio-temporels & \textbf{Exception} : vérité générale, fait encore vrai, \eng{must} $\to$ \eng{must/had to}. \\
\textbf{Conditionnel 1} & \eng{If + Pres. Simple, ... will + BV} & Réel, possible. \\
\textbf{Conditionnel 2} & \eng{If + Past Simple, ... would + BV} & Irréel présent/futur, hypothétique. \eng{If I \textbf{were} you}. \\
\textbf{Conditionnel 3} & \eng{If + Past Perfect, ... would have + V-en} & Irréel passé, regret, reproche. \\
\textbf{Wish / If only} & \eng{Wish + Past Simple} (présent) / \eng{Past Perfect} (passé) / \eng{would} (envie de changement) & \eng{I wish I \textbf{knew}} (je ne sais pas). \eng{I wish I \textbf{had known}}. \\
\textbf{Relative clauses} & \eng{who/which/that/whose/where/when} & \textbf{Defining} : pas de virgule, \eng{that} possible. \textbf{Non-defining} : virgules, \eng{that} interdit. \\
\hline
\end{tabularx}
\end{center}

\textbf{3. Méthodes express pour l'examen}
\begin{itemize}
    \item \textbf{Compréhension (Reading/Listening) :} Lire les questions $\to$ Survoler (skimming) $\to$ Chercher l'info (scanning) $\to$ Répondre en phrases complètes $\to$ Vérifier l'orthographe des mots du texte.
    \item \textbf{Grammaire (Transformation/Remise en forme) :} Identifier le temps requis (mots-clés : \eng{since, ago, while, by the time}) $\to$ Vérifier le sujet (3e pers. sg ?) $\to$ Appliquer la structure (auxiliaire + base verbale / V-ing / V-en) $\to$ Relire la phrase entière.
    \item \textbf{Vocabulaire (Derivation/Cloze) :} Identifier la catégorie grammaticale manquante (N, V, Adj, Adv) $\to$ Appliquer préfixes/suffixes (\eng{-tion, -ment, un-, -ful, -ly}) $\to$ Vérifier l'orthographe (double consonnes, \eng{y $\to$ i}).
    \item \textbf{Composition (Writing) :} Analyser le sujet (type : article, lettre, essay ; registre : formel/informel) $\to$ Plan (Intro $\to$ 2-3 arguments + exemples $\to$ Conclusion) $\to$ Rédiger en surveillant les connecteurs (\eng{Furthermore, However, Consequently}) $\to$ Relire (concordance des temps, \eng{s} du présent, majuscules).
\end{itemize}

\end{bilan}

\begin{aretenir}[Checklist avant le Bac]
$\square$ Je maîtrise la formation et l'emploi des \textbf{12 temps verbaux} (affirmatif, négatif, interrogatif).\\
$\square$ Je distingue \eng{Present Perfect} / \eng{Past Simple} et \eng{will} / \eng{be going to} / \eng{Present Continuous} (futur).\\
$\square$ Je construis correctement les \textbf{modaux au passé} (\eng{should have V-en}, \eng{must have V-en}, etc.).\\
$\square$ Je transforme une phrase active en passive (et inversement) dans \textbf{tous les temps}.\\
$\square$ J'applique le \textbf{backshifting} (concordance des temps) et le changement de pronoms/marqueurs au discours rapporté.\\
$\square$ Je structure les \textbf{3 types de conditionnels} et j'utilise \eng{Wish / If only} correctement.\\
$\square$ Je relie mes idées avec des \textbf{connecteurs logiques} variés (addition, concession, cause, conséquence).\\
$\square$ Je forme les mots (dérivation) selon la catégorie grammaticale requise (N, V, Adj, Adv).\\
$\square$ J'évite les \textbf{pièges classiques francophones} : \eng{actually} (en fait), \eng{eventually} (finalement), \eng{sensible} (raisonnable), \eng{library} (bibliothèque), \eng{actuellement} = \eng{currently}.\\
$\square$ Je sais rédiger une \textbf{lettre formelle}, un \textbf{article de presse} et un \textbf{essai argumenté} avec la structure attendue.\\
$\square$ Je gère mon temps : 40 min compréhension, 40 min grammaire/vocabulaire, 1h20 composition (brouillon + propre).
\end{aretenir}

\findecours

\end{document}
